\documentclass[12pt, spanish]{book}
\usepackage{epsfig, bm, epsf, float}
\usepackage{amsmath}
\usepackage{array}
\usepackage{booktabs}
\usepackage[usenames, dvipsnames]{color}
\usepackage{mathrsfs}
\usepackage[utf8]{inputenc}
\usepackage{graphicx}
\usepackage[usenames]{color}
\usepackage{times}
\usepackage{eqlist}
\usepackage[spanish]{babel}
\usepackage{amssymb}
\setlength{\parindent}{0cm}
\usepackage{titlesec}
\usepackage{cancel}
\usepackage{subfigure}
\usepackage{amsmath}
\usepackage{amssymb}
\usepackage{amsthm}
\usepackage{algorithm}
\usepackage{algpseudocode}
\floatname{algorithm}{Algoritmo}
\renewcommand{\algorithmicrequire}{\textbf{Entradas:}}
\renewcommand{\algorithmicensure}{\textbf{Salidas:}}
\newtheorem{teorema}{Teorema}
\newtheorem{corolario}{Corolario}
\renewcommand\proofname{Demostración}
\usepackage{cite}
\usepackage[top=3cm, bottom=3cm, left=3cm, right=3cm]{geometry}
\emergencystretch=2em % Permite estirar espacios para evitar texto fuera de márgenes
\usepackage[hyphens,obeyspaces,spaces]{url} % Debe ir ANTES de hyperref
\usepackage{hyperref}
\hypersetup{
    colorlinks=true, % colores en lugar de cajas
    linkcolor=black, % color de enlaces internos
    citecolor=black, % color de las citas
    urlcolor=black % color de enlaces de URL
}
% Comando \code: texto monoespaciado que permite cortes de línea automáticos
% en /, _, ., - (a diferencia de \texttt que no permite cortes)
\newcommand{\code}[1]{\texttt{\detokenize{#1}}}
\DeclareUrlCommand\codeinline{\urlstyle{tt}}

% Notación de Dirac
\newcommand{\ket}[1]{\left| #1 \right\rangle}
\newcommand{\bra}[1]{\left\langle #1 \right|}
\newcommand{\braket}[2]{\left\langle #1 \mid #2 \right\rangle}
\titleformat{\section}
  {\normalfont\Large\bfseries}
  {\thesection}
  {0.5em}
  {}

\titleformat{\subsection}
  {\normalfont\large\bfseries}
  {\thesubsection}
  {0.5em}
  {}

\bibliographystyle{IEEEtran} % estilo IEEE para referencias

\title{Universidad de Nariño\\
Facultad de Ciencias Exactas y Naturales\\
Departamento de Física\\[0.5cm]
% Cambiamos 'scale' por 'width' para tener una medida física absoluta
\includegraphics[width=6cm]{Imagenes/logoudenar1.png}
\\[0.5cm]
Evaluación de la viabilidad técnica y optimización del algoritmo de Shor en procesadores NISQ: un estudio comparativo de fidelidad y transpilación en IBM Quantum.\\}
\author{Christopher Andrey Villota Freire\\
Director: Jorge Hernán López Melo.\\[1.5cm]} % Reducido de 3.8cm a 1.5cm
\date{San Juan de Pasto, mayo de 2026}

\begin{document}
\maketitle
\newpage
\renewcommand{\contentsname}{Contenido}
\tableofcontents
\listoffigures \listoftables

\newpage

\chapter*{Agradecimientos}
\addcontentsline{toc}{chapter}{Agradecimientos}

\vspace*{0.5cm}

Al culminar esta etapa, miro hacia atrás y comprendo que ningún logro es enteramente propio: detrás de cada paso hay manos que sostienen, voces que alientan y presencias que acompañan en silencio. Por ello, estas líneas son apenas un humilde intento de devolver, en forma de palabras, lo mucho que he recibido a lo largo de este camino.\\

En primer lugar, agradezco a \textbf{Dios}, fuente inagotable de fuerza, sabiduría y serenidad. A Él debo la salud, la paciencia y la lucidez necesarias para sostenerme en los momentos de duda, y la gratitud profunda por cada oportunidad que ha puesto en mi camino. Sin su guía silenciosa, ninguna de estas páginas habría sido posible.\\

A mis padres, \textbf{Claudia Freire} y \textbf{Carlos Villota}, mi gratitud más sincera y eterna. Gracias por creer en mí incluso cuando yo mismo dudaba, por entregarme lo mejor de ustedes sin reservas y por enseñarme, con el ejemplo, que el esfuerzo honesto y la perseverancia son los cimientos sobre los que se construye una vida con sentido. Cada renuncia suya se ha transformado en una conquista mía, y este trabajo es una pequeña forma de honrar todo lo que han sembrado en mí.\\

A todos los \textbf{profesores} que dejaron una huella en mi formación, mi reconocimiento sincero. Gracias por compartir, con generosidad y vocación, no solo el conocimiento de su disciplina, sino también la pasión con que la habitan. Cada clase, cada explicación paciente y cada pregunta provocadora contribuyó a moldear la mirada con la que hoy me asomo a la física y al mundo.\\

De manera muy especial, agradezco a mi director de trabajo de grado, el profesor \textbf{Jorge Hernán López Melo}, por su acompañamiento riguroso, su confianza y su disposición permanente a lo largo de este proceso. Gracias por todo lo aprendido dentro y fuera del aula, por las largas conversaciones que trascendieron lo estrictamente académico y por mostrarme, con su ejemplo, que enseñar es también un acto de generosidad y compromiso. Su guía paciente fue fundamental para que este trabajo encontrara su rumbo, y las enseñanzas que llevo conmigo van mucho más allá de lo que estas páginas pueden contener.\\

Por último, pero no por ello menos importante, agradezco a la \textbf{Universidad de Nariño}, casa de saber que me abrió sus puertas y se convirtió en mi segundo hogar durante estos años. En sus aulas, pasillos y laboratorios aprendí no solo física, sino también el valor del pensamiento crítico, del trabajo colectivo y del compromiso con la región. Espero que este modesto aporte sea una pequeña retribución a la oportunidad invaluable que me brindó de formarme como físico y como persona.\\

A todos, mi gratitud profunda y sincera.\\

\begin{flushright}
\textit{Christopher Andrey Villota Freire}
\end{flushright}

\vspace*{\fill}

\chapter*{Dedicatoria}
\addcontentsline{toc}{chapter}{Dedicatoria}

\vspace*{1.5cm}

\begin{flushright}
\textit{A mis padres, Claudia Freire y Carlos Villota,\\
faros silenciosos que han iluminado cada uno de mis pasos.\\
Su amor incondicional, su esfuerzo callado y su fe inquebrantable\\
en mí son la raíz profunda de todo lo que he construido.\\
Cada página de este trabajo lleva grabada la huella\\
de su sacrificio y la calidez de su abrigo.\\
A ustedes les debo no solo lo que he logrado,\\
sino, sobre todo, la persona en la que me he convertido.}\\[0.8cm]

\textit{A mi familia, hilo invisible que sostiene cada uno de mis días.\\
A quienes celebraron mis triunfos\\
y abrazaron mis tropiezos, transformando\\
los momentos difíciles en oportunidades para crecer.\\
Su presencia ha sido el cimiento sereno\\
sobre el que he podido edificar mis sueños.}\\[0.8cm]

\textit{A mis amigos de la carrera, compañeros de ruta\\
en este largo y exigente camino por la física.\\
Gracias por hacer más llevadera cada jornada,\\
por convertir las dificultades en risas compartidas\\
y por recordarme, en los momentos justos,\\
que ningún reto se enfrenta verdaderamente solo.\\
Su compañía hizo que el aprendizaje fuera más humano\\
y que la universidad se sintiera como un hogar.}\\[0.8cm]

\textit{Y a cada persona que, con un gesto, una palabra oportuna\\
o un grano de arena ofrecido sin esperar nada a cambio,\\
contribuyó a forjar a quien hoy escribe estas líneas.\\
Este logro también es suyo.}\\[0.8cm]

\textit{Gracias, de corazón.}
\end{flushright}

\vspace*{\fill}

\chapter*{Resumen}

Este trabajo de grado evalúa la viabilidad técnica y el desempeño del algoritmo de Shor para la factorización del entero \(N = 15\) en procesadores cuánticos NISQ de IBM, comparando los resultados obtenidos en tres entornos de ejecución progresivamente más realistas: simulación ideal, simulación con modelo de ruido calibrado (\textit{Fake Backend} basado en \texttt{ibm\_torino}) y ejecución en el procesador real \texttt{ibm\_torino} (familia Heron~r1, 133 qubits superconductores). La implementación se realizó en Python utilizando el SDK Qiskit. Se aplicaron tres estrategias de optimización y supresión de errores: el algoritmo de ruteo SABRE para adaptar el circuito a la topología \textit{Heavy-Hex}, el Desacoplamiento Dinámico (DD) para reenfocar la decoherencia y el Pauli Twirling (PT) para convertir errores coherentes en ruido estocástico. El desempeño se cuantificó mediante la fidelidad de Hellinger y la probabilidad de éxito experimental (\(PST\)). Los resultados confirman la factorización experimental de \(N = 15 = 3 \times 5\) en hardware real, con una señal máxima de \(PST = 82.8\%\) y una tasa de éxito cercana al \(90\%\) en la extracción de los factores no triviales sobre las bases coprimas válidas. El nivel de optimización del transpilador resultó ser el factor más determinante para la viabilidad del algoritmo, con una ganancia aproximada de \(64\times\) en señal. Por contraste, las técnicas de supresión de errores mostraron un efecto contraproducente en el régimen estudiado: el Desacoplamiento Dinámico degradó la señal en más de 60 puntos porcentuales. Se documenta además que los \textit{Fake Backends} subestiman la calidad del hardware real en circuitos de profundidad moderada. El trabajo aporta evidencia experimental rigurosa y una metodología reproducible para futuros estudios en computación cuántica en la Universidad de Nariño.\\

\noindent\textbf{Palabras clave:} computación cuántica, algoritmo de Shor, era NISQ, IBM Quantum, transpilación, SABRE, Desacoplamiento Dinámico, Pauli Twirling, fidelidad de Hellinger.

\chapter*{Abstract}

This thesis evaluates the technical feasibility and performance of Shor's algorithm for the factorization of the integer \(N = 15\) on IBM NISQ quantum processors, comparing the results obtained in three progressively more realistic execution environments: ideal simulation, simulation with a calibrated noise model (\textit{Fake Backend} based on \texttt{ibm\_torino}), and execution on the real \texttt{ibm\_torino} processor (Heron~r1 family, 133 superconducting qubits). The implementation was carried out in Python using the Qiskit SDK. Three optimization and error suppression strategies were applied: the SABRE routing algorithm to adapt the circuit to the \textit{Heavy-Hex} topology, Dynamical Decoupling (DD) to refocus decoherence, and Pauli Twirling (PT) to convert coherent errors into stochastic noise. Performance was quantified using the Hellinger fidelity and the experimental probability of success (\(PST\)). The results confirm the experimental factorization of \(N = 15 = 3 \times 5\) on real hardware, with a maximum signal of \(PST = 82.8\%\) and a success rate close to \(90\%\) in the extraction of non-trivial factors over the valid coprime bases. The transpiler optimization level proved to be the most decisive factor for the viability of the algorithm, yielding an approximate \(64\times\) gain in signal. In contrast, the error suppression techniques showed a counterproductive effect in the studied regime: Dynamical Decoupling degraded the signal by more than 60 percentage points. It is also documented that \textit{Fake Backends} underestimate the quality of real hardware in moderate-depth circuits. This work provides rigorous experimental evidence and a reproducible methodology for future research in quantum computing at the University of Nariño.\\

\noindent\textbf{Keywords:} quantum computing, Shor's algorithm, NISQ era, IBM Quantum, transpilation, SABRE, Dynamical Decoupling, Pauli Twirling, Hellinger fidelity.

\chapter*{Glosario}
\addcontentsline{toc}{chapter}{Glosario}

\textbf{ALAP (\textit{As Late As Possible}):} criterio de planificación que retrasa cada operación lo máximo posible dentro del circuito, dejando los huecos de inactividad al inicio para insertar pulsos de protección.\\

\textbf{AQFT:} versión recortada de la Transformada Cuántica de Fourier en la que se descartan las rotaciones de ángulo más pequeño; reduce la profundidad del circuito a cambio de una pequeña pérdida de precisión.\\

\textbf{Backend:} dispositivo donde finalmente corre el circuito, ya sea una máquina física en un laboratorio de IBM o un programa que imita su comportamiento.\\

\textbf{Basis gates (compuertas nativas):} repertorio reducido de operaciones que el procesador sabe ejecutar directamente; toda compuerta del algoritmo se traduce a este conjunto antes de correr.\\

\textbf{Bitstring:} cadena de ceros y unos que resulta de medir varios qubits a la vez. Cada repetición del experimento produce una de estas cadenas.\\

\textbf{Coprimo:} se dice de dos enteros que solo comparten el 1 como divisor común; condición que debe cumplir la base \(a\) frente a \(N\) en el algoritmo de Shor.\\

\textbf{Crosstalk:} interferencia indeseada entre qubits vecinos: al accionar uno, los de al lado se contaminan ligeramente, alterando la operación pretendida.\\

\textbf{Eagle / Heron / Flamingo:} sucesivas familias de procesadores superconductores de IBM. Eagle tiene 127 qubits, Heron r1 (el utilizado en este trabajo) 133, y Flamingo es la generación anunciada como siguiente paso.\\

\textbf{ECR (\textit{Echoed Cross-Resonance}):} compuerta nativa de dos qubits en algunos procesadores de IBM; lógicamente equivale a la CNOT salvo por rotaciones locales adicionales.\\

\textbf{Fake Backend:} simulador local que imita un procesador real cargando una "fotografía" de sus parámetros de calibración; permite ensayar el algoritmo sin consumir tiempo de cómputo en la nube.\\

\textbf{Heurística:} regla práctica que busca una buena solución sin garantizar la óptima, útil cuando explorar todas las posibilidades resultaría inviable en tiempo de cómputo.\\

\textbf{Idle time (tiempo de inactividad):} intervalo durante el cual el qubit no participa en ninguna compuerta pero sigue expuesto al entorno y, por tanto, a la pérdida de coherencia.\\

\textbf{Job / Job ID:} envío individual de un circuito al servicio en la nube de IBM. Cada uno recibe un identificador único para recuperar después sus resultados.\\

\textbf{Layout:} asignación inicial que decide qué qubit lógico del algoritmo se coloca sobre qué qubit físico del chip antes de empezar a ejecutar.\\

\textbf{Little-endian:} convención de orden en la que el bit menos significativo se escribe primero. Qiskit la adopta, lo que obliga a invertir cadenas al traducir resultados a enteros.\\

\textbf{Open Plan:} modalidad gratuita de acceso a IBM Quantum, con un cupo mensual limitado de tiempo de procesador real.\\

\textbf{Outcome:} cada resultado distinto observado al medir el registro. Cuantos más outcomes únicos aparezcan, más dispersa está la probabilidad y peor es la señal del algoritmo.\\

\textbf{Pipeline:} secuencia ordenada de pasos en la que la salida de cada etapa alimenta a la siguiente; en transpilación describe las fases sucesivas de adaptación del circuito al hardware.\\

\textbf{POVM (\textit{Positive Operator-Valued Measure}):} formalismo general de medida cuántica, más amplio que las mediciones proyectivas, capaz de describir cualquier procedimiento experimental realista de extracción de información.\\

\textbf{Qiskit:} entorno de desarrollo en Python creado por IBM para construir, simular y enviar circuitos cuánticos a sus procesadores.\\

\textbf{QPU (\textit{Quantum Processing Unit}):} chip físico que ejecuta las operaciones cuánticas; análogo de la CPU en computación clásica.\\

\textbf{Routing (enrutamiento):} etapa de la transpilación que decide cómo mover la información entre qubits no vecinos, insertando SWAPs cuando hace falta.\\

\textbf{RSA:} criptosistema ampliamente desplegado en Internet cuya seguridad depende de la dificultad clásica de factorizar números muy grandes, problema que el algoritmo de Shor amenaza con resolver rápidamente.\\

\textbf{SamplerV2:} primitiva del entorno Qiskit que envía circuitos al hardware y devuelve la distribución estadística de los resultados de medición.\\

\textbf{SDK (\textit{Software Development Kit}):} colección de bibliotecas y herramientas con las que se programa para una plataforma específica; aquí se refiere al kit de Qiskit para hardware de IBM.\\

\textbf{Semilla (\textit{seed}):} número que inicializa un proceso pseudoaleatorio; fijarla garantiza que repetir el experimento con los mismos datos arroje siempre el mismo resultado.\\

\textbf{Shot:} cada disparo individual del circuito en el dispositivo. Promediar muchos shots permite reconstruir la distribución estadística de las salidas.\\

\textbf{Snapshot:} fotografía instantánea de los parámetros de calibración del procesador (tiempos, errores, conectividad) que se carga en los simuladores para imitar su comportamiento real.\\

\textbf{Trampas de iones:} alternativa tecnológica a los qubits superconductores, en la que átomos cargados se confinan con campos electromagnéticos y se controlan con láseres.\\

\textbf{Twirling:} técnica genérica de aleatorización que envuelve una operación entre otras escogidas al azar, con el fin de simetrizar el ruido y convertir errores complejos en otros más sencillos de tratar estadísticamente.\\

\textbf{Ventaja cuántica:} situación en la que un computador cuántico resuelve un problema sustancialmente más rápido que cualquier algoritmo clásico conocido; el caso de Shor frente a la factorización es uno de los ejemplos más espectaculares.

\chapter{Introducción}
\label{chap:introduccion}

La computación cuántica se ha consolidado como una de las fronteras más activas de la física contemporánea y la ciencia de la computación. Fundamentada en los principios de superposición y entrelazamiento de la mecánica cuántica, esta disciplina promete abordar problemas que resultan intratables para los computadores clásicos, con aplicaciones que abarcan desde la simulación de sistemas moleculares hasta la optimización combinatoria. Dentro de este panorama, el algoritmo de Shor ocupa un lugar paradigmático: demostró teóricamente que un computador cuántico puede factorizar un entero \(N\) en tiempo polinomial, una ventaja exponencial sobre los mejores algoritmos clásicos conocidos y una amenaza directa a los sistemas criptográficos de clave pública como RSA. Su materialización práctica, sin embargo, enfrenta un obstáculo físico fundamental: el hardware cuántico actual pertenece a la era NISQ (\textit{Noisy Intermediate-Scale Quantum}), caracterizada por un número intermedio de qubits y por la ausencia de corrección de errores cuánticos a escala completa. En este régimen, la decoherencia y las tasas de error de las compuertas lógicas restringen severamente la complejidad de los algoritmos que pueden ejecutarse con resultados distinguibles del ruido.\\

El presente trabajo de grado nace de la confluencia de dos circunstancias. Por un lado, la democratización del acceso a procesadores cuánticos reales a través de la plataforma IBM Quantum y su ecosistema de desarrollo Qiskit, que permite a estudiantes e investigadores ejecutar circuitos en hardware superconductor de última generación desde cualquier lugar del mundo. Por otro lado, la escasez de trabajos experimentales en computación cuántica reportados desde la Universidad de Nariño y la región. Esta investigación se propone cerrar dicha brecha, aportando evidencia experimental sobre el desempeño del algoritmo de Shor en dispositivos reales y generando una metodología reproducible que sirva como referencia para futuros estudios en la institución.\\

Para la evaluación experimental, este trabajo se enfoca en la instancia pedagógica canónica del algoritmo, la factorización de \(N = 15\). A pesar de su tamaño reducido, este caso constituye una prueba de estrés exigente para el hardware actual, pues exige circuitos profundos de exponenciación modular y la aplicación de la Transformada Cuántica de Fourier sobre un registro de varios qubits. La ejecución en un procesador real introduce además tres desafíos físicos interrelacionados: la conectividad restringida de la topología \textit{Heavy-Hex} de los procesadores IBM, que obliga al transpilador a insertar compuertas SWAP adicionales y amplificar la profundidad del circuito; la fidelidad limitada de las compuertas nativas de dos qubits, cuyos errores se acumulan con cada operación; y los tiempos de coherencia finitos, que imponen un límite estricto a la duración total del experimento. Para contrarrestar estas limitaciones se evaluaron tres estrategias complementarias: el algoritmo de ruteo SABRE, que minimiza heurísticamente la inserción de SWAPs durante la transpilación; el Desacoplamiento Dinámico (DD), que reenfoca la decoherencia de baja frecuencia mediante secuencias de pulsos en los intervalos de inactividad; y el Pauli Twirling (PT), que convierte errores coherentes en ruido estocástico promediable.\\

El diseño experimental se articuló en tres fases progresivamente más realistas. En la primera se verificó la corrección del circuito mediante simulaciones ideales, estableciendo una línea base libre de ruido y caracterizando el sobrecosto introducido por la topología del procesador. En la segunda, los circuitos optimizados se ejecutaron sobre un simulador con modelo de ruido calibrado (\textit{Fake Backend}) basado en el procesador \texttt{ibm\_torino}. En la tercera fase, los circuitos validados se enviaron al procesador real \texttt{ibm\_torino} (familia Heron~r1, 133 qubits superconductores) a través del IBM Quantum Open Plan, evaluando el impacto del nivel de optimización del transpilador, la selección del coeficiente clásico de pre-procesamiento y la activación de las técnicas de supresión de errores. Los resultados obtenidos permitieron demostrar experimentalmente la factorización de \(N = 15 = 3 \times 5\) en hardware real, cuantificar la degradación de la señal cuántica a lo largo de la cadena ideal~\(\to\)~simulador con ruido~\(\to\)~hardware, y obtener un hallazgo metodológico inesperado: en el régimen estudiado, la elección del nivel de optimización del transpilador resulta determinante para la viabilidad del algoritmo, mientras que las técnicas de supresión de errores no siempre son beneficiosas y pueden incluso degradar la señal cuando se aplican de forma indiscriminada.\\

El presente documento se organiza en siete capítulos. Tras esta introducción, el \textbf{Capítulo~2} formula el planteamiento del problema y la pregunta de investigación que guía el estudio, y el \textbf{Capítulo~3} enuncia los objetivos generales y específicos del trabajo. El \textbf{Capítulo~4} desarrolla el marco referencial, que incluye los antecedentes de implementaciones previas del algoritmo de Shor en diversas plataformas y un marco teórico que construye el formalismo necesario para comprender el algoritmo, el contexto NISQ, la arquitectura de IBM Quantum, los procesos de transpilación y supresión de errores y las métricas de evaluación empleadas. El \textbf{Capítulo~\ref{chap:metodologia}} expone el diseño metodológico y la implementación experimental en sus tres fases. El \textbf{Capítulo~\ref{chap:resultados}} presenta los resultados organizados en las mismas tres fases: simulación ideal, simulación con ruido y ejecución en hardware real. Finalmente, el \textbf{Capítulo~\ref{cap:conclusiones}} sintetiza las conclusiones del trabajo, identifica sus limitaciones y propone perspectivas futuras para la continuación de esta línea de investigación en la Universidad de Nariño.


\chapter{Planteamiento del problema}

Desde hace varias décadas, la comunidad científica ha anticipado una revolución
computacional basada en las propiedades de la mecánica cuántica \cite{Hidary2021,Singh2016}.
Los computadores cuánticos explotan fenómenos como la superposición y el
entrelazamiento \cite{Norsen2017,Zettili2009} para abordar problemas
intratables para la computación clásica, tales como la simulación de nuevos materiales
\cite{Ma2020}, el diseño de fármacos \cite{Kumar2024} y, notablemente, la
factorización de enteros grandes, lo cual representa una amenaza teórica para
los sistemas criptográficos de clave pública actuales, como RSA
\cite{Buchanan2017, Harrow2017}. Sin embargo, la materialización de estas promesas enfrenta un obstáculo
físico fundamental: la necesidad de un hardware tolerante a fallos (FTQC).
Aunque existen anuncios de procesadores con más de 1000 qubits físicos
\cite{AbuGhanem2025}, la industria aún se encuentra en la era de los Dispositivos
Cuánticos de Escala Intermedia Ruidosos (NISQ, por sus siglas en inglés)
\cite{IBM20251,AbuGhanem2025}. En este contexto, los qubits son susceptibles
a la decoherencia y las compuertas lógicas poseen tasas de error que se
acumulan rápidamente con la profundidad del circuito, lo que limita severamente
la ejecución de algoritmos complejos como el de Shor.\\

Esta limitación física resulta particularmente crítica en el algoritmo de Shor,
cuya ventaja exponencial teórica exige una profundidad de circuito considerable
debido a la exponenciación modular y la Transformada Cuántica de Fourier
(QFT) \cite{villarreal}. La ejecución directa de estas rutinas tiende a
exceder los tiempos de coherencia ($T_{1}, T_{2}$) de los dispositivos
actuales, produciendo salidas dominadas por ruido en lugar del espectro
esperado \cite{villarreal, Skosana2021}. En el plano experimental, esto
obliga a traducir el algoritmo teórico a circuitos físicos eficientes mediante
estrategias de compilación, transpilación, mitigación y supresión de errores que respeten
la topología del procesador y sus restricciones de ruido \cite{CLASSICQ2025}.
De este modo, la brecha entre la formulación teórica y su ejecución concreta
en hardware NISQ no constituye dos problemas independientes, sino la
manifestación unificada de una misma causa física.\\

En el ámbito local, la Universidad de Nariño no dispone de infraestructura
cuántica propia, pero el acceso remoto a la nube de IBM \cite{IBM20252}
habilita la investigación experimental de estas limitaciones. No obstante, persiste
un vacío en la comprensión práctica sobre cuánto se degradan estos
algoritmos al someterlos a las restricciones físicas reales. El problema
central no radica en la mera ejecución del algoritmo, sino en cuantificar la
brecha entre la simulación ideal y la ejecución ruidosa para casos abordables
en NISQ (como $N=15$), determinando si las técnicas actuales de transpilación
y supresión de errores permiten obtener resultados distinguibles del ruido.\\

En este contexto de hardware real pero limitado, surge la siguiente pregunta
de investigación:

¿Cuál es el desempeño y la viabilidad técnica de implementar versiones
optimizadas para NISQ del algoritmo de Shor para enteros pequeños en
procesadores cuánticos de IBM, comparando su fidelidad, profundidad de
circuito y probabilidad de éxito frente a simulaciones ideales y modelos
de ruido?

\chapter{Objetivos}

\section{Objetivo general}
Evaluar la viabilidad técnica y el desempeño de implementaciones educativas y
optimizadas del algoritmo de Shor para enteros pequeños ($N=15$) en procesadores
cuánticos NISQ de IBM, mediante un análisis comparativo de fidelidad, profundidad
de circuito y probabilidad de éxito entre simuladores con modelos de ruido y
la ejecución en hardware real.

\section{Objetivos específicos}
\begin{enumerate}
    \item Diseñar y caracterizar en Qiskit los bloques fundamentales del algoritmo
          (Transformada Cuántica de Fourier y Exponenciación Modular),
          analizando el impacto de estrategias de transpilación avanzadas (como
          el algoritmo de ruteo SABRE) en la reducción de la profundidad del
          circuito y el conteo de compuertas CNOT \cite{villarreal}.

    \item Implementar estrategias de optimización híbrida mediante pre-procesamiento
          clásico para identificar los coeficientes $a$ que minimicen teóricamente
          la complejidad del circuito de exponenciación modular, haciendo
          viable la ejecución experimental para casos de estudio como $N=15$
          en dispositivos NISQ \cite{villarreal}.

    \item Ejecutar las versiones optimizadas del algoritmo utilizando ``Fake
          Backends'' (simuladores locales con ruido real) y hardware de IBM,
          aplicando técnicas de supresión de errores (como \textit{Dynamical
              decoupling} y \textit{Pauli Twirling}) para mejorar la relación
          señal-ruido de los resultados \cite{villarreal}.

    \item Cuantificar la degradación de la señal cuántica comparando los resultados
          experimentales frente a las simulaciones ideales, utilizando la métrica
          de \textit{Fidelidad de Hellinger} (Medida de similitud) y \textit{Probabilidad
              de Éxito} (tasa de obtención del periodo correcto).
\end{enumerate}

\chapter{Marco referencial}

\section{Marco de antecedentes}

Desde finales del siglo XIX y durante las primeras décadas del siglo XX, la física experimentó una transformación profunda. Los problemas asociados a la radiación de cuerpo negro, el efecto fotoeléctrico y otros fenómenos condujeron a la formulación de la mecánica cuántica, disciplina que introdujo nociones como la cuantización, la superposición y el principio de incertidumbre \cite{Kleppner2000}. Este desarrollo histórico, junto con las tecnologías instrumentales que desencadenó (tales como el control preciso de estados cuánticos y técnicas de medida coherente), es la base indispensable que hace posible hoy manipular qubits y diseñar algoritmos cuánticos implementables en hardware real \cite{Kleppner2000,Nielsen2024}.\\

La unidad de información en este nuevo paradigma es el qubit, que puede existir en superposición y entrelazarse con otros, propiedades que habilitan algoritmos con ventajas teóricas sobre sus contrapartes clásicas \cite{Nielsen2024}. En particular, el algoritmo de Shor, propuesto en 1994, demostró que la factorización de enteros (problema central en criptografía asimétrica) puede resolverse en tiempo polinómico sobre un computador cuántico, lo que generó un interés inmediato en la verificación experimental de sus subrutinas críticas: estimación de fase, búsqueda del orden y la transformada cuántica de Fourier (QFT) \cite{Shor1994,Nielsen2024}. No obstante, la distancia entre la formulación matemática y la ejecución en hardware físico real se mantiene condicionada por limitaciones tecnológicas (tasas de error de las compuertas, decoherencia y conectividad), lo que ha orientado la investigación hacia versiones compiladas o reducidas de Shor y hacia esquemas híbridos para sortear las restricciones de los dispositivos NISQ (Noisy Intermediate-Scale Quantum) \cite{Vandersypen2001,Roffe2019}.\\

Las primeras implementaciones experimentales de Shor y sus bloques constitutivos se realizaron en plataformas diversas, aportando lecciones metodológicas vigentes. La demostración en NMR (resonancia magnética nuclear) que factorizó $N=15$ validó la factibilidad conceptual de Shor y permitió desarrollar técnicas de control y verificación trasladables a otras tecnologías \cite{Vandersypen2001}. En fotónica se llevaron a cabo versiones compiladas que corroboraron la generación de entrelazamiento multipartito y la ejecución de subrutinas de la QFT bajo condiciones reales \cite{Lanyon2007,Lu2007,Politi2009}, mientras que trabajos con trampas de iones documentaron rutas escalables para implementar operaciones aritméticas modulares y transformadas cuánticas complejas \cite{Rossi2022}. Estas implementaciones históricas establecieron dos convenciones prácticas esenciales: (i) la necesidad de reducir la profundidad de los circuitos y minimizar operaciones multi-qubit costosas, y (ii) la conveniencia de validar experimentalmente mediante métricas robustas (fidelidad, probabilidades de salida y chequeos de entrelazamiento) y comparaciones con simuladores ideales y ruidosos \cite{Vandersypen2001,Lanyon2007,Lu2007,Rossi2022}.\\

La democratización del acceso a QPUs superconductoras a través de servicios en la nube, particularmente la plataforma IBM Quantum y su ecosistema Qiskit, cambió radicalmente el panorama experimental: hoy es factible para equipos de investigadores y estudiantes enviar circuitos a procesadores remotos, ejecutar simuladores con modelos de ruido y ensayar técnicas de supresión y mitigación sin disponer del laboratorio físico \cite{Saxena2021,IBMqiskit}. Esta accesibilidad ha impulsado una proliferación de estudios que implementan versiones compiladas o reducidas de Shor directamente en backends de IBM, documentando el flujo completo de experimentación (desde el mapeo lógico-físico hasta la aplicación de mitigaciones en lectura y extrapolación a ruido cero) y ofreciendo buenas prácticas para garantizar la reproducibilidad \cite{Saxena2021,DeCusatis2021,IBM}. La evidencia reciente muestra que, con ajustes de compilación, supresión y mitigación apropiados, es posible obtener factorizaciones experimentales reproducibles para valores de $N$ pequeños o moderados en QPUs de IBM, aunque la probabilidad de éxito se degrada con la profundidad y el número de puertas de dos qubits \cite{Skosana2021,Rossi2022,Roffe2019}.\\

Entre las contribuciones más relevantes para la implementación en IBM destaca la demostración de order-finding para $N=21$ ejecutada sobre QPUs de IBM, la cual empleó una versión compilada de la estimación de fase y puertas Toffoli aproximadas para reducir la complejidad física, además, el trabajo verificó la presencia de entrelazamiento entre registros, ofreciendo un precedente metodológico claro para validar implementaciones en backends reales con recursos limitados \cite{Skosana2021}. Trabajos más recientes han refinado esta metodología, demostrando que la viabilidad de ejecutar Shor en hardware de IBM (como \textit{Eagle} o \textit{Heron}) depende críticamente de técnicas de pre-procesamiento clásico para seleccionar coeficientes $a$ que generen periodos cortos, y del uso de algoritmos de ruteo estocástico (SABRE) para adaptar el circuito a la topología \textit{Heavy-Hex} del chip \cite{villarreal}. Además, este estudio resalta la importancia de utilizar ``Fake Backends'' (simuladores con modelos de ruido reales) como paso intermedio de validación para predecir la similitud antes de consumir recursos de hardware \cite{villarreal, DeCusatis2021}.\\

La bibliografía metodológica y los tutoriales oficiales de Qiskit aportan, además, conocimiento operativo imprescindible: describen herramientas de compilación (incluyendo heurísticas para el mapeo de qubits y la reducción de compuertas), técnicas de supresión (Dynamical Decoupling, Pauli Twirling) o mitigación (corrección de readout, extrapolación a ruido cero) y prácticas de registro (anotar la versión del backend, snapshots y calibraciones) que permiten comparar de forma rigurosa simulador ideal, simulador con ruido y ejecución en hardware \cite{DeCusatis2021, IBM,IBMqiskit}. La convergencia de estas herramientas y resultados sugiere una hoja de ruta experimental consistente: escoger una variante de Shor apta para NISQ, documentar el backend y las calibraciones, aplicar protocolos de supresión o mitigación y cuantificar métricas clave (probabilidad de éxito, fidelidad, dependencia con la profundidad) mediante comparaciones controladas entre simuladores y QPUs reales \cite{Rossi2022,DeCusatis2021,Roffe2019}. A pesar de estos avances, la escalabilidad hacia enteros de relevancia práctica sigue limitada por restricciones técnicas (principalmente la imperfección de compuertas de dos qubits y la falta de corrección de errores a gran escala), por lo que la investigación aplicada continúa orientada a mejorar la relación coste-beneficio de las implementaciones y a desarrollar procedimientos reproducibles en la nube \cite{Roffe2019,Vandersypen2001,Rossi2022}.\\

En el ámbito nacional, Colombia ha mostrado señales de un interés institucional creciente hacia tecnologías cuánticas, con la llegada del primer computador cuántico a una universidad colombiana \cite{UniversidaddelosAndes2024} y la inclusión de estas tecnologías en convocatorias de Minciencias \cite{MinisteriodeCienciasdeColombia2025,FundacionAlejandroAngelEscobar2020}. Sin embargo, la revisión documental sugiere que en la región de Nariño y en la Universidad de Nariño existe aún una escasez de trabajos experimentales publicados que documenten implementaciones de Shor en QPUs de IBM \cite{udenar}. Este trabajo se fundamenta en estos antecedentes para llenar dicho vacío, adoptando las metodologías de optimización de circuito y validación en simuladores ruidosos propuestas en la literatura reciente para evaluar el desempeño real del algoritmo en el entorno académico local \cite{villarreal, udenar}. El objetivo es aportar evidencia experimental rigurosa y guías prácticas para la comunidad de Pasto y la Universidad de Nariño, aprovechando los recursos formativos y técnicos disponibles junto a la experiencia metodológica acumulada en la literatura \cite{IBMqiskit,Saxena2021}.

\newpage

% =========================================================================
% BLOQUE A — Sección 4.2 reducida (cuerpo principal)
% =========================================================================

\section{Fundamentos de la computación cuántica}
\label{sec:fundamentos}
 
Esta sección recopila los elementos del formalismo de la computación cuántica que seran necesarios para entender la implementación experimental del algoritmo de Shor sobre el hardware NISQ de IBM Quantum. El detalle adicional sobre el modelo de ruido y la mitigación de errores de lectura se difiere al Apéndice~\ref{app:ruido}.
 
\subsection{Qubit, notación y matrices de Pauli}
\label{sec:qubit}
 
La computación cuántica se formula sobre espacios de Hilbert complejos de dimensión finita \cite{Nielsen2024, Wong2022}. En la notación de Dirac, un vector de estado se denota $|\psi\rangle$ (\textit{ket}) y su dual $\langle\psi|$ (\textit{bra}), con producto interno $\langle\phi|\psi\rangle$.\\
 
El \textit{bit cuántico} o \textit{qubit} es la unidad fundamental de información cuántica. Su espacio de estados es $\mathcal{H} = \mathbb{C}^2$ y la base ortonormal estándar, denominada \textit{base computacional}, está dada por
\begin{equation}
    |0\rangle = \begin{bmatrix} 1 \\ 0 \end{bmatrix}, \qquad
    |1\rangle = \begin{bmatrix} 0 \\ 1 \end{bmatrix}.
    \label{eq:base_computacional}
\end{equation}
El estado más general de un qubit es una superposición coherente
\begin{equation}
    |\psi\rangle = \alpha|0\rangle + \beta|1\rangle, \qquad \alpha, \beta \in \mathbb{C}, \qquad |\alpha|^2 + |\beta|^2 = 1,
    \label{eq:qubit_general}
\end{equation}
donde $\alpha$ y $\beta$ son las \textit{amplitudes de probabilidad}: al medir en la base computacional, los resultados $|0\rangle$ y $|1\rangle$ se obtienen con probabilidades $|\alpha|^2$ y $|\beta|^2$ respectivamente \cite{Nielsen2024, Wong2022}.\\
 
Las matrices de Pauli (Figura ~\ref{fig:compuertas}), junto con la identidad, generan el espacio de operadores hermíticos sobre $\mathbb{C}^2$ \cite{Nielsen2024} y aparecen sistemáticamente en la descripción de compuertas y canales de ruido.\\

Cuando exista incertidumbre clásica adicional sobre el estado del sistema (por ejemplo, debido a la interacción con el entorno), se emplea el \textit{operador densidad} $\rho = |\psi\rangle\langle\psi|$ para estados puros, o $\rho = \sum_i p_i |\psi_i\rangle\langle\psi_i|$ con $\sum_i p_i = 1$ para ensambles estadísticos \cite{Nielsen2024, Watrous2018}. Este formalismo es indispensable para describir el ruido NISQ (Sección~\ref{sec:nisq}).

\subsection{Evolución, medición y sistemas compuestos}
\label{sec:evolucion_medicion_compuestos}
 
La evolución de un sistema cuántico cerrado se describe por una transformación unitaria $|\psi'\rangle = U|\psi\rangle$ con $U^{\dagger}U = I$, lo que garantiza la conservación de la norma; esta evolución unitaria se implementa en el modelo de circuitos como la aplicación sucesiva de \textit{compuertas cuánticas}. Las mediciones empleadas en este trabajo son \textit{proyectivas en la base computacional} $\{|0\rangle, |1\rangle\}$. Para sistemas de $n$ qubits, el espacio de estados es el producto tensorial $\mathcal{H}^{\otimes n} = \mathbb{C}^{2^n}$; un estado factorizado de $n$ qubits independientes se escribe $|\psi_1\rangle \otimes \cdots \otimes |\psi_n\rangle$ \cite{Nielsen2024}.
 
\subsection{Superposición y entrelazamiento}
\label{sec:superposicion_entrelazamiento}
 
El principio de superposición es consecuencia directa de la linealidad del espacio de estados: dada una base ortonormal $\{|e_i\rangle\}_{i=1}^{n}$, toda combinación lineal normalizada constituye un estado físico válido \cite{Nielsen2024, Wong2022}. La superposición permite que un registro de $n$ qubits codifique $2^n$ amplitudes simultáneamente, y es la base del paralelismo cuántico explotado por la transformada cuántica de Fourier en el algoritmo de Shor.\\
 
Un estado de un sistema bipartito $|\Psi\rangle \in \mathcal{H}_A \otimes \mathcal{H}_B$ es \textit{separable} si puede escribirse como $|\Psi\rangle = |\psi_A\rangle \otimes |\phi_B\rangle$, y \textit{entrelazado} en caso contrario \cite{Nielsen2024}. El entrelazamiento es un recurso fundamental en computación cuántica y, en el contexto experimental de este trabajo, su preservación a lo largo del circuito está limitada por la decoherencia (Sección~\ref{sec:nisq}).
 
\subsection{Modelo de circuitos y compuertas}
\label{sec:modelo_circuitos}
 
Las operaciones lógicas del algoritmo se implementan mediante \textit{compuertas cuánticas}, operadores unitarios sobre uno o varios qubits \cite{Nielsen2024}. Las compuertas de un solo qubit empleadas con mayor frecuencia se resumen en la Figura~\ref{fig:compuertas}; entre ellas destaca la \textit{compuerta Hadamard},
\begin{equation}
    H|0\rangle = \tfrac{1}{\sqrt{2}}(|0\rangle+|1\rangle), \qquad H|1\rangle = \tfrac{1}{\sqrt{2}}(|0\rangle-|1\rangle),
    \label{eq:hadamard}
\end{equation}

que genera superposiciones equitativas a partir de la base computacional y constituye el primer paso del registro de conteo en el algoritmo de Shor.\\
 
Las \textit{compuertas de rotación de fase}, centrales para la construcción de la transformada cuántica de Fourier (QFT), están dadas por
\begin{equation}
    R_k = \begin{bmatrix} 1 & 0 \\ 0 & e^{2\pi i/2^k} \end{bmatrix}.
    \label{eq:rotacion_fase}
\end{equation}
La generación de entrelazamiento y la aritmética modular del algoritmo de Shor requieren \textit{compuertas de dos qubits}. La más simple es la CNOT (\textit{Controlled-NOT}), que actúa como
\begin{equation}
    \text{CNOT} = \begin{bmatrix} 1 & 0 & 0 & 0 \\ 0 & 1 & 0 & 0 \\ 0 & 0 & 0 & 1 \\ 0 & 0 & 1 & 0 \end{bmatrix}, \qquad \text{CNOT}|a,b\rangle = |a, a \oplus b\rangle,
    \label{eq:cnot}
\end{equation}

\begin{figure}[h]
    \centering
    \includegraphics[width=0.5\textwidth]{Imagenes/compuertas.png}
    \vspace{-0.4cm}
    \caption{Nombres, símbolos y matrices unitarias para las compuertas comunes de un solo qubit \cite{Nielsen2024}.}
    \label{fig:compuertas}
\end{figure}

con $\oplus$ la suma módulo 2 \cite{Nielsen2024, Wong2022}. La CNOT es un caso particular de \textit{compuerta controlada-$U$}: dada una operación unitaria $U$ sobre un qubit, la operación controlada-$U$ aplica $U$ al qubit objetivo si y solo si el qubit de control está en $|1\rangle$ (Figura~\ref{fig:controlU}). La exponenciación modular $CU_{a,N}$, núcleo computacional del algoritmo de Shor, se construye precisamente como una secuencia de operaciones controladas de este tipo.
 
\begin{figure}[h]
    \centering
    \includegraphics[width=0.5\textwidth]{Imagenes/controlU.png}
    \caption{Representación de la operación controlada-$U$ \cite{Nielsen2024}.}
    \label{fig:controlU}
\end{figure}
 
\subsection{Ruido NISQ: \texorpdfstring{$T_1$, $T_2$}{T1, T2}, canales de error y errores de lectura}
\label{sec:nisq}
 
Las subsecciones precedentes asumen evolución unitaria ideal. En el hardware real, los qubits están inevitablemente acoplados a su entorno electromagnético y térmico, fenómeno que produce la \textit{decoherencia} y degrada la información cuántica \cite{Nielsen2024, Preskill}. En el régimen NISQ, la decoherencia, los errores de compuerta y los errores de lectura constituyen la principal limitación a la profundidad ejecutable de circuitos como el del algoritmo de Shor.\\
 
A nivel fenomenológico, la coherencia de un qubit se caracteriza por dos escalas temporales \cite{Nielsen2024, Preskill, Brand2022}:
\begin{itemize}
    \item El \textit{tiempo de relajación} $T_1$, asociado al decaimiento exponencial de la población del estado excitado $|1\rangle$ hacia el fundamental $|0\rangle$, $\rho_{11}(t) = \rho_{11}(0)\,e^{-t/T_1}$.
    \item El \textit{tiempo de decoherencia transversal} $T_2$, asociado al decaimiento de las coherencias fuera de la diagonal del operador densidad, $\rho_{01}(t) = \rho_{01}(0)\,e^{-t/T_2}$.
\end{itemize}
Ambos satisfacen la cota fundamental
\begin{equation}
    T_2 \leq 2T_1.
    \label{eq:T2_cota}
\end{equation}
 
El error introducido por una compuerta imperfecta se modela habitualmente componiendo la operación unitaria ideal con un \textit{canal despolarizante}, que en su forma estándar actúa como
\begin{equation}
    \mathcal{E}_{\text{depol}}(\rho) = (1-p)\,\rho + \frac{p}{d}\,I,
    \label{eq:canal_despolarizante}
\end{equation}
donde $d$ es la dimensión del espacio de Hilbert ($d=2$ para un qubit) y $p$ está directamente relacionado con la tasa de error de compuerta reportada por IBM en sus calibraciones diarias \cite{Nielsen2024, AbuGhanem2025}. Un mayor número de compuertas en el circuito se traduce en un canal despolarizante efectivo más severo y, por tanto, en una pérdida de fidelidad medible al final del algoritmo.\\
 
Finalmente, la medición proyectiva implementada en el hardware no es ideal: se modela mediante una \textit{matriz de confusión} (o de asignación) \cite{IBM_noise_models}
\begin{equation}
    \mathcal{A} = \begin{bmatrix} P(0|0) & P(0|1) \\ P(1|0) & P(1|1) \end{bmatrix},
    \label{eq:matriz_confusion}
\end{equation}
donde $P(m|s)$ es la probabilidad de registrar el resultado $m$ cuando el estado real del qubit es $|s\rangle$. La asimetría habitual entre $P(1|0)$ y $P(0|1)$ refleja la asimetría energética entre $|0\rangle$ y $|1\rangle$ en los qubits superconductores \cite{Preskill, DeCusatis2021}. La inversión de $\mathcal{A}$ permite mitigar el error de lectura a posteriori sobre las distribuciones de probabilidad medidas; el procedimiento concreto, junto con la representación de Kraus de los canales y la composición del canal despolarizante con la operación unitaria, se desarrolla en el Apéndice~\ref{app:ruido}.

% =========================================================================
% BLOQUE A — Sección 4.3 reescrita (cuerpo principal)
% =========================================================================

\section{El algoritmo de Shor}

El algoritmo de Shor \cite{Shor1994} resuelve, en un computador cuántico ideal y en tiempo polinomial, el problema de factorizar un entero compuesto impar $N$. Su importancia trasciende el interés algorítmico: la seguridad del criptosistema RSA, ampliamente desplegado en la infraestructura de Internet, descansa sobre la presunta dificultad clásica de este mismo problema. Los mejores algoritmos clásicos conocidos (en particular el \textit{General Number Field Sieve}) requieren un tiempo subexponencial en el número de bits $L = \lceil \log_2 N \rceil$, mientras que el algoritmo de Shor lo logra en $\mathcal{O}(L^3)$ operaciones cuánticas, una aceleración superpolinomial.\\

La idea central es indirecta. En lugar de atacar la factorización de forma directa, Shor reduce el problema a una tarea relacionada: la \textit{búsqueda de orden} \cite{Nielsen2024}. Dado un entero $a$ coprimo con $N$, se define el orden $r$ de $a$ módulo $N$ como el menor entero positivo que satisface
\[
    a^r \equiv 1 \pmod{N}.
\]
Una vez conocido $r$, una manipulación clásica elemental permite (con probabilidad bien acotada inferiormente) obtener un factor no trivial de $N$ a partir del cálculo del máximo común divisor entre $a^{r/2} \pm 1$ y $N$. La dificultad se traslada entonces a determinar $r$. Clásicamente esta tarea es exponencial en $L$, pero cuánticamente admite una solución eficiente porque $r$ es el período de la función $f(z) = a^z \bmod N$, y los algoritmos cuánticos pueden extraer periodicidades de funciones evaluadas en superposición mediante la transformada cuántica de Fourier.\\

La Figura~\ref{fig:esquema-shor} sintetiza la estructura del algoritmo en tres bloques, cuyos ingredientes se desarrollarán formalmente en las subsecciones siguientes.

 \begin{figure}[H]
     \centering
     \includegraphics[width=0.8\textwidth]{Imagenes/diagrama_shor.png}
     \caption{Esquema general del algoritmo de Shor: bloque clásico de pre-procesamiento, subrutina cuántica de búsqueda de orden y bloque clásico de post-procesamiento.}
     \label{fig:esquema-shor}
 \end{figure}

\newpage

\subsection{Transformada cuántica de Fourier (QFT)}
\label{sec:qft}

La transformada cuántica de Fourier (QFT) es el operador lineal que efectúa el cambio de la base computacional a la base de fase (o base de Fourier). Sobre un espacio de Hilbert de dimensión $N = 2^n$, la QFT se define por su acción sobre los estados ortonormales $\{|j\rangle : j = 0,\dots,N-1\}$ \cite{Nielsen2024}:
\begin{equation}
    \mathrm{QFT}|j\rangle = \frac{1}{\sqrt{N}} \sum_{k=0}^{N-1} e^{2\pi i\, j k / N}\, |k\rangle.
    \label{eq:qft_def}
\end{equation}

A partir de esta definición se obtiene, mediante la representación binaria de $j$ y $k$ y la separación de la parte entera y fraccionaria del producto $j\,2^{-l}$, la forma factorizada \cite{Nielsen2024}
\begin{equation}
    \mathrm{QFT}|j_1 j_2 \dots j_n\rangle = \frac{1}{2^{n/2}} \big(|0\rangle + e^{2\pi i 0.j_n}|1\rangle\big) \otimes \big(|0\rangle + e^{2\pi i 0.j_{n-1}j_n}|1\rangle\big) \otimes \dots \otimes \big(|0\rangle + e^{2\pi i 0.j_1j_2\dots j_n}|1\rangle\big).
    \label{eq:qft_producto}
\end{equation}
Esta representación revela que el entrelazamiento del estado de entrada se mapea a fases relativas locales en un estado completamente separable (Figura~\ref{fig:qft}), lo que dictamina la factibilidad de sintetizar el operador $\mathrm{QFT}$ en hardware mediante compuertas de Hadamard ($H$) locales y rotaciones de fase controladas $R_k = \mathrm{diag}(1, e^{2\pi i / 2^k})$, con compuertas $\mathrm{SWAP}$ al final para invertir el orden de los qubits y alinear la salida con la convención estándar del algoritmo \cite{Nielsen2024}.

\begin{figure}[h]
    \centering
    \includegraphics[width=\textwidth]{Imagenes/qft.png}
    \caption{Circuito eficiente para la transformada cuántica de Fourier. Este circuito se deriva fácilmente de la representación del producto (Ec.~\ref{eq:qft_producto}) para la transformada cuántica de Fourier. No se muestran las compuertas de intercambio al final del circuito, las cuales invierten el orden de los cúbits, ni los factores de normalización de $1/\sqrt{2}$ en la salida \cite{Nielsen2024}.}
    \label{fig:qft}
\end{figure}

\textbf{Limitaciones de la QFT en el régimen NISQ.}\\

Asintóticamente, la QFT exacta demanda $\mathcal{O}(n^2)$ compuertas. Sin embargo, en la arquitectura NISQ, las compuertas $R_k$ para valores grandes de $k$ inducen rotaciones con ángulos $\theta = 2\pi/2^k$ exponencialmente pequeños. Cuando la magnitud de estas rotaciones cae por debajo de la tasa de error intrínseca de las compuertas de dos qubits, estas operaciones no solo son redundantes, sino perjudiciales, ya que inyectan más ruido por decoherencia térmica que corrección de fase. Esto justifica la adopción de la transformada cuántica de Fourier aproximada (AQFT), donde se truncan las compuertas $R_k$ para $k > m$, reduciendo la profundidad del circuito a cambio de una degradación controlada de la fidelidad.

\subsection{Estimación de fase cuántica (QPE) y cotas de error}
\label{sec:qpe}

La QPE resuelve el problema de extraer el autovalor $e^{2\pi i \varphi}$ asociado a un autovector $|u\rangle$ de un operador unitario $U$. El circuito emplea dos registros: un registro de conteo de $t$ qubits inicializado en $|0\rangle^{\otimes t}$ y un registro de trabajo que contiene el estado $|u\rangle$ (Figura~\ref{fig:QPE_registros}).

\begin{figure}[h]
    \centering
    \includegraphics[width=\textwidth]{Imagenes/QPE_registros.png}
    \caption{Primera etapa del procedimiento de estimación de fase. Se han omitido los factores de normalización de $1/\sqrt{2}$, a la derecha \cite{Nielsen2024}.}
    \label{fig:QPE_registros}
\end{figure}

Tras la aplicación paralela de compuertas Hadamard al registro de conteo y la posterior aplicación de las operaciones controladas $c$-$U^{2^j}$, el estado del sistema evoluciona a
\begin{equation}
    |\psi_1\rangle = \frac{1}{2^{t/2}} \sum_{k=0}^{2^t-1} e^{2\pi i \varphi k} |k\rangle \otimes |u\rangle.
    \label{eq:qpe_estado}
\end{equation}
Este estado tiene exactamente la forma de una QFT aplicada al estado $|y\rangle$ (compárese con la Ec.~\ref{eq:qft_def}): es decir, $|\psi_1\rangle_{\text{conteo}} = \mathrm{QFT}|y\rangle$ cuando $\varphi = y/2^t$. Por consiguiente, la aplicación de la $\mathrm{QFT}^{\dagger}$ recupera determinísticamente el estado $|y\rangle$. Esta es la razón fundamental por la cual la estimación de fase funciona: la $\mathrm{QFT}^{\dagger}$ deshace la codificación de la fase en las amplitudes del registro de conteo y permite leer directamente el valor $y$ que codifica $\varphi = y/2^t$. El circuito completo de la estimación de fase cuántica se muestra en la Figura~\ref{fig:QPE_completa}.\\

\begin{figure}[htbp]
    \centering
    \includegraphics[width=0.6\textwidth]{Imagenes/QPE_completa.png}
    \caption{Esquema del procedimiento general de estimación de fase. Los \( t \) qubits superiores (la barra ‘/’ denota un haz de cables) \cite{Nielsen2024}.}
    \label{fig:QPE_completa}
\end{figure}

Si $\varphi$ puede expresarse exactamente en $t$ bits ($\varphi = y/2^t$), la ecuación anterior garantiza que la medición produce $|y\rangle$ con certeza. Si $\varphi$ no posee una expansión binaria finita exacta en $t$ bits, el proceso es probabilístico. La estimación de fase requiere acotar matemáticamente el error: para aproximar $\varphi$ a $n$ bits de precisión con una probabilidad de éxito mínima de $1-\epsilon$, el número de qubits de conteo $t$ requeridos obedece la cota \cite{Nielsen2024}
\begin{equation}
    t = n + \left\lceil \log_2\left( 2 + \frac{1}{2\epsilon} \right) \right\rceil.
    \label{eq:qpe_cota}
\end{equation}

En el hardware actual, el número estricto de qubits disponibles y su tiempo de coherencia transversal $T_2$ impiden satisfacer esta cota para precisiones altas, lo que obliga a emplear técnicas de aproximación y optimización \cite{villarreal}.\\

El procedimiento de la estimación de fase cuántica se resume formalmente en el Algoritmo~\ref{alg:qpe}.

\begin{algorithm}[H]
    \caption{Estimación de fase cuántica (QPE) \cite{Nielsen2024}}
    \label{alg:qpe}
    \begin{algorithmic}[1]
        \Require (1) Una caja negra que realiza la operación controlada $c$-$U^{2^j}$ para entero $j$, (2) un autoestado $|u\rangle$ de $U$ con autovalor $e^{2\pi i \varphi_u}$, y (3) $t = n + \lceil \log(2 + \frac{1}{2\epsilon}) \rceil$ qubits inicializados en $|0\rangle$.
        \Ensure Una aproximación de $n$ bits $\widetilde{\varphi}_u$ a $\varphi_u$.
        \Statex \textbf{Tiempo de ejecución:} $\mathcal{O}(t^2)$ operaciones y una llamada a la caja negra $c$-$U^{2^j}$. Éxito con probabilidad $\geq 1 - \epsilon$.
        \Statex
        \Statex \textbf{Procedimiento:}
        \State \textit{Estado inicial:} $|0\rangle^{\otimes t}|u\rangle$
        \State \textit{Crear superposición:} Aplicar $H^{\otimes t}$ al registro de conteo:
        \Statex \hspace{2em} $\to \dfrac{1}{\sqrt{2^t}} \displaystyle\sum_{j=0}^{2^t-1} |j\rangle |u\rangle$
        \State \textit{Aplicar caja negra:} Aplicar las operaciones controladas $c$-$U^{2^j}$:
        \Statex \hspace{2em} $\to \dfrac{1}{\sqrt{2^t}} \displaystyle\sum_{j=0}^{2^t-1} |j\rangle U^{j}|u\rangle = \dfrac{1}{\sqrt{2^t}} \displaystyle\sum_{j=0}^{2^t-1} e^{2\pi i j \varphi_u} |j\rangle |u\rangle$
        \State \textit{Aplicar QFT inversa:} Aplicar $\mathrm{QFT}^{\dagger}$ al registro de conteo:
        \Statex \hspace{2em} $\to |\widetilde{\varphi}_u\rangle|u\rangle$
        \State \textit{Medir:} Medir el registro de conteo para obtener $\widetilde{\varphi}_u$.
    \end{algorithmic}
\end{algorithm}


\subsection{Búsqueda de orden y exponenciación modular}
\label{sec:order_finding}

Para enteros positivos $x$ y $N$, $x < N$, coprimos, el orden $r$ de $x$ módulo $N$ es el menor entero positivo tal que $x^r \equiv 1 \pmod N$. El problema de búsqueda de orden consiste en determinar $r$ para un par $(x, N)$ dado. Es un problema difícil en un computador clásico, en el sentido de que no se conoce un algoritmo que lo resuelva usando recursos polinomiales en $L = \lceil \log_2 N \rceil$.\\  

El algoritmo cuántico de búsqueda de orden es el algoritmo de estimación de fase aplicado al operador unitario de multiplicación modular sobre el registro de trabajo \cite{Nielsen2024}:
\begin{equation}
    U_{x,N} |y\rangle \equiv |\,x y \pmod N\,\rangle, \qquad y \in \{0,1,\dots,N-1\}.
    \label{eq:U_def}
\end{equation}
Los estados definidos por
\begin{equation}
    |u_s\rangle = \frac{1}{\sqrt{r}} \sum_{k=0}^{r-1} e^{-2\pi i s k / r} |x^k \bmod N\rangle, \qquad s \in \{0, 1, \dots, r-1\},
    \label{eq:autovectores_U}
\end{equation}
son autovectores de $U_{x,N}$ con autovalores $\lambda_s = e^{2\pi i s / r}$ \cite{Nielsen2024}.\\

La preparación directa de los $|u_s\rangle$ requeriría conocer $r$ a priori, precisamente la cantidad que se desea determinar. Esta dificultad se evita mediante una propiedad estructural: el estado computacional $|1\rangle$ (entendido como la codificación binaria de $L$ bits del número 1), que es el estado inicial natural del registro de trabajo, admite la descomposición
\[
    |1\rangle = \frac{1}{\sqrt{r}}\sum_{s=0}^{r-1}|u_s\rangle,
\]
es decir, $|1\rangle$ es una superposición uniforme de todos los autovectores $|u_s\rangle$ \cite{Nielsen2024}. Como consecuencia, al preparar el registro de trabajo en $|1\rangle$ y ejecutar la QPE con $t = 2L + 1 + \lceil\log(2+1/2\epsilon)\rceil$ qubits en el registro de conteo, el algoritmo realiza la estimación de fase simultáneamente sobre todos los autovectores $|u_s\rangle$. Tras la medición se obtiene una estimación de la fase $\varphi \approx s/r$ con precisión de $2L+1$ bits, para un $s$ aleatorio uniformemente distribuido y con probabilidad de al menos $(1-\epsilon)/r$ \cite{Nielsen2024}. El esquema completo se representa en la Figura~\ref{fig:busqueda_orden}.\\

\begin{figure}[htbp]
    \centering
    \includegraphics[width=0.9\textwidth]{Imagenes/Busqueda_Orden.png}
    \caption{Circuito cuántico para el algoritmo de búsqueda de orden. Este circuito también se puede utilizar para la factorización \cite{Nielsen2024}.}
    \label{fig:busqueda_orden}
\end{figure}

\textbf{El cuello de botella: exponenciación modular}\\

La viabilidad técnica del algoritmo de Shor descansa en la implementación eficiente de la secuencia de operaciones controladas $c$-$U_{x,N}^{2^j}$ que actúan durante la estimación de fase. Su efecto combinado equivale a multiplicar el contenido del segundo registro por la exponencial modular $x^z \pmod{N}$, donde $z$ es el contenido del primer registro. El desarrollo detallado de esta estrategia, incluyendo los casos en que la búsqueda de orden falla, se presenta en el Apéndice~\ref{app:exponenciacion-modular}.\\

\textbf{Origen de la ventaja cuántica}\\

La exponenciación modular por sí sola no es más rápida en un computador cuántico que en uno clásico: ambos la realizan en $\mathcal{O}(L^3)$ operaciones. La ventaja cuántica radica en cómo se utiliza esta operación.
\begin{itemize}
    \item Clásicamente, hallar $r$ requiere evaluar $x, x^2, x^3, \ldots$ de forma secuencial hasta encontrar $x^r \equiv 1 \pmod{N}$, lo que en el peor caso requiere $\mathcal{O}(r)$ evaluaciones, exponencial en $L$.
    \item Cuánticamente, gracias a la superposición, el circuito evalúa $x^z \pmod{N}$ para todos los valores de $z$ simultáneamente en una sola ejecución; la $\mathrm{QFT}^{\dagger}$ extrae la periodicidad $r$ de las fases globales del estado resultante sin enumerar caso por caso. Esta es la esencia de la aceleración exponencial del algoritmo de Shor.
\end{itemize}

\textbf{Fracciones continuas}\\

La extracción de $r$ a partir de la fase medida se garantiza mediante la teoría clásica de fracciones continuas:
\begin{teorema}[Convergencia de fracciones continuas]
    \label{thm:fracciones_continuas}
    Supóngase que $s/r$ es un número racional tal que
    \begin{equation}
        \left| \frac{s}{r} - x \right| \leq \frac{1}{2r^2}.
        \label{eq:fracciones_continuas}
    \end{equation}
    Entonces $s/r$ es un convergente de la expansión en fracción continua de $x$.
\end{teorema}
Aplicado a la medición $x = y/2^t$, este teorema implica que si el número de qubits de conteo cumple $t = 2L + 1$, la condición de convergencia se satisface y $r$ puede hallarse en tiempo $\mathcal{O}(L^3)$. La descripción detallada del procedimiento de extracción —expansión en fracción continua, recursión de los convergentes (Ecs.~\ref{eq:fraccion_continua_expansion} y~\ref{eq:convergentes}) y verificación clásica del candidato— se presenta en el Apéndice~\ref{app:fracciones-continuas-teoria} y se retoma en el contexto del trabajo experimental en Sección~\ref{sec:postprocesamiento}.\\

El procedimiento completo de la búsqueda cuántica de orden se formaliza en el Algoritmo~\ref{alg:order_finding}.


\begin{algorithm}[H]
    \caption{Búsqueda cuántica de orden \cite{Nielsen2024}}
    \label{alg:order_finding}
    \begin{algorithmic}[1]
        \Require (1) Una caja negra $U_{x,N}$ que realiza la transformación $|j\rangle|k\rangle \to |j\rangle|x^j k \bmod N\rangle$ para $x$ coprimo con el número de $L$ bits $N$, (2) $t = 2L + 1 + \lceil \log(2 + \frac{1}{2\epsilon}) \rceil$ qubits inicializados en $|0\rangle$, y (3) $L$ qubits inicializados en el estado $|1\rangle$.
        \Ensure El menor entero $r > 0$ tal que $x^r \equiv 1 \pmod{N}$.
        \Statex \textbf{Tiempo de ejecución:} $\mathcal{O}(L^3)$ operaciones. Éxito con probabilidad $\mathcal{O}(1)$.
        \Statex
        \Statex \textbf{Procedimiento:}
        \State \textit{Estado inicial:} $|0\rangle^{\otimes t}|1\rangle$ \hfill \textit{estado inicial}
        \State \textit{Crear superposición:}
        \Statex \hspace{2em} $\to \dfrac{1}{\sqrt{2^t}} \displaystyle\sum_{j=0}^{2^t-1} |j\rangle|1\rangle$ \hfill \textit{crear superposición}
        \State \textit{Aplicar $U_{x,N}$:}
        \Statex \hspace{2em} $\to \dfrac{1}{\sqrt{2^t}} \displaystyle\sum_{j=0}^{2^t-1} |j\rangle|x^j \bmod N\rangle$ \hfill \textit{aplicar $U_{x,N}$}
        \Statex \hspace{2em} $\approx \dfrac{1}{\sqrt{r\,2^t}} \displaystyle\sum_{s=0}^{r-1}\sum_{j=0}^{2^t-1} e^{2\pi i s j/r} |j\rangle|u_s\rangle$ \hfill \textit{expandir en autovectores}
        \State \textit{Aplicar QFT inversa al primer registro:}
        \Statex \hspace{2em} $\to \dfrac{1}{\sqrt{r}} \displaystyle\sum_{s=0}^{r-1} |\widetilde{s/r}\rangle|u_s\rangle$ \hfill \textit{aplicar $\mathrm{QFT}^{\dagger}$}
        \State \textit{Medir el primer registro:}
        \Statex \hspace{2em} $\to \widetilde{s/r}$ \hfill \textit{medir primer registro}
        \State \textit{Aplicar algoritmo de fracciones continuas} a $\widetilde{s/r}$ para extraer $r$.
        \Statex \hspace{2em} $\to r$ \hfill \textit{fracciones continuas}
    \end{algorithmic}
\end{algorithm}

\textbf{Implicaciones para la era NISQ}\\

En implementaciones pedagógicas orientadas a la evaluación en plataformas NISQ como IBM Quantum, el caso $N = 15$ es la cota menor práctica actual. Incluso para este caso, las implementaciones estándar de la exponenciación modular son transpiladas a miles de compuertas CNOT, aniquilando la fidelidad del estado mucho antes de la ejecución de la $\mathrm{QFT}^{\dagger}$. Por ende, en este trabajo se emplean arquitecturas altamente optimizadas con conocimiento previo del estado de la subrutina (compilación dependiente de $a$) \cite{villarreal}, minimizando el \textit{gate count} de dos qubits para habilitar el estudio experimental de fidelidad distinguible del ruido térmico.

\subsection{Reducción clásica: de la factorización a la búsqueda de orden}
\label{sec:reduccion_clasica}

Una vez establecidos los bloques fundamentales que intervienen en el algoritmo de Shor, queda por establecer la conexión entre la factorización prima de un entero compuesto y dicha maquinaria. El problema de factorización es equivalente al problema de búsqueda de orden, en el sentido de que un algoritmo polinomial para la búsqueda de orden se transforma directamente en un algoritmo polinomial para la factorización \cite{Nielsen2024}.\\

La reducción procede en dos pasos. Primero, se demuestra que es posible calcular un factor de $N$ a partir de cualquier solución no trivial $x \neq \pm 1 \pmod{N}$ de la ecuación $x^2 \equiv 1 \pmod{N}$. Segundo, se demuestra que un $y$ elegido al azar y coprimo con $N$ tiene, con probabilidad acotada inferiormente, un orden $r$ par que satisface $y^{r/2} \neq \pm 1 \pmod{N}$, de modo que $x \equiv y^{r/2} \pmod{N}$ es una solución no trivial de la ecuación anterior. Estos resultados se formalizan en los siguientes teoremas \cite{Nielsen2024}.\\

\begin{teorema}[Reducción de factorización a búsqueda de orden]
    \label{thm:red_fac}
    Supóngase que $N$ es un número compuesto de $L$ bits, y $x$ es una solución no trivial a la ecuación $x^2 = 1 \pmod{N}$ en el rango $1 \leq x \leq N$, es decir, ni $x = 1 \pmod{N}$ ni $x = N - 1 \equiv -1 \pmod{N}$. Entonces al menos uno de
    \begin{equation}
        \gcd(x - 1, N) \quad \text{y} \quad \gcd(x + 1, N),
    \end{equation}
    es un factor no trivial de $N$ que puede ser calculado utilizando $\mathcal{O}(L^3)$ operaciones.
    \label{eq:factores}
\end{teorema}

La demostración se apoya en la identidad $(a^{r/2}-1)(a^{r/2}+1) \equiv 0 \pmod N$, que implica que $N$ divide al producto. Si ninguna de las dos diferencias es múltiplo de $N$ (es decir, $a^{r/2} \not\equiv \pm 1 \pmod N$), entonces cada una debe compartir un factor primo con $N$. El cálculo del máximo común divisor ($gcd$) se realiza eficientemente con el algoritmo de Euclides en tiempo polinomial.

\begin{teorema}[Probabilidad de la factorización]
    \label{thm:prob_fac}
    Supóngase que $N = p_1^{\alpha_1} \dots p_m^{\alpha_m}$ es la factorización prima de un entero positivo compuesto impar. Sea $x$ un entero elegido uniformemente al azar, sujeto a los requisitos de que $1 \leq x \leq N - 1$ y $x$ es coprimo con $N$. Sea $r$ el orden de $x$ módulo $N$. Entonces
    \begin{equation}
        p(r \text{ es par y } x^{r/2} \neq -1 \pmod{N}) \geq 1 - \frac{1}{2^m}.
    \end{equation}
    \label{eq:probabilidad}
\end{teorema}
\vspace{-0.8cm}
Los Teoremas~\ref{thm:red_fac} y~\ref{thm:prob_fac} se combinan para producir un algoritmo que, con alta probabilidad, devuelve un factor no trivial de cualquier número compuesto $N$. Todos los pasos pueden realizarse de manera eficiente en una computadora clásica salvo (por lo que se sabe hoy en día) la subrutina de búsqueda de orden, que clásicamente exige recursos exponenciales. Repitiendo el procedimiento se obtiene una factorización prima completa de $N$ \cite{Nielsen2024}.\\

El procedimiento completo del algoritmo de Shor se formaliza en el Algoritmo~\ref{alg:shor}.

\begin{algorithm}[H]
    \caption{Reducción de la factorización a la búsqueda de orden \cite{Nielsen2024}}
    \label{alg:shor}
    \begin{algorithmic}[1]
        \Require Un número compuesto $N$
        \Ensure Un factor no trivial de $N$
        \Statex \textbf{Tiempo de ejecución:} $\mathcal{O}((\log N)^3)$ operaciones. Éxito con probabilidad $\mathcal{O}(1)$.
        \Statex
        \If{$N$ es par}
        \State \Return el factor $2$
        \EndIf
        \State Determinar si $N = a^b$ para enteros $a \geq 1$, $b \geq 2$. Si es así, \Return el factor $a$.
        \State Elegir aleatoriamente $x$ en el rango $1 \leq x \leq N-1$.
        \If{$\gcd(x, N) > 1$}
        \State \Return el factor $\gcd(x, N)$
        \EndIf
        \State Usar la \textbf{subrutina cuántica de búsqueda de orden} (Algoritmo~\ref{alg:order_finding}) para hallar el orden $r$ de $x$ módulo $N$.
        \If{$r$ es par \textbf{y} $x^{r/2} \not\equiv -1 \pmod{N}$}
        \State Calcular $\gcd(x^{r/2} - 1, N)$ y $\gcd(x^{r/2} + 1, N)$.
        \State \Return el factor no trivial encontrado.
        \Else
        \State El algoritmo falla. Repetir desde el paso 9.
        \EndIf
    \end{algorithmic}
\end{algorithm}

Los pasos 1-3 y 4 del algoritmo devuelven un factor o, de lo contrario, aseguran que $N$ es un entero impar con más de un factor primo. Estos pasos se realizan en $\mathcal{O}(1)$ y $\mathcal{O}(L^3)$ operaciones, respectivamente. Los pasos 5-8 devuelven un factor o producen un elemento $x$ de $\{0, 1, 2, \dots, N - 1\}$ elegido al azar. El paso 9 llama a la subrutina de búsqueda de orden, calculando el orden $r$ de $x$ módulo $N$. El paso 10 completa el algoritmo: el Teorema~\ref{thm:prob_fac} garantiza que con probabilidad de al menos un medio $r$ es par y $x^{r/2} \neq -1 \pmod N$, y entonces el Teorema~\ref{thm:red_fac} garantiza que $\gcd(x^{r/2} - 1, N)$ o $\gcd(x^{r/2} + 1, N)$ es un factor no trivial de $N$.

\subsection{Ejemplo guiado: factorización de \texorpdfstring{$N=15$}{N=15}}
\label{sec:ejemplo_n15}

Para aterrizar el formalismo previo en el caso de estudio del trabajo, se recorre a continuación el Algoritmo~\ref{alg:shor} aplicado a $N = 15$, identificando explícitamente los valores numéricos en cada paso.\\

\begin{enumerate}
\item \textit{Comprobar si $N$ es par:} como $N = 15$ es impar, el algoritmo no termina aquí.

\item \textit{Comprobar si $N$ es una potencia perfecta:} se examinan las posibles representaciones $N = a^b$ con $a \geq 2$ y $b \geq 2$. Para $b = 2$, $\sqrt{15} \approx 3{,}87 \notin \mathbb{Z}$; para $b = 3$, $\sqrt[3]{15} \approx 2{,}47 \notin \mathbb{Z}$; valores mayores de $b$ requerirían $a < 2$. Por lo tanto, $N = 15$ no es potencia perfecta.

\item \textit{Elegir una base $a$ coprima con $N$:} el grupo multiplicativo $\mathbb{Z}_{15}^{*}$ está formado por los enteros del intervalo $[1, 14]$ coprimos con 15:
\[
    \mathbb{Z}_{15}^{*} = \{1, 2, 4, 7, 8, 11, 13, 14\}.
\]
La base $a = 1$ es trivial (orden $r = 1$, no aprovechable); las demás constituyen candidatos válidos. En este trabajo se evalúan experimentalmente las bases representativas $a \in \{4, 7, 11, 14\}$. Para fijar el ejemplo, se toma $a = 7$.

\item \textit{Verificar coprimalidad:} como $\gcd(7, 15) = 1$, no se obtiene un factor por esta vía y se procede a la subrutina cuántica.

\item \textit{Subrutina cuántica de búsqueda de orden:} el registro de trabajo se inicializa en $|1\rangle$, el registro de conteo en $|0\rangle^{\otimes t}$ con $t = 2L + 1 = 9$ (puesto que $L = \lceil \log_2 15 \rceil = 4$), y se ejecuta el Algoritmo~\ref{alg:order_finding} con el operador $U_{7,15}$. La subrutina devuelve, tras el post-procesamiento por fracciones continuas, el orden $r = 4$. Esto se verifica calculando explícitamente las potencias de 7 módulo 15:
\begin{align*}
    7^1 & \equiv 7  \pmod{15},                       \\
    7^2 & = 49 \equiv 4 \pmod{15},                   \\
    7^3 & = 7 \cdot 4 = 28 \equiv 13 \pmod{15},      \\
    7^4 & = 7 \cdot 13 = 91 = 6 \cdot 15 + 1 \equiv 1 \pmod{15}.
\end{align*}
El menor entero positivo que satisface $7^r \equiv 1 \pmod{15}$ es, en efecto, $r = 4$.

\item \textit{Comprobar las condiciones del Teorema~\ref{thm:red_fac}:} el orden $r = 4$ es par. Se evalúa entonces $7^{r/2} = 7^2 \equiv 4 \pmod{15}$. Como $4 \not\equiv -1 \equiv 14 \pmod{15}$, se cumple la condición $x^{r/2} \not\equiv -1 \pmod N$ y la rama exitosa del algoritmo se activa.

\item \textit{Calcular los factores:} se evalúan los máximos comunes divisores
\begin{align*}
    \gcd(7^{r/2} - 1,\; 15) & = \gcd(48, 15) = 3, \\
    \gcd(7^{r/2} + 1,\; 15) & = \gcd(50, 15) = 5.
\end{align*}
Ambos resultan ser factores no triviales de $N$, lo que conduce a la factorización
\[
    15 = 3 \times 5.
\]
\end{enumerate}

Este recorrido ilustra el flujo abstracto del algoritmo. La traducción de la subrutina cuántica del Paso~5 a un circuito ejecutable en hardware IBM Quantum, incluyendo la compilación dependiente de la base $a$, la elección del número efectivo de qubits y la estructura de la QFT, se desarrolla en el Capítulo~5.

\section{El contexto NISQ y la arquitectura de IBM Quantum}
\label{sec:contexto_nisq}

Antes de abordar las estrategias de transpilación y supresión de errores, es imprescindible comprender el entorno de hardware en el que se ejecutarán los circuitos del algoritmo de Shor. Esta sección presenta un panorama del régimen computacional actual, la plataforma de hardware empleada en este trabajo y las restricciones físicas que determinan la viabilidad de los experimentos.

%% ----------------------------------------------------------------
%% 4.4.1  LA ERA NISQ
%% ----------------------------------------------------------------
\subsection{La era NISQ: qubits físicos frente a qubits lógicos}

El término \textit{Noisy Intermediate-Scale Quantum} (NISQ) fue introducido por Preskill \cite{Preskill} para designar la generación actual de procesadores cuánticos, caracterizada por dos rasgos definitorios: (i) un número intermedio de qubits (de decenas a unos pocos miles) y (ii) la \textit{ausencia de corrección cuántica de errores} (QEC) a escala completa. En este régimen, cada qubit que opera en el chip es un \textit{qubit físico} expuesto directamente al ruido del entorno, en contraste con el \textit{qubit lógico} de la computación cuántica tolerante a fallos (\textit{Fault-Tolerant Quantum Computing}, FTQC), donde la información se codifica redundantemente en cientos o miles de qubits físicos mediante códigos correctores de errores \cite{Roffe2019}.\\

Las limitaciones centrales de los dispositivos NISQ pueden resumirse en tres ejes que operan simultáneamente \cite{Preskill, Lammers2025}:

\begin{enumerate}
    \item \textit{Decoherencia y tiempos de coherencia finitos:} como se formalizó en la Sección~\ref{sec:nisq}, los qubits superconductores poseen tiempos de relajación $T_1$ y de desfase $T_2$ del orden de decenas a cientos de microsegundos \cite{AbuGhanem2025, errores_lectura, Preskill}. Esto impone un límite superior estricto a la profundidad del circuito: una vez que el tiempo de ejecución total del circuito transpilado se aproxima a $T_2$, la señal cuántica se degrada irreversiblemente \cite{Preskill}.

    \item \textit{Errores intrínsecos de las compuertas:} las compuertas de un qubit alcanzan fidelidades $\gtrsim 99.9\%$ en los procesadores actuales de IBM \cite{compuertas_CNOT,AbuGhanem2025}, pero las compuertas de dos qubits (la compuerta ECR nativa, equivalente a CNOT con rotaciones locales)\cite{IBM_ECRGate2025} presentan tasas de error típicas del orden de $\epsilon_{2q} \sim 10^{-3}$ en el mejor de los casos \cite{AbuGhanem2025, DeCusatis2021,errores_lectura}. El error global del circuito se acumula multiplicativamente: para un circuito con $N_g$ compuertas de dos qubits, la fidelidad del estado de salida decae aproximadamente como $(1 - \epsilon_{2q})^{N_g}$, lo que hace que circuitos con unos pocos cientos de compuertas ECR produzcan distribuciones de medición esencialmente indistinguibles del ruido \cite{szasz2023numerical}. Este comportamiento multiplicativo ha sido confirmado tanto en contextos de síntesis de circuitos \cite{szasz2023numerical}
          como en modelos analíticos de acumulación de errores \cite{deburuaga2024fidelity}.

    \item \textit{Errores de lectura (\textit{readout}):} la medición de cada qubit también introduce errores asimétricos modelados por la matriz de confusión (Ec.~\ref{eq:matriz_confusion}), con tasas típicas de $1$--$5\%$ por qubit \cite{DeCusatis2021}.
\end{enumerate}

%% ----------------------------------------------------------------
%% 4.4.2  ARQUITECTURA DE IBM QUANTUM
%% ----------------------------------------------------------------
\subsection{Arquitectura de IBM: qubits superconductores transmon}

Los procesadores de IBM Quantum se basan en qubits superconductores de tipo \textit{transmon} (\textit{transmission-line shunted plasma oscillation qubit}), que constituyen la plataforma de hardware más extendida en la industria cuántica actual \cite{AbuGhanem2025}.\\

El principio físico fundamental es la cuantización de la carga y el flujo magnético en un circuito superconductor que incorpora una \textit{unión Josephson}: un elemento no lineal cuya energía depende de la diferencia de fase superconductora a través de la unión. Esta no linealidad permite que los dos niveles de energía más bajos del circuito ($|0\rangle$ y $|1\rangle$) estén separados por una frecuencia de transición $\omega_{01}$ distinta de las transiciones superiores, de modo que el sistema se comporta efectivamente como un sistema de dos niveles (qubit), a diferencia de un oscilador armónico lineal donde todos los niveles están igualmente espaciados \cite{AbuGhanem2025, Nielsen2024}.\\

\begin{figure}[htbp]
    \centering
    \includegraphics[width=0.7\textwidth]{Imagenes/transmon.png}
    \caption{Imágenes de un cúbit transmon sintonizable en frecuencia. (a) El transmon completo consta de una capacitancia grande (en forma de "$+$") en paralelo con un SQUID conectado a tierra. En (b) y (c) se muestra un acercamiento del SQUID y una unión Josephson, respectivamente. \cite{transmon_ieee2023}}
    \label{fig:transmon}
\end{figure}

Las operaciones (compuertas) sobre los qubits transmon se realizan mediante \textit{pulsos de microondas} calibrados a la frecuencia de resonancia de cada qubit ($\omega_{01} \sim 4$--$5\;\text{GHz}$) \cite{AbuGhanem2025}. Una rotación arbitraria de un qubit se logra modulando la amplitud, la fase y la duración del pulso. La base de compuertas nativas (\textit{basis gates}) para algunos de los procesadores IBM Eagle/Heron es: \textit{RZ($\lambda$)}, \textit{SX}, \textit{X}, \textit{ECR}(equivalente a CNOT con rotaciones locales), junto con otras como \textit{CZ}, \textit{ID}, \textit{DELAY}, \textit{MEASURE}, etc \cite{IBMcloud2025}.
%% ----------------------------------------------------------------
%% 4.4.3  TOPOLOGÍA HEAVY-HEX Y CONECTIVIDAD
%% ----------------------------------------------------------------
\subsection{Topología Heavy-Hex y el cuello de botella de la conectividad}

En un procesador cuántico ideal con conectividad \textit{all-to-all}, cualquier compuerta de dos qubits podría aplicarse directamente entre dos qubits cualesquiera del chip. Sin embargo, los procesadores superconductores de IBM emplean una topología de conectividad restringida denominada \textit{Heavy-Hex} \cite{AbuGhanem2025, IBMcloud2025}. En esta arquitectura, los qubits se disponen en los vértices y aristas de una red hexagonal pesada (\textit{heavy hexagonal lattice}), donde cada qubit está acoplado a un máximo de dos o tres vecinos según su posición en la red. Los procesadores de la familia Eagle (127 qubits), Heron (133 qubits) y los más recientes de la hoja de ruta de IBM utilizan variantes de esta topología \cite{AbuGhanem2025}.\\

La consecuencia operativa más importante de esta conectividad restringida es la necesidad de insertar \textit{compuertas SWAP} durante la transpilación. Cuando el circuito lógico requiere una compuerta de dos qubits entre qubits que no son vecinos físicos en la topología, el transpilador debe mover la información cuántica a través de la red insertando operaciones SWAP, cada una de las cuales equivale a tres compuertas ECR nativas \cite{compuertas_SWAP}:
\begin{equation}
    \text{SWAP} = \text{ECR}_{ij}\;\text{ECR}_{ji}\;\text{ECR}_{ij},
    \label{eq:swap_ecr}
\end{equation}
triplicando el número de compuertas de dos qubits ruidosas por cada SWAP insertado \cite{villarreal, IBMcloud2025}. Este fenómeno constituye el \textit{principal cuello de botella físico} para la ejecución de algoritmos profundos como Shor: un circuito lógico con $N_{\text{ECR}}$ compuertas de dos qubits puede inflarse a $3N_{\text{SWAP}} + N_{\text{ECR}}$ compuertas de dos qubits tras la transpilación, donde $N_{\text{SWAP}}$ depende de la distancia topológica entre los qubits involucrados y de la calidad del algoritmo de ruteo empleado.\\

\begin{figure}[htbp]
    \centering
    \includegraphics[width=\textwidth]{Imagenes/fez_ibm.png}
    \caption{Topología Heavy-Hex del procesador cuántico IBM \textit{ibm\_fez} (panel derecho) de la familia Heron r2 con 156 qubits. A la izquierda los datos de calibración mas importantes \cite{IBM20251}.}
    \label{fig:heavy-hex-topology-ibm-fez}
\end{figure}

\section{Transpilación y algoritmo SABRE}
\label{sec:transpilacionysabre}

La ejecución exitosa de algoritmos cuánticos profundos, como el algoritmo de Shor, en arquitecturas NISQ requiere una adaptación exhaustiva del circuito lógico a las restricciones físicas del hardware y la aplicación de técnicas simultáneas de supresión y mitigación de ruido. El marco de desarrollo Qiskit proporciona las abstracciones de transpilación, supresión y mitigación que posibilitan la transición de la teoría a la experimentación real \cite{IBMqiskit, IBMcloud2025}. A continuación, se detalla el formalismo de estas estrategias que fueron implementadas sobre los procesadores IBM.

\subsection{El proceso de transpilación y niveles de optimización}

La ejecución de algoritmos cuánticos en hardware real de IBM Quantum, requiere la transformación de una representación de alto nivel y abstracta del circuito a un conjunto de instrucciones físicas que respete las limitaciones del \textit{backend} objetivo. Este proceso es llevado a cabo por el transpilador cuántico (transpiler), una herramienta fundamental que actúa como compilador, adaptando la descripción lógica del algoritmo a la arquitectura física subyacente \cite{IBM_transpile_2025}. Dicha adaptación es indispensable para sortear restricciones como la conectividad limitada entre qubits (topología del chip), el conjunto nativo de compuertas (\textit{basis gates}) y las características de calibración específicas de cada dispositivo \cite{IBM_defaults_2025}.\\

El proceso de transpilación en Qiskit se estructura en un \textit{pipeline} compuesto por seis etapas secuenciales, cada una con un propósito específico para transformar el circuito de entrada en uno optimizado y compatible con el hardware \cite{IBM_stages_2025}. Estas etapas son:

\begin{enumerate}
    \item \textit{Inicialización (init):} configura el transpilador con los parámetros iniciales y las propiedades del \textit{backend} objetivo.
    \item \textit{Asignación lógica a física (layout):} establece un mapeo inicial entre los qubits virtuales (lógicos) del circuito y los qubits físicos del dispositivo. Este mapeo busca, en sus etapas iniciales, minimizar la distancia para futuras operaciones.
    \item \textit{Enrutamiento (routing):} dado que las compuertas de dos qubits solo pueden ejecutarse en pares físicamente conectados, esta etapa inserta compuertas SWAP para enrutar los qubits lógicos a posiciones donde las compuertas no locales puedan ser aplicadas.
    \item \textit{Traducción (translation):} convierte todas las compuertas del circuito a un conjunto de compuertas base (\textit{basis gates}) que son físicamente realizables por el \textit{backend}, como las mencionadas (\texttt{rz}, \texttt{sx}, \texttt{x}, \texttt{ecr}).
    \item \textit{Optimización (optimization):} aplica una serie de transformaciones para reducir el número de compuertas, la profundidad del circuito y otras métricas, sin alterar la función unitaria del algoritmo.
    \item \textit{Programación (scheduling):} asigna los tiempos de ejecución de cada instrucción en el hardware, considerando efectos como la calibración dinámica (e.g., compuertas RZ virtuales) y los tiempos de relajación de los qubits.
\end{enumerate}

Para controlar el balance entre el tiempo de compilación y la calidad de la optimización, el transpilador ofrece cuatro niveles de optimización, configurados mediante el parámetro \texttt{optimization\_level}, cuyo valor puede ser 0, 1, 2 o 3 \cite{IBM_optimization_2025}. Cada nivel activa un conjunto diferente de rutinas de optimización, impactando significativamente en el número final de compuertas, la profundidad del circuito y el tiempo de compilación. En la Tabla ~\ref{tab:opt_levels} se presenta un resumen comparativo de estos niveles.

\begin{table}[htbp]
    \centering
    \caption{Niveles de optimización en el transpilador de Qiskit \cite{IBM_transpile_2025}.\\}
    \label{tab:opt_levels}
    \begin{tabular}{p{0.7cm} p{6.3cm} p{1.6cm} p{1.6cm} p{1.6cm}}
        \toprule
        \textbf{Nivel} & \textbf{Transformaciones clave activadas}                                                                                                                                                                                                             & \textbf{Número-compuertas} & \textbf{Depth}          & \textbf{Tiempo-compilación} \\
        \midrule
        0              & Solo etapas de traducción y enrutamiento mínimos. No se aplican optimizaciones de reducción.                                                                                                                                                          & Alto                       & Alta                    & Mínimo                      \\
        1              & Optimizaciones ligeras: cancelación de compuertas adyacentes (e.g., \texttt{Not} seguida de \texttt{Not}), y plegado de compuertas.                                                                                                                   & Reducción moderada         & Reducción moderada      & Bajo                        \\
        2              & Optimizaciones medias: incluye las del nivel 1, más consolidación de bloques unitarios (\texttt{ConsolidateBlocks}) y cancelación conmutativa (\texttt{CommutativeCancellation}). Re-síntesis de bloques para usar compuertas nativas más eficientes. & Reducción significativa    & Reducción significativa & Moderado                    \\
        3              & Optimizaciones exhaustivas: incluye todas las del nivel 2, junto con técnicas de re-enrutamiento y optimización dependientes del ruido (\textit{noise-adaptive}) como \texttt{SABRE} con parámetros más agresivos.                                    & Reducción máxima           & Reducción máxima        & Alto                        \\
        \bottomrule
    \end{tabular}
\end{table}

Dentro de las etapas de asignación (\textit{layout}) y enrutamiento (\textit{routing}), el algoritmo SABRE juega un papel crucial, especialmente en los niveles de optimización más altos. Su eficacia, como se verá en la siguiente sección, radica en la capacidad de realizar un \textit{layout} inicial adaptativo y un proceso de enrutamiento que considera las calibraciones en tiempo real del dispositivo, buscando un equilibrio entre la inserción de compuertas SWAP y la fidelidad de la ejecución \cite{IBM_SABRE_2025}.\\

\begin{figure}[htbp]
    \centering
    \includegraphics[width=\textwidth]{Imagenes/ejemplo_transpilacion.png}
    \caption{Ejemplo de transpilación de un circuito cuántico simple en Qiskit para un procesador IBM con topología Heavy-Hex \cite{IBM_transpile_2025}.
        A la izquierda se muestra el circuito lógico original.
        A la derecha aparece el circuito transpilado, donde se han insertado compuertas nativas del hardware y se han mapeado los qubits lógicos $q_0$ y $q_1$ a los qubits físicos 58 y 57, respectivamente.
        Nótese cómo incluso un circuito sencillo requiere varias compuertas de dos qubits (\textit{ECR}) y rotaciones adicionales para adaptarse a la conectividad restringida del procesador.}
    \label{fig:transpilacion-ejemplo-qiskit}
\end{figure}

En resumen, la transpilación es un proceso multi-etapa que adapta un algoritmo cuántico abstracto a las realidades físicas de un dispositivo NISQ. Los niveles de optimización de Qiskit proporcionan un mecanismo de control fino sobre el balance entre el esfuerzo computacional de compilación y la calidad del circuito resultante.

\subsection{Algoritmo SABRE}

En los procesadores superconductores de IBM, una compuerta de dos qubits sólo es físicamente realizable entre qubits unidos por una arista del grafo de acoplamiento del chip \cite{IBM_SABRE2026}. Esta restricción geométrica fractura la correspondencia entre el circuito lógico y el circuito físico: dos qubits que el algoritmo trata como vecinos en el espacio de Hilbert pueden estar separados por varios sitios sobre la red \textit{Heavy-Hex}. Cuando esto ocurre, el transpilador se ve obligado a transportar los estados cuánticos a través del cristal mediante compuertas SWAP, cada una descompuesta a su vez en varias compuertas entrelazantes nativas. Cada inserción amplifica simultáneamente la profundidad del circuito y la probabilidad acumulada de error, comprometiendo la fidelidad del estado final y volviendo prohibitivamente costosas las subrutinas largas como la exponenciación modular del algoritmo de Shor \cite{Li2019SABRE}.\\

El algoritmo SABRE (\textit{SWAP-based BidiREctional heuristic search}) \cite{Li2019SABRE} aborda este problema de transporte representando el hardware como un grafo de conectividad
\begin{equation}
    G=(V,E),
\end{equation}

cuyos vértices son los qubits físicos del procesador y cuyas aristas codifican los acoplamientos permitidos por el diseño del chip. Sobre este grafo se define un mapeo
\begin{equation}
    \pi:\{q_1,q_2,\dots,q_n\} \to \{Q_1,Q_2,\dots,Q_N\},
\end{equation}
que ubica cada qubit lógico del algoritmo en un sitio físico del cristal. A partir de la geometría del chip, SABRE precalcula una matriz de distancias $D[i][j]$ correspondiente al camino más corto entre los qubits $Q_i$ y $Q_j$ \cite{Li2019SABRE}. Físicamente, $D[i][j]$ es una métrica de coste experimental: estima el número mínimo de intercambios necesarios para que dos estados cuánticos puedan interactuar coherentemente.\\

La filosofía del método consiste en renunciar a una optimización global y elegir, en cada paso, el SWAP más conveniente mediante una heurística local. Para ello, el algoritmo organiza las compuertas entrelazantes del circuito según sus dependencias y aísla un \emph{frente de interacciones cuánticas inmediatas} $F$: el conjunto de operaciones que están listas lógicamente, pero cuya ejecución física depende del mapeo vigente. Si alguna compuerta de $F$ ya conecta sitios vecinos del cristal, se ejecuta de inmediato. En caso contrario, SABRE genera SWAP candidatos únicamente en torno a los qubits implicados en ese frente, evitando explorar reconfiguraciones subóptimas en regiones distantes del procesador \cite{Li2019SABRE}. Esta restricción geométrica concentra el esfuerzo allí donde el circuito tiene una demanda física inminente de conectividad.\\

La versión más simple de la heurística favorece el SWAP que minimiza la distancia total entre los pares de qubits del frente \cite{Li2019SABRE}:
\begin{equation}
    H_{\mathrm{basic}}=\sum_{g\in F} D\left[\pi(gate.q_1)\right]\left[\pi(gate.q_2)\right],
    \label{eq:sabre_basic}
\end{equation}

donde $\pi$ representa el mapeo resultante tras aplicar el intercambio y $q_1$, $q_2$ son los qubits lógicos involucrados en la compuerta $g$. Una disminución de $H_{\mathrm{basic}}$ indica que los estados que deben entrelazarse han quedado geométricamente más próximos sobre la red. Sin embargo, una elección óptima para el frente puede degradar la conectividad de las compuertas inmediatamente posteriores. Por ello SABRE incorpora una \emph{previsión topológica de acoplamientos} mediante un conjunto extendido $E$ que contiene los sucesores cercanos de $F$, así como un mecanismo que penaliza temporalmente reutilizar qubits recién intervenidos para impedir que las operaciones se acumulen en una misma región del cristal superconductor. La heurística completa se expresa como \cite{Li2019SABRE}:

\begin{align}
    H & = \max\big(decay(SWAP.q_1), decay(SWAP.q_2) \big) \notag \\
      & \quad \times \Bigg\{
    \frac{1}{|F|} \sum_{\text{gate} \in F}
    D\big[\pi(gate.q_1)\big]\big[\pi(gate.q_2)\big] \notag       \\
      & \qquad + W \cdot \frac{1}{|E|} \sum_{\text{gate} \in E}
    D\big[\pi(gate.q_1)\big]\big[\pi(gate.q_2)\big]
    \Bigg\}
    \label{eq:sabre_completa}
\end{align}

donde $q_1$ y $q_2$ son los qubits lógicos intercambiados por el SWAP candidato, $W$ ($0\leq W<1$) regula el peso relativo de la previsión, y $\mathrm{decay}(\cdot)$ penaliza los movimientos sobre qubits recién utilizados. Cada término admite una lectura física directa: el primero resuelve las interacciones urgentes del frente; el segundo evita comprometer las interacciones futuras al elegir un atajo demasiado apresurado; y el factor multiplicativo distribuye espacialmente la actividad entrelazante, impidiendo la concentración operacional sobre una misma región del cristal y favoreciendo SWAP paralelizables que no incrementan la profundidad total del circuito.\\

\begin{figure}[htbp]
    \centering
    \includegraphics[width=0.8\textwidth]{Imagenes/sabre_ejemplo.png}
    \caption{Ejemplo de búsqueda heurística basada en SWAP \cite{Li2019SABRE}.}
    \label{fig:sabre_ejemplo}
\end{figure}

Un rasgo particularmente potente de SABRE es que no se limita al transporte local de estados, sino que también optimiza la colocación inicial de los qubits sobre el cristal. El algoritmo explota la reversibilidad de la mecánica cuántica: ejecuta un recorrido hacia adelante, registra el mapeo final, transpila el circuito invertido tomando ese resultado como condición de contorno y, finalmente, refina con esa información el mapeo inicial del circuito original. Esta bidireccionalidad es esencial en circuitos como el de Shor, donde optimizar únicamente las primeras compuertas resulta insuficiente: la distribución física inicial debe anticipar las interacciones más exigentes que aparecerán a lo largo de la exponenciación modular \cite{Li2019SABRE}.\\

\begin{figure}[htbp]
    \centering
    \includegraphics[width=\textwidth]{Imagenes/sabre_mapeo.png}
    \caption{Actualización inicial del mapeo mediante la técnica de recorrido inverso \cite{Li2019SABRE}.}
    \label{fig:sabre_mapeo}
\end{figure}

En síntesis, SABRE reduce la sobrecarga de transporte cuántico mediante dos mecanismos geométricamente complementarios \cite{Li2019SABRE}: restringe la búsqueda a los SWAP físicamente relevantes para la conectividad inmediata, y selecciona entre ellos con una métrica de distancia que aproxima el coste experimental real, ya que menos intercambios implican menos compuertas entrelazantes y, por tanto, menor acumulación de decoherencia. Para esta investigación, dicha filosofía resulta especialmente apropiada: el cuello de botella en la implementación de Shor sobre hardware NISQ no es de naturaleza computacional, sino física. Cada compuerta adicional debilita la señal útil del estado cuántico antes de la medición, de modo que un algoritmo de enrutamiento que minimiza profundidad y número de operaciones entrelazantes incrementa directamente la viabilidad experimental del circuito sobre el procesador \cite{IBM_SABRE2026}.


% =========================================================================
% BLOQUE A — Sección 4.6 reducida (cuerpo principal)
% =========================================================================

\section{Estrategias de supresión en compilaciones complejas}
\label{sec:supresion_errores}

En los dispositivos NISQ, incluso los circuitos relativamente cortos del algoritmo de Shor acumulan errores con rapidez: las compuertas de dos qubits (ECR o CX) introducen errores coherentes a cada aplicación, y los intervalos en los que un qubit permanece inactivo lo exponen a la decoherencia del entorno. Sin acceso a corrección de errores completa, se recurre a técnicas de \textit{supresión de errores} que actúan durante la compilación y la ejecución del circuito con bajo costo adicional. En este trabajo se emplean dos técnicas complementarias: el \textit{Desacoplamiento Dinámico} (DD) y el \textit{Pauli Twirling} (PT). Ambas son compatibles con el flujo de transpilación (incluyendo SABRE) y se activan automáticamente desde Qiskit como configuraciones del \texttt{SamplerV2}.

\subsection{Desacoplamiento Dinámico (Dynamical Decoupling)}

Durante los tiempos de espera entre compuertas (\textit{idle times}), un qubit superconductor permanece expuesto a interacciones no controladas con su entorno: fluctuaciones térmicas, \textit{crosstalk} estático con qubits vecinos y acoplamiento con modos del resonador \cite{Ezzell2023, AbuGhanem2025}. Estas interacciones, agrupadas en un acoplamiento sistema--baño $H_{SB}$, hacen que la fase relativa del qubit evolucione de forma impredecible (\textit{dephasing}), degradando la coherencia que el algoritmo necesita para producir interferencia constructiva \cite{ViolaKnillLloyd1999}. El Desacoplamiento Dinámico no elimina físicamente este acoplamiento: lo \textit{promedia a cero} mediante la aplicación rápida de pulsos de control sobre el qubit.\\

La intuición proviene del \textit{eco de espín} de Hahn (1950) en resonancia magnética nuclear \cite{Hahn1950}: si durante un intervalo $\tau$ el qubit acumula una fase errónea $+\phi$ debido al ruido, un pulso $\pi$ alrededor del eje $x$ (compuerta $X$) invierte el signo de la evolución del ruido, de modo que durante el siguiente intervalo $\tau$ se acumula $-\phi$ y ambas contribuciones se cancelan a primer orden en $\tau$. La generalización moderna consiste en aplicar secuencias periódicas de pulsos durante el intervalo de inactividad, de manera que en el límite de pulsos rápidos el Hamiltoniano efectivo de error se promedia a un término que actúa exclusivamente sobre el baño y deja al qubit desacoplado del entorno a primer orden \cite{ViolaKnillLloyd1999, Ezzell2023}.\\

Sobre esta base, una de las secuencia mas relevante es XY4, propuesta originalmente por Maudsley en NMR \cite{Maudsley1986}, que alterna pulsos $\pi$ alrededor de los ejes $x$ e $y$:
\begin{equation}
    \text{XY4} \equiv Y - \tau - X - \tau - Y - \tau - X - \tau.
    \label{eq:xy4}
\end{equation}
XY4 destaca sobre alternativas más simples (como la secuencia $X$-$X$, equivalente a CPMG) porque responde a una propiedad física concreta: al rotar simultáneamente alrededor de dos ejes ortogonales del espacio de Pauli, XY4 simetriza errores tanto de defasamiento ($\sigma_z$) como de \textit{bit-flip} ($\sigma_x$) e interacciones cruzadas ($\sigma_y$), mientras que $X$-$X$ solo cancela los términos que anticonmutan con $X$ y deja sin atender los acoplamientos $\sigma_x$ estáticos \cite{Ezzell2023, Rahman2024}. XY4 es, por tanto, una secuencia universal de primer orden cuya residual escala como $\mathcal{O}(\tau^2)$.

\begin{figure}[htbp]
    \centering
    \includegraphics[width=0.8\textwidth]{Imagenes/DD_secuencia.png}
    \caption{Diagrama esquemático de secuencias DD. (a) Dos secuencias CPMG concatenadas. (b) Dos secuencias XY4 concatenadas. Todas las secuencias se muestran para un período de inactividad de \(4\tau\). CPMG usa solo puertas \(X\) (rectángulo azul), mientras que XY4 usa puertas $X$ e $Y$ (rectángulos verdes) \cite{Rahman2024}.}
    \label{fig:dd_secuencia}
\end{figure}

\subsection{Pauli Twirling (PT)}

A diferencia del DD, que actúa sobre los intervalos de inactividad, el Pauli Twirling aborda los errores que ocurren \textit{durante} la ejecución de las propias compuertas de dos qubits. Una compuerta real implementa una unitaria $\tilde{\mathcal{U}} = \mathcal{U}\cdot\mathcal{E}$ ligeramente desviada de la ideal $\mathcal{U}$, donde $\mathcal{E}$ representa una pequeña rotación con dirección y magnitud fijas. Al tratarse de un error \textit{coherente} ---repetible y no aleatorio--- las desviaciones se acumulan de forma constructiva, es decir, cuadráticamente con el número $N_g$ de compuertas, mientras que los errores estocásticos lo hacen linealmente \cite{Wallman2016, ibm_error_mitigation}. En circuitos profundos como los del algoritmo de Shor, este crecimiento cuadrático constituye el modo de fallo dominante.\\

El Pauli Twirling, también denominado \textit{compilación aleatorizada} \cite{Wallman2016}, transforma los errores coherentes en errores estocásticos. La operación consiste en envolver cada compuerta de dos qubits $\mathcal{U}$ entre un par de operadores de Pauli $(P_{\text{antes}}, P_{\text{después}})$ con $P \in \{I, X, Y, Z\}^{\otimes 2}$, elegidos aleatoriamente de un conjunto finito de combinaciones válidas. Tras promediar sobre múltiples instancias aleatorizadas del circuito, las componentes fuera de la diagonal del canal de error en la base de Pauli se cancelan y el ruido residual queda descrito por un canal estocástico de Pauli, cuya acumulación es lineal en $N_g$ \cite{CaiBenjamin2018, ibm_error_mitigation}.\\

La condición física que hace posible esta operación sin alterar el algoritmo es que el par de Pauli debe preservar la operación lógica de la compuerta original. Para una compuerta nativa de dos qubits como ECR (o su equivalente CNOT), los pares válidos satisfacen
\begin{equation}
    P_{\text{después}} \cdot \mathcal{U} \cdot P_{\text{antes}} = e^{i\phi}\, \mathcal{U},
    \label{eq:pt_preservation}
\end{equation}
para alguna fase global $\phi$, irrelevante para las probabilidades de medición \cite{ibm_error_mitigation}. Cada compuerta de dos qubits admite, por tanto, un conjunto finito de pares válidos del que el motor de ejecución selecciona aleatoriamente en cada instancia.\\

DD y PT actúan sobre fuentes de error diferentes: el primero protege a los qubits durante los intervalos de inactividad, el segundo aleatoriza los errores de las compuertas activas. El análisis detallado se reporta en el Apéndice~\ref{app:dd-formal} y en el~\ref{app:dd-formal}.

\section{Métricas de Evaluación}
\label{sec:metricas_evaluacion}

Para evaluar qué tan bien funciona un algoritmo cuántico ejecutado en hardware real, se necesitan indicadores numéricos que permitan comparar el resultado obtenido con el resultado esperado. En un mundo ideal, esta comparación se haría reconstruyendo completamente el estado cuántico final del sistema mediante Tomografía de Estado Cuántico (QST) \cite{Nielsen2024}. Sin embargo, para un sistema de $m$ qubits, QST requiere la medición de $4^m - 1$ observables independientes. En el caso del algoritmo de Shor para $N=15$ con $m = 9$ qubits de control, esto implica más de $2.6 \times 10^5$ mediciones distintas, una demanda que excede tanto los tiempos de coherencia de los dispositivos NISQ como los límites de acceso del IBM Quantum Open Plan \cite{IBM20251}.\\

Ante esta limitación, se adoptan métricas que trabajan directamente con las distribuciones de probabilidad obtenidas al medir en la base computacional, sin necesidad de reconstruir la matriz de densidad completa. En este trabajo se emplean dos métricas complementarias: la \textit{fidelidad de Hellinger}, que cuantifica la similitud estadística entre la distribución ideal y la experimental, y la \textit{probabilidad de éxito experimental} ($PST$), que mide la fracción de mediciones que caen en los resultados correctos del algoritmo.

\subsection{Fidelidad de Hellinger}

La idea intuitiva detrás de la fidelidad es simple: dados dos estados cuánticos, ¿qué tan parecidos son? Si un circuito ideal produce el estado $\rho$ y el hardware real produce el estado $\sigma$, la fidelidad cuántica mide qué tanto se ``solapan'' ambos estados. Formalmente, la fidelidad de Uhlmann--Jozsa se define como \cite{Fuchs1995, Nielsen2024}:
\begin{equation}
    F(\rho, \sigma) = \left( \mathrm{Tr} \sqrt{\sqrt{\rho}\,\sigma\,\sqrt{\rho}} \right)^2.
    \label{eq:fidelidad_uhlmann}
\end{equation}
Esta cantidad satisface $0 \leq F \leq 1$, donde $F = 1$ si y solo si los estados son idénticos, y $F = 0$ cuando son completamente ortogonales (distinguibles con certeza).\\

Cuando el estado de referencia es puro (como ocurre al comparar con una simulación ideal) la expresión se simplifica notablemente. Si $\rho = |\psi\rangle\langle\psi|$, la Ecuación~\ref{eq:fidelidad_uhlmann} se reduce a \cite{Fuchs1995, Nielsen2024}:
\begin{equation}
    F(|\psi\rangle\langle\psi|, \sigma) = \langle\psi|\sigma|\psi\rangle.
    \label{eq:fidelidad_estado_puro}
\end{equation}
Es decir, la fidelidad se convierte simplemente en la probabilidad de que, al medir el estado experimental $\sigma$, se obtenga el estado ideal $|\psi\rangle$. Esta interpretación conecta directamente la fidelidad con una cantidad físicamente medible. El problema práctico es que calcular la Ecuación~\ref{eq:fidelidad_estado_puro} requiere conocer la matriz de densidad $\sigma$, lo cual exige la tomografía completa que ya se descartó. La solución consiste en trabajar con las distribuciones de probabilidad que sí están disponibles: las frecuencias relativas de cada cadena binaria obtenidas tras muchas mediciones del circuito.\\

Sea $P = \{p_x\}$ la distribución de probabilidad ideal (obtenida del simulador sin ruido) y $Q = \{q_x\}$ la distribución experimental, donde cada $x$ es una cadena de $m$ bits. La \textit{fidelidad clásica} o \textit{fidelidad de Hellinger} se define como \cite{Nielsen2024, DeCusatis2021}:
\begin{equation}
    \mathcal{F}_H(P, Q) = \left( \sum_{x \in \{0,1\}^m} \sqrt{p_x \cdot q_x} \right)^2.
    \label{eq:fidelidad_clasica}
\end{equation}
Si bien la cantidad dentro del paréntesis, $\sum_x \sqrt{p_x q_x}$, no es una métrica en el sentido matemático, funciona perfectamente como medida de similitud acotada. El rango de $\mathcal{F}_H$ es $[0, 1]$. Un valor $\mathcal{F}_H \to 1$ indica que las distribuciones ideal y experimental son prácticamente indistinguibles; un valor $\mathcal{F}_H \to 0$ señala que el ruido ha destruido completamente la estructura de probabilidades del algoritmo.\\

La relación entre $\mathcal{F}_H$ y la fidelidad cuántica $F$ está dada por un resultado establecido en la teoría de la información cuántica. Fuchs demostró que la fidelidad clásica obtenida a partir de cualquier medición (formalizada como un POVM) siempre proporciona una cota superior de la fidelidad cuántica \cite{Fuchs1995, Wong2022}:
\begin{equation}
    F(\rho, \sigma) \leq \mathcal{F}_H(P, Q).
    \label{eq:fuchs_caves}
\end{equation}
Esta desigualdad tiene una consecuencia práctica importante: si la fidelidad de Hellinger calculada a partir de las mediciones experimentales resulta baja, la fidelidad cuántica real es necesariamente aún más baja. En otras palabras, $\mathcal{F}_H$ ofrece una estimación \textit{optimista} de la similitud entre los estados ideal y experimental.\\

\subsection{Probabilidad de Éxito Experimental ($PST$)}

Mientras que la fidelidad de Hellinger ofrece una visión global de la similitud entre distribuciones, la $PST$ responde a una pregunta más directa y operativa: ¿qué fracción de las mediciones produce un resultado que conduce a la factorización correcta?\\

El algoritmo de Shor funciona porque la estimación de fase cuántica (QPE) concentra la probabilidad de medición en un conjunto discreto de posiciones $\mathcal{Y}_{\text{ideal}}$ dentro del registro de conteo. Estas posiciones corresponden a las fases $\varphi_s = s/r$ de los eigenestados del operador de multiplicación modular $U_x$, y en el histograma aparecen como picos bien definidos en las posiciones \cite{Nielsen2024}:
\begin{equation}
    y_s = \left\lfloor \frac{s \cdot 2^m}{r} \right\rceil \pmod{2^m}, \quad s \in \{0, 1, \dots, r-1\}.
    \label{eq:picos_qpe}
\end{equation}
La $PST$ se define como la suma de las probabilidades experimentales medidas exclusivamente en estas posiciones \cite{villarreal}:
\begin{equation}
    PST = \sum_{y \in \mathcal{Y}_{\text{ideal}}} q_y.
    \label{eq:pst_definicion}
\end{equation}

En una simulación ideal sin ruido, toda la probabilidad se concentra en los picos del QPE, de modo que $PST = 1$. En el régimen ruidoso, la decoherencia y los errores de compuerta redistribuyen la probabilidad hacia cadenas binarias que no corresponden a ninguna fase válida, reduciendo la $PST$. Como señala \cite{Bagourd2026}, para que el algoritmo de Shor sea útil, la masa de probabilidad en las posiciones correctas debe superar de forma estadísticamente significativa la línea base que produciría un espectro completamente uniforme (aleatorio), la cual vale $b = 1/r$. En el régimen ideal, cada pico del QPE deposita al menos $4/\pi^2 \approx 0.405$ \cite{ibm_phase_estimation_procedure} de su masa dentro de la ventana de fracciones continuas, garantizando un éxito constante por ejecución \cite{Bagourd2026}. Cuando el ruido del hardware reduce la $PST$ por debajo de estos umbrales, el algoritmo pierde su ventaja cuántica y el número de repeticiones necesarias crece hasta hacer inviable la factorización. \\

La $PST$ complementa a $\mathcal{F}_H$ de forma natural: mientras la fidelidad mide la similitud global entre distribuciones, la $PST$ cuantifica directamente la utilidad algorítmica de los resultados experimentales. Es posible obtener una $\mathcal{F}_H$ moderada pero una $PST$ aceptable si el ruido redistribuye la probabilidad de forma relativamente uniforme entre los estados no deseados, manteniendo los picos correctos como los más probables. En la práctica, ambas métricas se reportan conjuntamente para ofrecer una evaluación completa del rendimiento del circuito en cada entorno de ejecución.


% =====================================================================
% 5.CAPITULO 5
% =====================================================================

\chapter{Diseño metodológico e implementación experimental}
\label{chap:metodologia}

Este capítulo documenta el diseño experimental empleado para evaluar el algoritmo de Shor con $N = 15$. La exposición se organiza en torno a las decisiones que estructuran el estudio: la delimitación del aporte original respecto al repositorio base, la selección clásica de parámetros, la construcción del circuito ideal, el protocolo experimental en tres fases concebido para aislar fuentes de error sucesivas, su transpilación, las técnicas de supresión evaluadas y, finalmente, el post-procesamiento clásico y las métricas adoptadas. Los detalles propios de la implementación, se trasladan al Apéndice~\ref{app:arquitectura-software-cap}.\\

\textbf{Nota.} A partir de aquí, la variable $x$ (base en la búsqueda de orden) del algoritmo de Shor se reemplaza por $a$, dado que en la práctica esta última es el parámetro de entrada efectivo del algoritmo.

% =====================================================================
% 5.1 DELIMITACIÓN DEL APORTE ORIGINAL
% =====================================================================
\section{Delimitación del aporte respecto al repositorio base}
\label{sec:aporte_original}

El repositorio proporcionado por \cite{villarreal} aportó la arquitectura general de clases y un flujo de ejecución gobernado por archivos de configuración. Sin embargo, dicha base fue concebida sobre versiones anteriores del SDK de Qiskit y contenía dependencias, nomenclaturas y elecciones metodológicas incompatibles o subóptimas respecto al ecosistema actual de IBM Quantum. El detalle técnico de cada migración (clases sustituidas, módulos reorganizados, conversión \textit{bitstring}-a-entero y verificaciones implementadas) se documenta en el Apéndice~\ref{app:migracion-qiskit}.

% =====================================================================
% 5.3 PRE-PROCESAMIENTO CLÁSICO Y SELECCIÓN DE PARÁMETROS
% =====================================================================
\section{Pre-procesamiento clásico y selección de parámetros}
\label{sec:preprocesamiento_parametros}

El algoritmo de Shor posee una naturaleza fundamentalmente híbrida: antes de invocar la subrutina cuántica, se ejecuta una serie de verificaciones clásicas que pueden resolver la factorización sin consumir recursos cuánticos. La envoltura clásica implementa esta lógica en un flujo secuencial, coherente con el procedimiento formal descrito en el Algoritmo~\ref{alg:shor} del marco teórico. Solo cuando ninguna de estas verificaciones produce resultados, se procede a la fase cuántica.

\subsection{Selección de la base $a$ y caracterización del grupo multiplicativo}

Para el caso de estudio $N = 15$, el grupo multiplicativo $\mathbb{Z}_{15}^* = \{2, 4, 7, 8, 11, 13, 14\}$ contiene siete elementos coprimos con $15$. Una rutina de pre-cálculo clásico determina, para cada base candidata, su orden multiplicativo $r$ ---evaluando sucesivamente $a^k \bmod N$ para $k = 1, 2, \ldots, N-1$ hasta encontrar el primer $k$ tal que $a^k \equiv 1 \pmod{N}$--- y verifica las condiciones de no-degeneración del Algoritmo~\ref{alg:shor}. Los resultados de esta caracterización, fundamentales para el diseño experimental, se resumen en la Tabla~\ref{tab:caracterizacion_a}.

\begin{table}[htbp]
    \centering
    \caption{Caracterización del grupo $\mathbb{Z}_{15}^*$: bases, órdenes y factores obtenidos.\\}
    \label{tab:caracterizacion_a}
    \begin{tabular}{cccccc}
        \toprule
        \textbf{Base $a$} & \textbf{Orden $r$} & \textbf{$r$ par} & $a^{r/2} \bmod 15$ & $\gcd(a^{r/2}-1, 15)$ & $\gcd(a^{r/2}+1, 15)$ \\
        \midrule
        2  & 4 & Sí & 4  & 3                & 5                \\
        4  & 2 & Sí & 4  & 3                & 5                \\
        7  & 4 & Sí & 4  & 3                & 5                \\
        8  & 4 & Sí & 4  & 3                & 5                \\
        11 & 2 & Sí & 11 & 3                & 5                \\
        13 & 4 & Sí & 4  & 3                & 5                \\
        14 & 2 & Sí & 14 & \textit{trivial} & \textit{trivial} \\
        \bottomrule
    \end{tabular}
\end{table}

La rutina automatiza este análisis: filtra las bases cuyo periodo $r$ es par, verifica la condición $a^{r/2} \not\equiv -1 \pmod{N}$ y descarta las que producen factores triviales (1 o $N$). La condición de \textit{degeneración} descrita se refiere exclusivamente al post-procesamiento clásico de fracciones continuas, no al circuito cuántico propiamente dicho, lo que justifica la inclusión diagnóstica de $a = 14$ en la Fase~3 (Sección~\ref{sec:protocolo_fases}).\\

Desde la perspectiva de la complejidad del circuito cuántico, la selección de $a$ tiene consecuencias directas. Las bases con período $r = 2$ (como $a = 4$) requieren que la estimación de fase distinga entre solo dos fases ($s/r \in \{0/2, 1/2\}$), lo que genera circuitos de menor profundidad efectiva. En contraste, las bases con $r = 4$ (como $a = 7$ o $a = 13$) exigen resolver cuatro fases y producen operadores de exponenciación modular más complejos. El diseño experimental explota esta diferencia evaluando sistemáticamente ambos grupos para cuantificar el impacto de la profundidad del circuito en la degradación por ruido.

\subsection{Dimensionamiento de los registros cuánticos}

El tamaño de los registros cuánticos se calcula dinámicamente a partir de $N$. Esta configuración es la utilizada en todos los experimentos reportados en el Capítulo~\ref{chap:resultados}.

\begin{table}[h]
\centering
\caption{Configuración de registros cuánticos para $N = 15$.}
\label{tab:registros}
% Ajuste del espaciado vertical
\renewcommand{\arraystretch}{1.5} 
\begin{tabular}{ccc}
\hline
\textbf{Registro} & \textbf{Fórmula} & \textbf{Valor} \\
\hline
Trabajo (qubits) & $L = \lceil \log_2 N \rceil = \lceil \log_2 15 \rceil$ & 4 qubits \\
Conteo / control & $t = 2L + 1$ & 9 qubits \\
\hline
\textbf{Total lógico} & $L + t$ & \textbf{13} qubits \\
Registro clásico & mediciones del control & 9 bits \\
\hline
\end{tabular}
\end{table}

% =====================================================================
% 5.4 CONSTRUCCIÓN DEL CIRCUITO IDEAL
% =====================================================================
\section{Construcción del circuito ideal}
\label{sec:construccion_circuitos}

La síntesis del circuito de Shor se organiza como una secuencia de tres bloques funcionales que reproducen el Algoritmo~\ref{alg:qpe} de estimación de fase cuántica: la creación de la superposición inicial, la aplicación del oráculo de exponenciación modular controlada y la extracción de fase mediante la QFT inversa, tal como se muestra en la Figura~\ref{fig:shor_generalizado}.

\begin{figure}[htbp]
    \centering
    \includegraphics[width=\textwidth]{Imagenes/circuito_shor_generalizado.png}
    \caption{Representación esquemática del algoritmo de Shor para la búsqueda de periodo. El registro superior ($n=9$ qubits) se inicializa en una superposición uniforme mediante compuertas Hadamard ($H$). El registro inferior ($m=4$ qubits) se inicializa en el estado computacional $\ket{1}$ mediante una compuerta $X$. La cascada de oráculos representa la exponenciación modular controlada $c$-$U_{a^{2^j} \pmod N}$, donde cada base $b = a^{2^j} \pmod N$ ha sido pre-computada clásicamente. El proceso finaliza con una Transformada de Fourier Cuántica Inversa ($QFT^\dagger$) y la medición del registro de control.}
    \label{fig:shor_generalizado}
\end{figure}

\subsection{Síntesis del oráculo de exponenciación modular}

El núcleo computacional del circuito reside en la secuencia de operaciones controladas $c$-$U_{a}^{2^j}$ para $j = 0, 1, \ldots, t-1$. La construcción de cada operador unitario $U_{a}^{2^j}$ sigue un procedimiento en dos pasos.\\

\begin{enumerate}
    \item \textit{Pre-cómputo clásico.} Para cada qubit de control $j$, se calcula clásicamente la constante $b = a^{2^j} \bmod N$ mediante exponenciación modular por cuadrados repetidos. Esta operación, ejecutada en tiempo clásico $\mathcal{O}(L^2)$, evita la necesidad de implementar cuánticamente la secuencia completa de cuadrados repetidos descrita en la Ec.~\ref{eq:repeated_squaring}, lo cual habría requerido miles de compuertas adicionales.\\
    \item \textit{Construcción de la matriz de permutación.} Con la constante $b$ ya conocida, se construye una matriz de permutación $P_b$ de dimensión $2^L \times 2^L$ que codifica el mapeo $|x\rangle \to |bx \bmod N\rangle$ para todos los estados $x \in \{0, 1, \ldots, N-1\}$. Los estados con $x \geq N$, que no tienen significado aritmético en el contexto del algoritmo, se mapean a sí mismos, es decir, $P_b[x][x] = 1$ para $x \geq N$, garantizando que la matriz sea una permutación unitaria completa. La matriz resultante se encapsula como una compuerta unitaria que posteriormente se convierte en una compuerta controlada.\\ 
\end{enumerate}

Este enfoque de síntesis ---también denominado \textit{compiled Shor's algorithm} en la literatura \cite{Skosana2021,Lanyon2007}--- tiene una ventaja decisiva para la implementación NISQ: el transpilador recibe una única compuerta unitaria controlada de $L+1$ qubits en lugar de un circuito aritmético profundo. Aunque el transpilador deberá descomponer esta unitaria en compuertas nativas del hardware, el conteo final de compuertas de dos qubits es significativamente menor que el de una implementación aritmética genérica \cite{smolin2013oversimplifying}.

% =====================================================================
% 5.2 PROTOCOLO EXPERIMENTAL EN TRES FASES
% =====================================================================
\section{Protocolo experimental en tres fases}
\label{sec:protocolo_fases}

El protocolo experimental se diseñó como una secuencia de tres entornos de ejecución progresivamente más realistas, de modo que la diferencia entre los resultados de dos fases consecutivas pueda atribuirse a una única fuente de degradación añadida.\\

\textbf{Fase 1: simulación ideal}\\

El objetivo de esta fase es doble: verificar la corrección lógica del circuito reproduciendo las distribuciones predichas por la teoría de QPE y establecer la línea base cuantitativa contra la cual se medirá la degradación en las fases posteriores. Como variante complementaria se ejecuta el mismo circuito sobre el simulador inicializado a partir del \textit{snapshot} del procesador real con el modelo de ruido explícitamente desactivado: esta variante fuerza la topología Heavy-Hex sin aplicar decoherencia y aísla así el \textit{overhead} topológico del impacto del ruido.\\

\textbf{Fase 2: simulación con ruido (\textit{Fake Backends})}\\

El \textit{backend} se instancia a partir del \textit{snapshot} de calibración del procesador \texttt{ibm\_torino}, lo que carga automáticamente: los errores incoherentes (despolarización y relajación térmica) parametrizados por las tasas de error y los tiempos $T_1^{(i)}$, $T_2^{(i)}$ individuales de cada qubit; las matrices de confusión asimétricas de los errores de lectura; y el mapa de acoplamiento Heavy-Hex con las duraciones nominales de cada compuerta nativa. La limitación física fundamental de esta fase, que condiciona la lectura de sus resultados, es que el \textit{snapshot} captura exclusivamente ruido incoherente: no modela los errores coherentes (sobre-rotaciones sistemáticas) ni el \textit{crosstalk} dinámico entre qubits adyacentes que sí están presentes en el procesador físico. En consecuencia, los resultados de la Fase~2 deben interpretarse como una \textit{cota inferior} de la degradación esperable en hardware, y no como un predictor exacto de la Fase~3.\\

\textbf{Fase 3: ejecución en hardware real}\\

Los circuitos se envían al procesador \texttt{ibm\_torino} (Heron~r1, 133~qubits superconductores transmon en topología Heavy-Hex) a través del servicio en la nube de IBM Quantum bajo el plan de acceso abierto. La selección de este \textit{backend} obedece a dos criterios. Primero, por consistencia metodológica: el \textit{snapshot} empleado en la Fase~2 es directamente el de \texttt{ibm\_torino}, de modo que cualquier diferencia entre las Fases~2 y~3 puede atribuirse al ruido coherente y al \textit{crosstalk} no modelados, y no a discrepancias de arquitectura, topología o conjunto de compuertas nativas. Segundo, por disponibilidad: durante el período experimental, los demás procesadores Heron presentaban colas con decenas de \textit{jobs} en espera, lo que habría comprometido la agilidad necesaria para iterar sobre múltiples configuraciones. Las características nominales del \textit{backend} se reportan en el Capítulo~\ref{chap:resultados} y el desglose técnico ---identificadores de \textit{job}, configuración exacta de \texttt{SamplerV2} y \texttt{SamplerOptions}, organización de los estudios A2, A3 y A6/A7--- se documenta en el Apéndice~\ref{app:scripts-fases}.\\

En las tres fases, los resultados de cada ejecución se almacenan en formato JSON (incluyendo los identificadores de \textit{job} en el caso del hardware) para garantizar la trazabilidad y permitir el post-procesamiento asíncrono. El flujo de trabajo experimental se esquematiza en la Figura~\ref{fig:protocolo_fases}.

\newpage

\begin{figure}[H]
    \centering
    \includegraphics[width=\textwidth]{Imagenes/diagrama_protocolo.png}
    \caption{Diagrama de flujo del protocolo experimental por fases.}
    \label{fig:protocolo_fases}
\end{figure}

Las diferencias entre fases obedecen a un compromiso entre rigor estadístico y costo computacional: en Fase~1, el bajo costo del simulador ideal permite barrer las cuatro bases diagnósticas y los cuatro niveles de optimización con $4096$~\textit{shots}; en Fase~2, el costo del modelo de ruido (aproximadamente dos órdenes de magnitud superior) obliga a restringir el barrido a las dos bases representativas $a \in \{4, 7\}$ y a los niveles extremos $\texttt{opt} \in \{0, 3\}$; en Fase~3, el cupo de \textit{jobs} disponibles bajo el plan de acceso abierto fija el alcance experimental.\\

\textit{Bases diagnósticas:} la base $a = 14$ produce factores triviales tras el post-procesamiento clásico (cf. Sección~\ref{sec:preprocesamiento_parametros}), pero su circuito es instancia plenamente válida de la estimación de fase cuántica con orden $r = 2$. Su inclusión en la Fase~3 no busca factorizar $N$, sino caracterizar la respuesta del procesador para circuitos cortos: en estas ejecuciones las métricas relevantes son exclusivamente la fidelidad de Hellinger $\mathcal{F}_H$ y la $PST$, que cuantifican la calidad de la señal cuántica con independencia del éxito del filtro clásico.

% =====================================================================
% 5.5 TRANSPILACIÓN AL HARDWARE
% =====================================================================
\section{Transpilación}
\label{sec:transpilacion}

Las decisiones metodológicas relevantes en esta etapa son las siguientes. Primero, se barrió el nivel de optimización del transpilador en su rango completo (\texttt{opt}~$\in \{0, 1, 2, 3\}$) para cuantificar el impacto progresivo de las transformaciones sobre el circuito final: el nivel 0 se limita a la traducción mínima al conjunto nativo y sirve como referencia; el nivel 3 activa consolidación de bloques unitarios, cancelación conmutativa y re-síntesis con conciencia del modelo de ruido del \textit{backend}. Segundo, tanto la asignación inicial de qubits como el enrutamiento posterior se delegaron al algoritmo SABRE descrito en la Sección~\ref{sec:transpilacionysabre} del marco teórico. Tercero, la semilla del transpilador se fijó (\path{seed_transpiler}$=457$) para garantizar la reproducibilidad estadística de las heurísticas. Cuarto, durante la síntesis se permitió una aproximación controlada al operador exacto fijando \path{approximation_degree}$=0{,}7$, con la justificación expuesta en Apéndice~\ref{app:sabre-avanzado}. Para las ejecuciones en hardware se añadió, además, una planificación temporal \textit{as late as possible}, que concentra los intervalos \textit{idle} al inicio del circuito y deja espacio efectivo para la inserción de las secuencias de Desacoplamiento Dinámico descritas en la Sección~\ref{sec:supresion_implementada}.\\

Como métrica primaria de la calidad del circuito transpilado se adoptó la profundidad de dos qubits, $D_{2Q}$, definida como la longitud del camino crítico considerando exclusivamente las compuertas que operan sobre dos qubits. Esta es la métrica predictivamente más significativa para hardware NISQ: cada compuerta CZ contribuye un error del orden de $\epsilon_{2Q} \sim 5 \times 10^{-3}$ en el mejor de los casos en Heron~r1, de modo que $D_{2Q}$ controla, en primer orden, la fidelidad acumulada del circuito completo. La tabulación de $D_{2Q}$ por nivel de optimización permite, además, separar el efecto del \textit{layout} del efecto del \textit{routing} cuando se compara una asignación trivial frente a SABRE.\\

La lista detallada de parámetros del transpilador, los valores nominales de los hiperparámetros avanzados de SABRE (número de iteraciones, ensayos de \textit{layout} y de SWAP) y el procedimiento concreto por el cual se inyectan en el \textit{pipeline} de transpilación se documentan en el Apéndice~\ref{app:sabre-avanzado}.



% =====================================================================
% 5.7 TÉCNICAS DE SUPRESIÓN DE ERRORES
% =====================================================================
\section{Técnicas de supresión de errores}
\label{sec:supresion_implementada}

El estudio del impacto de las técnicas de supresión de errores se organizó como un diseño factorial $2 \times 2$ que combina la activación o desactivación de Desacoplamiento Dinámico (DD) y de Pauli Twirling (PT). Las cuatro variantes resultantes (sin supresión, solo DD, solo PT, DD$+$PT) se evaluaron sobre la misma transpilación del circuito, de modo que las diferencias observadas en las métricas se atribuyen exclusivamente al efecto de las técnicas.\\

\textbf{Selección de la secuencia DD}\\

El servicio de ejecución admite varias secuencias de pulsos para DD, entre ellas \texttt{XX}, \texttt{XpXm} y \texttt{XY4}. En este trabajo se adoptó \texttt{XY4} en sustitución de \texttt{XpXm}, que era la opción heredada del repositorio base. La secuencia \texttt{XY4} cancela a primer orden tanto el ruido de fase como el de bit a través de cuatro pulsos en direcciones ortogonales del plano $XY$, mientras que \texttt{XpXm} se limita a refocalizar el desfase de un solo eje. Cuando DD está habilitado, el servicio identifica automáticamente los intervalos \textit{idle} del circuito transpilado, calcula cuántos pulsos completos caben en cada intervalo según la duración del pulso $X$ nativo del dispositivo e inserta la secuencia centrada simétricamente. Esto es especialmente relevante en el circuito de Shor, donde la aplicación secuencial de los $c\text{-}U_a^{2^j}$ genera largos periodos de inactividad en los qubits de control que no están interactuando en un paso dado.\\

\textbf{Configuración del Pauli Twirling}\\

En lugar del modo pasivo (en el que cada compuerta de dos qubits se aleatoriza con un único par de Pauli), se seleccionó la estrategia \texttt{active}, que aplica el \textit{twirling} únicamente a los qubits efectivamente involucrados en cada compuerta y delega al \textit{runtime} la generación automática de un número de instancias aleatorias coherente con el total de \textit{shots}. Esta configuración es la recomendada por la primitiva V2 cuando el objetivo es convertir errores coherentes sistemáticos en ruido estocástico simétrico (Ec.~\ref{eq:twirling_exacto}) y simultáneamente preservar la economía de \textit{shots}.\\

La parametrización exacta de los objetos de configuración de DD y PT (nombres de los \textit{flags}, valores escogidos para cada uno y el procedimiento de inyección en \texttt{SamplerOptions}) se documenta en el Apéndice~\ref{app:supresion-config}.

% =====================================================================
% 5.8 POST-PROCESAMIENTO Y MÉTRICAS
% =====================================================================
\section{Post-procesamiento clásico y métricas de evaluación}
\label{sec:postprocesamiento}

Tras el muestreo, los \textit{bitstrings} obtenidos se procesan clásicamente para extraer el orden $r$ y, a partir de él, los factores no triviales de $N$. El procedimiento, común a las tres fases, sigue seis pasos:

\begin{enumerate}
    \item Recuperación de la distribución de conteos del registro clásico de salida.
    \item Selección de los $k$ resultados con mayor frecuencia (típicamente $k = 10$), ordenados de forma descendente.
    \item Para cada \textit{bitstring} seleccionado, conversión a entero $y$ aplicando la convención \textit{little-endian} (Apéndice~\ref{app:migracion-qiskit}).
    \item Cálculo de la fase estimada $\tilde{\varphi} = y / 2^t$.
    \item Aplicación del algoritmo de fracciones continuas con denominador acotado por $N$ para obtener los convergentes $p_k / q_k$ con $q_k < N$. El denominador $q_k$ es un candidato al orden $r$.
    \item Verificación clásica: $a^{q_k} \equiv 1 \pmod{N}$, $q_k$ par y factores no triviales $\gcd(a^{q_k/2} \pm 1, N) \notin \{1, N\}$.
\end{enumerate}

Cuando el filtro clásico devuelve un par de factores no triviales, la ejecución se reporta como exitosa. Los detalles de implementación se documentan en el Apéndice~\ref{app:postproc-modulos}.\\

\textbf{Métricas adoptadas}\\

La evaluación cuantitativa de cada ejecución se realiza con dos métricas complementarias, definidas formalmente en el marco teórico: la fidelidad de Hellinger $\mathcal{F}_H$ y la probabilidad de éxito experimental $PST$. Ambas se calculan comparando la distribución experimental contra la distribución ideal obtenida del simulador sin ruido.\\

Para interpretar los resultados que se presentan en el Capítulo~\ref{chap:resultados} se adoptan umbrales operacionales explícitos sobre ambas métricas. En primer lugar, siguiendo el criterio teórico de \cite{Bagourd2026} discutido en el marco teórico, una $PST$ se considera \textit{aceptable} únicamente si supera de manera estadísticamente significativa la línea base uniforme $1/r$, que vale $0{,}5$ para $r = 2$ y $0{,}25$ para $r = 4$. El caso ideal $PST = 1{,}0$ corresponde a toda la masa de probabilidad concentrada sobre los picos correctos del QPE; en consecuencia, el umbral mínimo práctico se fija en $PST > 1/r + \delta$ con $\delta = 0{,}1$, margen que se establece como cota de significancia estadística compatible con el número de \textit{shots} empleado en cada fase. En cuanto a la fidelidad de Hellinger, se adoptan los umbrales habituales en la literatura de \textit{benchmarking} de circuitos NISQ: $\mathcal{F}_H > 0{,}9$ indica \textit{alta fidelidad}, el rango $0{,}7 \leq \mathcal{F}_H \leq 0{,}9$ señala \textit{fidelidad moderada} y $\mathcal{F}_H < 0{,}7$ se interpreta como \textit{degradación severa} de la distribución respecto a la ideal. La interpretación combinada de ambos umbrales constituye el criterio central de análisis en el Capítulo~\ref{chap:resultados}.\\

Todo el código fuente, los archivos de configuración empleados en cada experimento y los identificadores de los \textit{jobs} de IBM Quantum se encuentran documentados en el repositorio público del proyecto \cite{compuertas_CNOT}, garantizando la reproducibilidad completa del estudio.

\chapter{Resultados}
\label{chap:resultados}

% ============================================================
% SECCIÓN 6.1 — RESULTADOS DE LA SIMULACIÓN IDEAL
% Versión compacta sin tablas redundantes (Tablas 6.1–6.5 omitidas;
% Tabla de fracciones continuas trasladada al Apéndice C).
% Referencias exclusivas a Figuras 6.1–6.6.
% ============================================================
\section{Resultados de la simulación ideal}
\label{sec:resultados_ideal}

La validación del algoritmo de Shor para $N = 15$ se aborda en tres fases progresivas (simulación ideal, simulación con ruido calibrado de \texttt{FakeTorino} y ejecución en \texttt{ibm\_torino}). Esta sección reporta la primera fase, ejecutada en el simulador \texttt{AerSimulator} de Qiskit.

% ============================================================
% 6.1.1  CARACTERIZACIÓN DE BLOQUES FUNDAMENTALES
% ============================================================
\subsection{Caracterización de los bloques fundamentales}
\label{sec:caracterizacion_bloques}

Antes de evaluar el circuito completo, se verificó individualmente la corrección de cada sub-circuito constituyente: la Transformada Cuántica de Fourier ($\mathrm{QFT}$) y la exponenciación modular controlada ($c$-$U_{a^{2^k}}$). La verificación numérica de la unitariedad arrojó $\lVert U^{\dagger} U - I \rVert_F < 10^{-13}$ para todas las variantes evaluadas (QFT directa e inversa, descompuesta y transpilada), valor compatible con la precisión de punto flotante. Para el operador de exponenciación modular se verificó adicionalmente la propiedad de periodicidad $(U_a)^r = I$, obteniendo $\lVert (U_a)^r - I \rVert_F = 0$ en todas las bases. El detalle exhaustivo se documenta en el Apéndice~\ref{app:caracterizacion-bloques}.\\

Un resultado estructural relevante para todo lo que sigue es la asimetría entre bases (Fig. ~\ref{fig:circuitos_a4_a7}). Para cada $a$, solo los bloques con $a^{2^k} \bmod N \neq 1$ contribuyen a la profundidad del circuito; los demás colapsan a la identidad. Esto explica por qué las bases con $r = 2$ producen circuitos significativamente menos profundos que $a = 7$ (con $r = 4$): para $r = 2$ únicamente el bloque $k = 0$ es no trivial, mientras que para $a = 7$ los bloques $k = 0$ y $k = 1$ lo son simultáneamente ($7^1 \bmod 15 = 7$ y $7^2 \bmod 15 = 4$). Esta diferencia tiene una consecuencia directa para la viabilidad experimental en hardware NISQ: la base $a = 7$, al exigir el doble de compuertas entrelazantes, será la más susceptible a la decoherencia en las fases experimentales con ruido. La selección clásica de $a$ como estrategia de pre-procesamiento (Sección~\ref{sec:preprocesamiento_parametros}) queda así cuantitativamente justificada.

\begin{figure}[H]
    \centering
    \includegraphics[width=0.8\textwidth]{Imagenes/circuitos_a4_a7.png}
    \caption{Diagrama esquemático del circuito \texttt{RegisterQC} para las bases $a=4$ (superior) y $a=7$ (inferior) con $N=15$. La duplicación visible de bloques entrelazantes en $a=7$ refleja directamente que su orden $r = 4$ activa dos potencias no triviales de $U_a$.}
    \label{fig:circuitos_a4_a7}
\end{figure}

Tras transpilar a la base lógica $\{\mathrm{CX}, \mathrm{U3}\}$ sin restricción topológica, el costo entrelazante refleja con nitidez la asimetría anterior: las bases con $r = 2$ requieren del orden de $230$ compuertas de dos qubits y $\sim 460$ capas de profundidad total, mientras que $a = 7$ demanda alrededor de $463$ compuertas y $\sim 919$ capas, prácticamente el doble. La $\mathrm{QFT}$ aporta apenas $\sim 30$ capas a la profundidad de dos qubits, en concordancia con su complejidad asintótica $\mathcal{O}(t^2)$ frente a la $\mathcal{O}(L^3)$ de la exponenciación modular, que domina abrumadoramente el costo total (Apéndice \ref{app:caracterizacion-bloques}).

% ============================================================
% 6.1.2  SIMULACIÓN IDEAL PURA (ALL-TO-ALL)
% ============================================================
\subsection{Simulación ideal pura: línea base sin restricción topológica}
\label{sec:resultados_ideal_puro}

La primera ejecución del circuito completo se realizó con conectividad \textit{all-to-all}, es decir, sin restricción topológica. El propósito es doble: establecer la profundidad intrínseca del circuito y confirmar que, en ausencia de ruido, la probabilidad de señal alcanza el valor teórico del 100\%.\\

\begin{figure}[H]
    \centering
    \includegraphics[width=0.95\textwidth]{simulacion_ideal/histogramas_picos_qpe_N15.png}
    \caption{Histogramas de medición del registro de conteo en simulación ideal del algoritmo de Shor ($N = 15$, $t = 9$, 4096 \textit{shots}) para las cuatro bases estudiadas. Las líneas verticales punteadas marcan los picos teóricos $y_s$ predichos por el QPE; las barras coloreadas muestran los conteos observados.}
    \label{fig:histogramas_picos}
\end{figure}

El resultado central es contundente: en las 16 configuraciones evaluadas, la totalidad de las $M = 4096$ mediciones cae sobre los picos teóricos del QPE, esto es, $PST = 100\%$. Los histogramas de la Figura~\ref{fig:histogramas_picos} muestran que la masa de probabilidad se concentra con probabilidad $\approx 1/2$ en cada uno de los dos picos $\{0, 256\}$ para las bases con $r = 2$, y con probabilidad $\approx 1/4$ en cada uno de los cuatro picos $\{0, 128, 256, 384\}$ para $a = 7$. Esta concordancia cuantitativa con la predicción del QPE valida que los operadores unitarios $U_a$ y la $\mathrm{QFT}^{\dagger}$ fueron implementados sin errores lógicos. Cualquier desviación observable respecto a este patrón ideal podrá imputarse, en las fases posteriores, al ruido de compuerta y la decoherencia (\texttt{FakeTorino}) o a los efectos adicionales del hardware real (\textit{crosstalk}, errores coherentes, deriva de calibraciones).\\

Por otro lado, la profundidad del circuito por base y nivel de optimización se muestra prácticamente invariante (Figura~\ref{fig:ideal_puro}): las diferencias entre Opt=0 y Opt=3 son de a lo sumo unas pocas unidades. Esta invariancia es físicamente esperada en conectividad \textit{all-to-all}: el transpilador no necesita insertar compuertas SWAP, y la cadena entrelazante proviene directamente de la descomposición de los $U_a$ controlados, ya cercana a su mínimo estructural. Lo que sí disminuye monótonamente con Opt es el conteo total de compuertas (Apéndice \ref{app:caracterizacion-bloques}), debido a que se ejecutan pases de optimización a partir de Opt=1. Aislar esta contribución es metodológicamente relevante, pues permite atribuir cualquier variación posterior de $D_{2Q}$ al proceso de ruteo.

\begin{figure}[H]
    \centering
    \includegraphics[width=0.9\textwidth]{simulacion_ideal/analisis_ideal_puro_N15.png}
    \caption{Resultados de la simulación ideal pura (\textit{all-to-all}). Izquierda: histogramas de medición del porcentaje de señal para cada base y nivel de Opt. Derecha: profundidad de dos qubits $D_{2Q}$ por base y nivel de optimización.}
    \label{fig:ideal_puro}
\end{figure}

\subsubsection*{El caso $a = 14$: factores triviales}

Para $a = 14 \equiv -1 \pmod{15}$ el circuito cuántico funciona correctamente (el QPE recupera $r = 2$ con señal del 100\%), pero la fase clásica de extracción produce factores triviales:
\begin{equation}
    \gcd(14^{r/2} - 1, 15) = \gcd(13, 15) = 1, \qquad \gcd(14^{r/2} + 1, 15) = \gcd(15, 15) = 15.
\end{equation}
Esta situación es una propiedad algebraica intrínseca del algoritmo de Shor, no un error computacional: cuando $a \equiv -1 \pmod{N}$, $a^{r/2} \equiv -1 \pmod{N}$ viola la condición necesaria $a^{r/2} \not\equiv -1 \pmod{N}$ del Algoritmo~\ref{alg:shor}. En tal caso, el protocolo prescribe reiniciar con un nuevo $a$~\cite{Nielsen2024}.

% ============================================================
% 6.1.3  SIMULACIÓN IDEAL EN FAKETORINO
% ============================================================
\subsection{Transpilación a la topología de FakeTorino}
\label{sec:resultados_faketorino}

La segunda fase del análisis transpila \texttt{RegisterQC} a la topología Heavy-Hex de \texttt{FakeTorino} y la ejecuta en modo ideal, sin modelo de ruido. Este experimento es metodológicamente crucial: aísla y cuantifica el \textit{overhead} de profundidad que introduce el ruteo (la inserción de compuertas SWAP exigidas por la conectividad del chip), sin que ese efecto se confunda con la degradación que el ruido producirá más adelante.\\

La Figura~\ref{fig:ideal_faketorino} hace visible este \textit{overhead}. A pesar del incremento significativo en la profundidad por la inserción de SWAPs, la señal $PST$ se mantiene en el 100\% para todas las bases, confirmando que el \textit{overhead} topológico no introduce errores lógicos en ausencia de ruido. La discrepancia entre la profundidad ideal pura y la transpilada cuantifica el costo exclusivo del mapeo geométrico. Por otro lado la profundidad de compuertas de dos qubits, si disminuye con el nivel de optimización, lo cual se puede atribuir a los pases de \textit{Opt} que busca minimizar el número de compuertas SWAP.

\begin{figure}[H]
    \centering
    \includegraphics[width=\textwidth]{simulacion_ideal/analisis_ideal_faketorino_N15.png}
    \caption{Resultados de la simulación ideal del algoritmo de Shor para $N = 15$ transpilado a la topología Heavy-Hex de \texttt{FakeTorino}.}
    \label{fig:ideal_faketorino}
\end{figure}

% ============================================================
% 6.1.4  COMPARACIÓN TRIVIAL vs SABRE
% ============================================================
\subsection{Impacto de la estrategia de mapeo: trivial vs.\ SABRE}
\label{sec:resultados_comparacion_transpilacion}

\begin{figure}[H]
    \centering
    \includegraphics[width=\textwidth]{simulacion_ideal/comparar_transpilacion_N15.png}
    \caption{Comparación de las métricas de transpilación entre el mapeo trivial y SABRE para el algoritmo de Shor ($N = 15$) sobre la topología Heavy-Hex de \texttt{FakeTorino}.}
    \label{fig:comparar_transpilacion}
\end{figure}

Para evaluar cuantitativamente la efectividad del algoritmo SABRE (Ec.~\ref{eq:sabre_completa}) frente al mapeo trivial ($\pi(q_i) = \text{qubit}_i$), se transpiló el circuito bajo ambas estrategias de \textit{layout} para las cuatro bases y los cuatro niveles de optimización, manteniendo el ruteo SABRE en ambos casos. La probabilidad de señal se conservó en $PST = 100\%$ y la fidelidad de Hellinger en $\mathcal{F}_H \geq 0.9996$ para las 32 configuraciones (esta cota se adopta como umbral de referencia para las fases siguientes, cualquier configuración que produzca en \texttt{FakeTorino} o en \texttt{ibm\_torino} una fidelidad sensiblemente inferior deberá atribuirse a ruido físico y no al tamaño finito de la muestra), lo que confirma un punto importante: en simulación ideal, la estrategia de asignación no afecta la corrección lógica del circuito, sino exclusivamente el costo físico de su ejecución.\\

La Figura~\ref{fig:comparar_transpilacion} muestra que en todas las configuraciones SABRE mejora al mapeo trivial, pero la magnitud del beneficio depende de manera \emph{no monótona} del nivel de optimización. Esta es, quizá, la observación más fina de la sección. La leve regresión observada en algunos casos al pasar de Opt=2 a Opt=3 (por ejemplo, $a=11$, donde la reducción de $D_{2Q}$ cae del 17.7\% al 9.3\%) es consistente con la naturaleza estocástica de los pases adicionales activados en este nivel: Opt=3 no garantiza una solución superior a Opt=2; ambas deben entenderse como muestreos de una misma distribución de mapeos válidos.\\


% ============================================================
% 6.1.6  POST-PROCESAMIENTO CLÁSICO POR FRACCIONES CONTINUAS
% ============================================================
\subsection{Post-procesamiento clásico: extracción por fracciones continuas}
\label{sec:resultados_fracciones}

Establecida en la Subsección~\ref{sec:resultados_ideal_puro}, resta verificar que la cadena clásica descrita en la Sección~\ref{sec:construccion_circuitos} extrae efectivamente el período $r$ y los factores no triviales de $N$ para cada pico informativo ($y \neq 0$).

\begin{table}[H]
    \centering
    \caption{Cadena de post-procesamiento clásico para los 
    resultados con $y \neq 0$ obtenidos en simulación ideal pura 
    ($N = 15$, $t = 9$, $M = 4096$ \textit{shots}).\\}
    \label{tab:fracciones_continuas}
    \small
    \begin{tabular}{cccccccc}
        \hline
        \textbf{$a$} & \textbf{$r$} & \textbf{$y$} & 
        $\boldsymbol{\tilde{\varphi}}$ & \textbf{FC} & 
        $\boldsymbol{p_k/q_k}$ & $\boldsymbol{a^{q_k} \bmod N}$ & 
        \textbf{Resultado clásico} \\
        \hline
        4  & 2 & 256 & 0.500 & $[0; 2]$    & $1/2$ & 1           & $\gcd = 3,\, 5$ \quad $\checkmark$  \\
        7  & 4 & 128 & 0.250 & $[0; 4]$    & $1/4$ & 1           & $\gcd = 3,\, 5$ \quad $\checkmark$  \\
        7  & 4 & 384 & 0.750 & $[0; 1, 3]$ & $3/4$ & 1           & $\gcd = 3,\, 5$ \quad $\checkmark$  \\
        7  & 4 & 256 & 0.500 & $[0; 2]$    & $1/2$ & 4 ($\neq 1$) & Convergente no válido \quad $\dagger$ \\
        11 & 2 & 256 & 0.500 & $[0; 2]$    & $1/2$ & 1           & $\gcd = 5,\, 3$ \quad $\checkmark$  \\
        14 & 2 & 256 & 0.500 & $[0; 2]$    & $1/2$ & 1           & Triviales: $\gcd = 1,\, 15$ \quad $\triangle$ \\
        \hline
    \end{tabular}
\end{table}

El protocolo recupera el orden correcto en todos los casos no degenerados: $r = 2$ para $a \in \{4, 11\}$ y $r = 4$ para $a = 7$, conduciendo en todos ellos a la factorización $15 = 3 \times 5$. Para $a = 14$ se recupera correctamente $r = 2$, pero la fase clásica produce factores triviales, ya discutido en la Subsección~\ref{sec:resultados_ideal_puro}.\\

Existe, sin embargo, un único modo de falla intermedio que merece comentario por su valor pedagógico. Para $a = 7$ y la medición $y = 256$, la fase $\tilde{\varphi} = 1/2$ produce la expansión $[0; 2]$ y un único convergente no trivial $q_1 = 2$. Este denominador respeta la cota $1 < q_1 < N$ pero no sobrevive la verificación modular: $7^2 \bmod 15 = 4 \neq 1$. La explicación es directa: el valor $y = 256$ corresponde al índice $s = 2$ con $r = 4$, y la fracción $s/r = 2/4$ se reduce a $1/2$ en términos mínimos, por lo que el algoritmo de Euclides devuelve un \emph{divisor} de $r$, no $r$ mismo. Este caso justifica empíricamente por qué la verificación modular es un segundo criterio indispensable, y no una redundancia: aproximadamente el 24\% de las mediciones simuladas para $a = 7$ produce un $q_k$ que respeta la cota de denominador pero no corresponde al orden buscado. En el experimento, este recurso no resultó necesario, pues las mediciones $y \in \{128, 384\}$ recuperan directamente $r = 4$ con probabilidad combinada $\approx 0.508$.\\

Un último elemento diferencia estructuralmente a $a = 7$ del resto. La recuperación de su período exige convergentes de profundidad $k = 2$ (la expansión $[0; 1, 3]$ que produce $3/4$), mientras que las bases con $r = 2$ se resuelven en $k = 1$. Como la precisión con que $\tilde{\varphi}$ aproxima $s/r$ se degrada bajo ruido y los convergentes de mayor profundidad exigen mayor precisión de fase para satisfacer la cota del Teorema de Convergencia, $a = 7$ acumula dos fuentes simultáneas de vulnerabilidad: mayor profundidad del circuito cuántico y mayor profundidad del árbol clásico de convergentes. Esta doble sensibilidad anticipa que en las Secciones~\ref{sec:resultados_ruido} y~\ref{sec:resultados_hardware} $a = 7$ exhibirá la degradación más pronunciada del conjunto.

% ============================================================
% SECCIÓN 6.2 — RESULTADOS DE LA SIMULACIÓN CON RUIDO
% ============================================================
\section{Resultados de la simulación con ruido: FakeTorino}
\label{sec:resultados_ruido}

Tras validar el circuito en condiciones ideales (Sección~\ref{sec:resultados_ideal}), se evalúa ahora cómo se comporta el algoritmo cuando se le somete a un modelo de ruido realista. Esta fase intermedia entre la simulación pura y el hardware constituye el primer encuentro de \texttt{RegisterQC} con la fragilidad del régimen NISQ: la pregunta ya no es si el circuito es lógicamente correcto, sino si la información cuántica sobrevive lo suficiente como para que el post-procesamiento clásico siga siendo capaz de extraer los factores de $N = 15$.

% ============================================================
% 6.2.2  CHOQUE CON EL RUIDO: IDEAL FRENTE A REALISTA
% ============================================================
\subsection{Choque con el ruido: comparación frente a la simulación ideal}
\label{subsec:comparativo_ideal_ruido}

Para esta fase se concentró el barrido en las dos bases más informativas, $a = 4$ ($r = 2$) y $a = 7$ ($r = 4$), evaluadas en los niveles extremos del transpilador (\texttt{opt}=0 y \texttt{opt}=3). El contraste con la simulación ideal de la Sección~\ref{sec:resultados_ideal} es inmediato. Allí los histogramas mostraban toda la masa de probabilidad alojada en los picos teóricos, con líneas verticales nítidas y el resto del eje rigurosamente vacío. Aquí (Figura~\ref{fig:histogramas_picos_ruido}), en cambio, los picos teóricos aparecen visiblemente \emph{achicados} y, en su lugar, surge un \emph{piso de ruido} continuo: una multitud de estados que en condiciones ideales eran imposibles ahora aparecen con probabilidad pequeña pero no nula, repartiendo masa por todo el registro de conteo. Esta migración de probabilidad desde los picos hacia el resto del espectro es la firma visual de la decoherencia y de los errores de compuerta acumulados a lo largo del circuito profundo, y se aprecia con claridad en el panel central de la Figura~\ref{fig:analisis_ruido_faketorino}, donde se junta el $PST$ ideal con el ruidoso.

\begin{figure}[H]
    \centering
    \includegraphics[width=\textwidth]{simulacion_ruido/histograma_picos_ruido_N15.png}
    \caption{Histogramas de la estimación de fase del algoritmo de Shor ($N=15$, $t=9$, $M=512$ \textit{shots}) con \texttt{FakeTorino}, para las cuatro configuraciones evaluadas. Las líneas verticales discontinuas marcan los picos teóricos $y_s = s\,2^m/r$ y la línea horizontal punteada indica la altura teórica ideal $1/r$. Los recuadros internos reproducen cada distribución en escala logarítmica para hacer visible el \emph{piso de ruido} y las contribuciones de baja probabilidad que la escala lineal oculta.}
    \label{fig:histogramas_picos_ruido}
\end{figure}

\begin{figure}[H]
    \centering
    \includegraphics[width=\textwidth]{simulacion_ruido/analisis_ruido_faketorino_N15.png}
    \caption{Simulación con ruido del algoritmo de Shor para $N=15$ en \texttt{FakeTorino} (sin supresión). Panel izquierdo: evolución del $PST$ por nivel de optimización. Panel central: $PST$ ideal (tenue) frente al ruidoso (sólido). Panel derecho: fidelidad de Hellinger $\mathcal{F}_H$ resultante.}
    \label{fig:analisis_ruido_faketorino}
\end{figure}

La caída cuantitativa es severa pero asimétrica entre bases, y se manifiesta de manera coherente en ambas métricas. Para $a = 4$ con \texttt{opt}=3, el $PST$ cae del 100\% al $4.3\%$ y $\mathcal{F}_H$ se desploma de $\geq 0.9996$ a $0.0010$: prácticamente toda la señal cuántica se ha disuelto en el piso de ruido y la distribución empírica es estadísticamente indistinguible del fondo uniforme sobre los $2^9 = 512$ resultados posibles. Para $a = 7$, en cambio, el $PST$ se mantiene cerca del $59\%$ con \texttt{opt}=3 y la fidelidad asciende a $\mathcal{F}_H \approx 0.41$, valores sorprendentemente altos si se recuerda que su circuito tiene casi el doble de profundidad entrelazante.\\

Esta inversión (el circuito más profundo conserva mejor la señal) parece, a primera vista, paradójica, pero tiene una explicación enteramente geométrica. Para $a = 4$, la distribución teórica está concentrada en solo dos picos altos ($p = 0.5$ cada uno), de modo que cualquier fuga de probabilidad hacia el piso de ruido afecta proporcionalmente más a la señal y, en consecuencia, también a la fidelidad de la distribución completa. Para $a = 7$, en cambio, la distribución teórica ya está repartida entre cuatro picos más bajos ($p = 0.25$ cada uno), lo que la hace estructuralmente más resiliente al efecto uniformizante del ruido: la métrica $\mathcal{F}_H$ es sensible a la forma de la distribución de referencia, y una distribución teórica más uniforme es inherentemente más tolerante al ruido despolarizante, que tiende precisamente a uniformizar la distribución experimental. Esta asimetría es importante metodológicamente porque demuestra que la resiliencia al ruido no depende solo de la profundidad del circuito (que jugaría en contra de $a = 7$), sino también de la geometría de la distribución que se quiere medir (que esta vez juega a favor).\\

A pesar de esta degradación, el resultado más importante del experimento es positivo: el algoritmo clásico de fracciones continuas extrae correctamente los factores $\{3, 5\}$ en todas las configuraciones evaluadas, incluso en el caso extremo de $a = 4$ con $PST \approx 4\%$ y $\mathcal{F}_H \approx 0.001$. Lo que se observa, en términos físicos, es que basta con que los picos teóricos asomen por encima del piso de ruido, aunque sea apenas, para que el filtro clásico sea capaz de seleccionarlos: la fidelidad de la distribución global puede haberse desplomado, pero la información mínima necesaria para identificar el período sigue presente en su soporte. La Sección~\ref{subsec:fracciones_ruido} analiza con más detalle este mecanismo de rescate.

% ============================================================
% 6.2.3  EFECTO DE LAS TÉCNICAS DE MITIGACIÓN
% ============================================================
\subsection{Efecto de las técnicas de supresión de errores}
\label{subsec:efecto_mitigacion_ruido}

Para evaluar las estrategias de supresión formuladas en la Sección~\ref{sec:supresion_implementada} se ejecutó el circuito en cuatro configuraciones: \textit{baseline} (sin supresión), Desacople Dinámico (DD), Pauli Twirling (PT) y la combinación DD + PT. La Figura~\ref{fig:comparacion_mitigacion_ruido} sintetiza el experimento para ambas bases.\\

El resultado es, a primera vista, contraintuitivo: las técnicas de supresión no mejoran ---e incluso degradan ligeramente--- las métricas de señal en el simulador. Para $a = 7$, la línea base sin supresión alcanza $PST \approx 56.6\%$ y $\mathcal{F}_H \approx 0.40$, mientras que la combinación DD + PT reduce ambos indicadores a $PST \approx 46.5\%$ y $\mathcal{F}_H \approx 0.29$. La caída concertada de las dos métricas en la misma dirección descarta que se trate de una anomalía: reflejan que, en el simulador estocástico, las técnicas de supresión añaden sobrecosto sin recuperar señal de manera sustancial.

\begin{figure}[H]
    \centering
    \includegraphics[width=1.0\textwidth]{simulacion_ruido/comparar_mitigacion_N15.png}
    \caption{Evaluación de técnicas de supresión de errores en simulación ruidosa: \textit{baseline}, DD, PT activo y DD + PT.}
    \label{fig:comparacion_mitigacion_ruido}
\end{figure}

Lejos de ser un fracaso de las técnicas, este comportamiento es exactamente lo que cabía esperar y constituye una observación metodológica importante. Tanto el desacople dinámico como el Pauli twirling están diseñados para combatir fenómenos físicos específicos del hardware real que el modelo de \texttt{FakeTorino} no reproduce, pues su ruido es puramente estocástico. Al insertar pulsos y compuertas auxiliares sobre un ruido que ya es estocástico, las técnicas no encuentran errores coherentes que estabilizar y, en cambio, alargan ligeramente el circuito, sumando una pequeña cantidad de error adicional que las dos métricas registran de manera consistente.
% ============================================================
% 6.2.4  POST-PROCESAMIENTO CLÁSICO BAJO RUIDO
% ============================================================
\subsection{Post-procesamiento bajo ruido: la robustez del filtro clásico}
\label{subsec:fracciones_ruido}

El éxito de la extracción de factores en condiciones tan adversas merece una explicación detallada, porque ilustra una de las virtudes más importantes del algoritmo de Shor en el régimen NISQ: la asimetría entre lo cuántico (frágil) y lo clásico (robusto). Aun cuando el piso de ruido domina visualmente el histograma y la fidelidad de la distribución global cae varios órdenes de magnitud por debajo de la línea base ideal, las cuatro mediciones más probables para $a = 7$ con DD + PT siguen siendo, en orden de frecuencia, exactamente los cuatro picos teóricos del QPE: $y \in \{0, 256, 128, 384\}$, correspondientes a las fases $\varphi \in \{0,\, 1/2,\, 1/4,\, 3/4\}$. Juntas concentran cerca del 47\% de los \textit{shots}; el resto de la masa se distribuye entre cientos de mediciones espurias, cada una de ellas con frecuencia individual muy baja.\\

Físicamente, lo que ocurre es que el ruido no es capaz de generar coherentemente una distribución de picos espurios consistente con un orden incorrecto: lo único que hace es desordenar la masa entre estados aislados, y esos estados aislados no satisfacen la verificación clásica. La cuántica aporta el indicio (la concentración relativa en los picos correctos, capturada conjuntamente por $PST$ y $\mathcal{F}_H$) y la clásica aporta el filtro (la verificación modular), de modo que el algoritmo completo sigue funcionando incluso cuando la calidad de la distribución cuántica es, en términos absolutos, muy pobre.


%% ================================================================
%%  SECCIÓN 6.3 — RESULTADOS EN HARDWARE REAL: IBM TORINO
%% ================================================================
\section{Ejecución experimental en hardware real: IBM Torino}
\label{sec:resultados_hardware}

En esta sección se presentan los resultados obtenidos al ejecutar el algoritmo de Shor para $N = 15$ en el procesador cuántico \texttt{ibm\_torino}, perteneciente a la familia Heron~r1 de IBM. Los experimentos constituyen la Fase~3 del diseño metodológico descrito en la Sección~\ref{sec:protocolo_fases} y responden directamente al objetivo de ejecutar las versiones optimizadas del algoritmo en hardware real, aplicando técnicas de supresión de errores y cuantificando la degradación de la señal cuántica respecto a las simulaciones ideales y ruidosas.
%% ================================================================
%%  6.3.1 — DESCRIPCIÓN DE LA EJECUCIÓN EXPERIMENTAL
%% ================================================================
\subsection{Descripción de la ejecución experimental}
\label{subsec:config_hardware}

Todos los circuitos fueron enviados al procesador \texttt{ibm\_torino} a través del servicio IBM Quantum bajo el plan de acceso abierto (\textit{Open Plan}), utilizando \texttt{Qiskit~Runtime}. La Tabla~\ref{tab:config_hardware} resume la configuración experimental base.

\begin{table}[H]
    \centering
    \caption{Configuración del backend y parámetros experimentales.\\}
    \label{tab:config_hardware}
    \begin{tabular}{ll}
        \hline
        \textbf{Parámetro}             & \textbf{Valor}                                                  \\
        \hline
        Procesador (QPU)               & \texttt{ibm\_torino} (Heron~r1, 133 qubits)                     \\
        Topología                      & Heavy-Hex                                                       \\
        Compuertas nativas 2Q          & CZ                                                              \\
        Compuertas nativas 1Q          & RZ, SX, X                                                       \\
        $T_2$ característico           & $\sim 250\;\mu$s                                                \\
        Fidelidad mediana CZ           & $\sim 99.5\%$ ($\epsilon_{\text{CZ}} \approx 5 \times 10^{-3}$) \\
        Circuito lógico                & RegisterQC ($t = 9$ qubits de control, 4 de trabajo)            \\
        $N$                            & 15                                                              \\
        \texttt{approximation\_degree} & 0.7                                                             \\
        Shots por ejecución            & 4096                                                            \\
        Semilla del transpilador       & 457                                                             \\
        \hline
    \end{tabular}
\end{table}

El parámetro \texttt{approximation\_degree}~$= 0.7$ del transpilador resultó determinante para la viabilidad de la ejecución: este valor permite al compilador realizar una síntesis unitaria aproximada, sacrificando algo de fidelidad a cambio de reducir la profundidad del circuito, en particular el número de compuertas CNOT. 
La Figura~\ref{fig:approx_depth} ilustra cuantitativamente esta reducción para todas las bases del estudio: la profundidad 2Q se comprime aproximadamente en un orden de magnitud al pasar de \texttt{approx}$~=1.0$ a \texttt{approx}$~=0.7$, llevando todos los circuitos al régimen de viabilidad temporal.

\begin{figure}[H]
    \centering
    \includegraphics[width=0.75\textwidth]{ejecucion_real/profundidad2q_vs_base_a.png}
    \caption{Impacto del parámetro \texttt{approximation\_degree} sobre la profundidad de compuertas de dos qubits ($D_{2Q}$) tras la transpilación a la topología Heavy-Hex de \texttt{ibm\_torino}, para las bases $a$ ejecutadas en hardware. La configuración exacta (\texttt{approx}~$=1.0$) produce profundidades entre 430 y 910~capas CZ, mientras que la configuración aproximada (\texttt{approx}~$=0.7$) las reduce a entre 12 y 100~capas, una compresión de un orden de magnitud. Esta reducción es la que hace viable la ejecución dentro del tiempo de coherencia $T_2 \sim 250\,\mu$s del procesador.}
    \label{fig:approx_depth}
\end{figure}

\textbf{Organización del experimento}\\

Por las restricciones de cupo del plan de acceso abierto (cf. Sección~\ref{sec:protocolo_fases}), el barrido factorial completo definido en el Capítulo~\ref{chap:metodologia} no pudo replicarse de manera exhaustiva en hardware.

\begin{table}[H]
    \centering
    \caption{Resumen de \textit{jobs} ejecutados en \texttt{ibm\_torino}.\\}
    \label{tab:jobs_hardware}
    \footnotesize
    \begin{tabular}{llcccccc}
        \hline
        \textbf{Estudio} & \textbf{Config.} & $a$ & \textbf{Depth 2Q} & $PST$ (\%) & $\mathcal{F}_H$ & \textbf{Factores} & \textbf{Job ID}   \\
        \hline
        V1 (inicial)     & PT, opt=3        & 4   & 116               & 75.5       & ---             & $3 \times 5$      & \texttt{d673l15b} \\
        V1 (inicial)     & PT, opt=3        & 7   & 244               & 67.5       & ---             & $3 \times 5$      & \texttt{d672p6gq} \\
        V1 (inicial)     & PT, opt=3        & 14  & 117               & 75.1       & ---             & $1, 15$           & \texttt{d672j2pv} \\
        \hline
        A2 (opt level)   & opt=0            & 4   & 770               & 1.1        & 0.011           & $3 \times 5$      & \texttt{d6agiong} \\
        A2 (opt level)   & opt=1            & 4   & 526               & 3.1        & 0.031           & $3 \times 5$      & \texttt{d6agiu97} \\
        A2 (opt level)   & opt=2            & 4   & 116               & 70.8       & 0.529           & $3 \times 5$      & \texttt{d6agj2p7} \\
        \hline
        A3 (base $a$)    & PT, opt=3        & 2   & 225               & 42.4       & 0.243           & $3 \times 5$      & \texttt{d6aghkvg} \\
        A3 (base $a$)    & PT, opt=3        & 8   & 212               & 79.4       & 0.484           & $3 \times 5$      & \texttt{d6aghrvg} \\
        A3 (base $a$)    & PT, opt=3        & 11  & 75                & 35.5       & 0.309           & $5 \times 3$      & \texttt{d6agi07g} \\
        A3 (base $a$)    & PT, opt=3        & 13  & 212               & 45.3       & 0.214           & $3 \times 5$      & \texttt{d6agi595} \\
        \hline
        A6/A7 (mitg.)    & Sin mitg.        & 4   & 116               & 82.8       & 0.547           & $3 \times 5$      & \texttt{d6agiah7} \\
        A6/A7 (mitg.)    & Solo DD          & 4   & 116               & 20.3       & 0.154           & No hallados       & \texttt{d6agiesn} \\
        A6/A7 (mitg.)    & PT + DD          & 4   & 116               & 0.8        & 0.007           & $3 \times 5$      & \texttt{d6agik97} \\
        \hline
    \end{tabular}
\end{table}

En su lugar, los 13 \textit{jobs} ejecutados se organizaron como cuatro estudios complementarios que, en conjunto, cubren las dimensiones más informativas del diseño: \textbf{V1} actúa como una calibración inicial del flujo de envío, fijando $\texttt{opt}=3$ con \textit{Pauli Twirling} activado para tres bases representativas ($a \in \{4, 7, 14\}$); \textbf{A2} es un barrido del nivel de optimización del transpilador para la base óptima $a = 4$; \textbf{A3} es un barrido de la base $a$ con $\texttt{opt}=3$ fijo; y \textbf{A6/A7} es la sub-exploración del diseño $2 \times 2$ de técnicas de supresión sobre el circuito de $a = 4$, $\texttt{opt}=3$. La organización detallada de los \textit{scripts} y configuraciones empleados en cada estudio se documenta en el Apéndice~\ref{app:scripts-fases}; los identificadores de cada \textit{job} se registran en la Tabla~\ref{tab:jobs_hardware} para garantizar la reproducibilidad de los resultados.

%% ================================================================
%%  6.3.2 — RESULTADOS SIN MITIGACIÓN DE ERRORES
%% ================================================================
\subsection{Resultados sin supresión de errores}
\label{subsec:resultados_crudos_hw}

En primer lugar se analizan los resultados obtenidos con la configuración base sin técnicas de supresión activas, es decir, sin \textit{Dynamical Decoupling} ni \textit{Pauli Twirling}, lo cual permite caracterizar el comportamiento intrínseco del procesador frente al algoritmo de Shor.

\subsubsection{Impacto del nivel de optimización del transpilador}

Según la Tabla \ref{tab:jobs_hardware}, el nivel de optimización del transpilador de Qiskit resultó ser la variable de control con mayor impacto sobre la señal experimental en hardware real. El estudio A2 fijó $a = 4$ y barrió formalmente $\texttt{opt} \in \{0, 1, 2\}$, manteniendo el resto de la configuración constante. La cuarta configuración ($\texttt{opt}=3$) no se ejecutó como parte del estudio A2 propiamente dicho, sino que su valor de referencia para esta base se toma del estudio A6/A7 sin supresión ($PST = 82.8\%$, $\mathcal{F}_H = 0.547$, mismo circuito de 116~CZ que la transpilación con $\texttt{opt}=2$); esta nota metodológica es importante para no sobre-interpretar la comparación entre los cuatro niveles, ya que A2 y A6/A7 corresponden a ventanas de calibración distintas y, por tanto, una pequeña parte de la diferencia $\texttt{opt}=2 \to \texttt{opt}=3$ puede atribuirse a deriva natural y no al transpilador.\\

La transición del nivel~0 al nivel~2 produce una reducción del $84.9\%$ en la profundidad 2Q (de 770 a 116 capas CZ) y del $90.5\%$ en el número de compuertas CZ (de 1220 a 116). Esta reducción se traduce en un incremento de la señal $PST$ de $1.1\%$ a $70.8\%$, y de la fidelidad de $\mathcal{F}_H = 0.011$ a $\mathcal{F}_H = 0.529$. El número de \textit{outcomes} únicos (bitstrings distintos observados) disminuye de 512 a 57, lo cual refleja la concentración de la probabilidad en los picos teóricos en lugar de dispersarse uniformemente por ruido.\\

Con $\texttt{opt} = 0$ (layout trivial, routing \textit{basic}), el circuito requiere 113~SWAPs estimados para adaptar la topología lógica a la malla Heavy-Hex del chip, cada uno descompuesto en 3~compuertas CZ. El circuito resultante ($D_{2Q} = 770$) excede ampliamente el umbral donde la acumulación multiplicativa de errores destruye la señal:
\begin{equation}
    \mathcal{F}_{\text{est}} \sim (1 - \epsilon_{\text{CZ}})^{1220} \approx (0.995)^{1220} \approx 0.002,
    \label{eq:fidelidad_opt0}
\end{equation}
lo que es consistente con la señal observada de $1.1\%$, prácticamente indistinguible de una distribución uniforme.\\

En contraste, $\texttt{opt} = 2$ con SABRE elimina completamente la necesidad de SWAPs adicionales y reduce la profundidad a 116~capas CZ. La estimación correspondiente es:
\begin{equation}
    \mathcal{F}_{\text{est}} \sim (0.995)^{116} \approx 0.56,
    \label{eq:fidelidad_opt2}
\end{equation}
valor coherente con la fidelidad observada $\mathcal{F}_H = 0.529$. Esto confirma que el modelo de decaimiento exponencial $\mathcal{F} \sim (1 - \epsilon_{2q})^{N_g}$ constituye una estimación razonable para circuitos de profundidad moderada en el procesador Heron~r1. La transición de $\texttt{opt}=2$ a $\texttt{opt}=3$ aporta una mejora adicional modesta (de $70.8\%$ a $82.8\%$ en la línea base sin supresión de A6/A7), atribuible al pase de re-síntesis con conciencia del ruido del \textit{backend} que activa el nivel~3 y que reordena las cadenas de CZ minimizando los acoplamientos a qubits con peor fidelidad.

\subsection{Efecto de la base $a$ en la señal de hardware}
\label{subsec:efecto_a}

El estudio A3 evaluó la señal experimental para cuatro bases coprimas adicionales de $N = 15$ ($a \in \{2, 8, 11, 13\}$), todas con $\texttt{opt} = 3$ y $\texttt{approx} = 0.7$. Los resultados se completan con las ejecuciones iniciales V1 ($a \in \{4, 7, 14\}$) y se sintetizan en la Tabla~\ref{tab:bases_hw}. Una precisión metodológica relevante para la lectura de esta tabla: tanto V1 como A3 se ejecutaron con \textit{Pauli Twirling} activado como condición operativa por defecto del flujo de envío adoptado en este trabajo. En consecuencia, la caracterización por base no aísla el efecto de la base \emph{en abstracto}, sino el efecto de la base \emph{en presencia de PT}; esta distinción no compromete la comparación interna entre filas de la tabla (todas comparten la misma condición), pero sí debe tenerse presente al comparar estos números con la línea base sin supresión del estudio A6/A7 para $a=4$ y al interpretar la magnitud absoluta del $PST$ por base.\\

Los factores no triviales $\{3, 5\}$ fueron extraídos exitosamente en seis de las siete bases ejecutadas. La excepción, $a = 14$, produce factores triviales $\{1, 15\}$, como se demostró en la simulación ideal (Sección~\ref{sec:resultados_ideal_puro}).
La Figura~\ref{fig:distribuciones_a3} muestra las distribuciones empíricas obtenidas en hardware para estas cuatro bases superpuestas a la distribución teórica esperada. 
La asimetría entre los picos se explica por el sesgo de relajación térmica ($T_1$) del procesador, que favorece el decaimiento hacia el estado $|0\rangle^{\otimes n}$.

\begin{table}[H]
    \centering
    \caption{Resultados por base $a$ en hardware real (\texttt{ibm\_torino}, $\texttt{opt} = 3$, $\texttt{approx} = 0.7$, PT activado en todos los \textit{jobs}).\\}
    \label{tab:bases_hw}
    \begin{tabular}{cccccccc}
        \hline
        $a$ & $\text{ord}(a,15)$ & Depth 2Q & CZ  & $PST_{\text{HW}}$ (\%) & $\mathcal{F}_H^{\text{HW}}$ & Factores            \\
        \hline
        4   & 2                  & 116      & 116 & 75.5                   & ---                         & $3 \times 5$        \\
        8   & 4                  & 212      & 212 & 79.4                   & 0.484                       & $3 \times 5$        \\
        14  & 2                  & 117      & 117 & 75.1                   & ---                         & $1, 15$ (triviales) \\
        7   & 4                  & 244      & 244 & 67.5                   & ---                         & $3 \times 5$        \\
        13  & 4                  & 212      & 212 & 45.3                   & 0.214                       & $3 \times 5$        \\
        2   & 4                  & 225      & 225 & 42.4                   & 0.243                       & $3 \times 5$        \\
        11  & 2                  & 75       & 76  & 35.5                   & 0.309                       & $5 \times 3$        \\
        \hline
    \end{tabular}
\end{table}

\begin{figure}[H]
    \centering
    \includegraphics[width=\textwidth]{ejecucion_real/distribuciones_picos_hardware.png}
    \caption{Distribuciones de probabilidad observadas en \texttt{ibm\_torino} (verde) frente a la distribución teórica ideal (azul) para las cuatro bases del estudio A3 ($a \in \{2, 8, 11, 13\}$, $\texttt{opt}=3$, $\texttt{approx}=0.7$, PT activado, 4096~\textit{shots}). En cada panel se reportan únicamente los \textit{bitstrings} asociados a las fases teóricas $\varphi = s/r$ del período correspondiente: $\{0, 1/4, 1/2, 3/4\}$ para $r=4$ ($a=2,8,13$) y $\{0, 1/2\}$ para $r=2$ ($a=11$).}
    \label{fig:distribuciones_a3}
\end{figure}

Según la Tabla~\ref{tab:bases_hw}, la mayor $PST$ del estudio corresponde a $a = 8$ ($79.4\%$), y no a $a = 4$. Las tres bases con período $r = 2$ ejecutadas en hardware ($a = 4, 14, 11$) presentan $PST$ que varían entre $35.5\%$ y $75.5\%$, lo que muestra que la ventaja teórica de un período corto no garantiza por sí sola la mejor señal en el procesador. La correlación profundidad–señal no es estrictamente monótona: $a = 4$ y $a = 14$ comparten una profundidad similar ($\sim 116$~CZ) y alcanzan $PST$ muy próximos ($75.5\%$ y $75.1\%$ respectivamente), pero $a = 8$ las supera con casi el doble de profundidad ($212$~CZ). Esto indica que la calidad de la señal también depende de la calibración local de los qubits asignados por el transpilador y de la propagación coherente de errores en cada cadena CZ. El caso de $a = 11$ resulta particularmente ilustrativo: a pesar de tener la menor profundidad 2Q del estudio ($75$~capas), su $PST$ es solo del $35.5\%$, lo que se atribuye a una variabilidad significativa en la fidelidad CZ del cableado físico seleccionado en esa ejecución concreta.

\subsubsection{Caso óptimo: $a = 4$}

El caso $a = 4$ ($r = 2$) resultó ser la elección óptima en general, para la demostración del algoritmo de Shor en hardware real. La Tabla~\ref{tab:distribucion_a4} presenta los 10 \textit{bitstrings} más frecuentes de las 4096 mediciones.

\begin{table}[H]
    \centering
    \caption{Distribución de probabilidad para $a = 4$, $r = 2$ en \texttt{ibm\_torino} (4096 \textit{shots}). Los estados $|000000000\rangle$ y $|100000000\rangle$ corresponden a las fases teóricas $\varphi = 0$ y $\varphi = 1/2$, respectivamente.\\}
    \label{tab:distribucion_a4}
    \begin{tabular}{ccccccc}
        \hline
        \# & $|y\rangle$         & Conteo & $p_{\text{obs}}$ & $\varphi = y/2^9$ & $r_{\text{cand}}$ & Factores              \\
        \hline
        1  & $|000000000\rangle$ & 2897   & 0.707            & 0.000             & ---               & ---                   \\
        2  & $|010000000\rangle$ & 262    & 0.064            & 0.250             & 4                 & triviales             \\
        3  & $|000100000\rangle$ & 202    & 0.049            & 0.063             & ---               & ---                   \\
        4  & $|100000000\rangle$ & 196    & 0.048            & 0.500             & 2                 & $\mathbf{3 \times 5}$ \\
        5  & $|000000010\rangle$ & 110    & 0.027            & 0.004             & ---               & ---                   \\
        \hline
    \end{tabular}
\end{table}

\begin{figure}[H]
    \centering
    \includegraphics[width=0.5\textwidth]{ejecucion_real/distribucion_probabilidad_torino_a4.png}
    \caption{Distribuciones de probabilidad medidas en \texttt{ibm\_torino} para $a=4$ (período $r=2$ extraído correctamente) en el estudio inicial V1 (\textit{Pauli Twirling} activado, $\texttt{opt}=3$, $\texttt{approx}=0.7$, 4096~\textit{shots}). El eje horizontal muestra el resultado de la medición normalizado al intervalo $[0,1]$.}
    \label{fig:distribuciones_v1}
\end{figure}

El pico dominante en $|000000000\rangle$ ($\varphi = 0$, conteo = 2897, $p = 70.7\%$) y el pico en $|100000000\rangle$ ($\varphi = 1/2$, conteo = 196, $p = 4.8\%$) corresponden a las fases teóricas del período $r = 2$. La señal total en estos dos picos, que define el $PST$ para $r = 2$, suma $75.5\%$. A partir de la fase $1/2$ se obtiene la derivación clásica $\gcd(4^{1} \pm 1, 15) = \{3,5\}$, que permanece inalterada.

\subsubsection{Viabilidad temporal del circuito}

Un criterio fundamental para la ejecución en hardware NISQ es que la duración del circuito sea significativamente menor que los tiempos de coherencia del procesador. Para los circuitos ejecutados:
\begin{equation}
    T_{\text{circuito}} \approx D_{2Q} \times \tau_{\text{CZ}} \approx 116 \times 660\;\text{ns} \approx 76.6\;\mu\text{s},
\end{equation}
donde $\tau_{\text{CZ}} \approx 660$~ns es la duración típica de una compuerta CZ en el procesador Heron~r1. El cociente $T_{\text{circuito}} / T_2 \approx 76.6 / 250 \approx 0.31$ confirma que el circuito opera holgadamente dentro del tiempo de coherencia. Para $a = 7$ ($D_{2Q} = 244$), esta razón asciende a $\sim 0.64$, aún dentro del margen viable pero con menor holgura, lo que explica el incremento de ruido observado.

El tiempo cuántico real reportado por el \textit{runtime} fue de 3~segundos para todos los \textit{jobs}, lo que incluye la inicialización de qubits, la ejecución del circuito y la lectura de 4096 \textit{shots}.

%% ================================================================
%%  6.3.3 — IMPACTO DE LAS TÉCNICAS DE SUPRESIÓN
%% ================================================================
\subsection{Impacto de las técnicas de supresión de errores}
\label{subsec:mitigacion_hardware}

El estudio A6/A7 evaluó el efecto del \textit{Dynamical Decoupling} (DD) y el \textit{Pauli Twirling} (PT) sobre la señal experimental en hardware real, comparando tres configuraciones contra la línea base sin supresión, todas con el mismo circuito ($a = 4$, $\texttt{opt} = 3$, $\texttt{approx} = 0.7$, 116~compuertas CZ). Los resultados se presentan en la Tabla~\ref{tab:mitigacion_hw}.

\begin{table}[H]
    \centering
    \caption{Comparación de técnicas de supresión de errores en hardware real (\texttt{ibm\_torino}, $a = 4$, 116~CZ, 4096 \textit{shots}). Se reportan la señal $PST$, la fidelidad $\mathcal{F}_H$, el número de \textit{outcomes} únicos y los factores extraídos.\\}
    \label{tab:mitigacion_hw}
    \begin{tabular}{lccccl}
        \hline
        Configuración  & $PST$ (\%)    & $\mathcal{F}_H$ & Outcomes & Factores     \\
        \hline
        Sin supresión & \textbf{82.8} & \textbf{0.547}  & 30       & $3 \times 5$ \\
        Solo DD (XY4)  & 20.3          & 0.154           & 163      & No hallados  \\
        PT + DD (XY4)  & 0.8           & 0.007           & 469      & $3 \times 5$ \\
        \hline
    \end{tabular}
\end{table}

\subsubsection{Desacoplamiento Dinámico}

El resultado más destacado --- y en apariencia contraintuitivo --- de la Tabla~\ref{tab:mitigacion_hw} es que el \textit{Dynamical Decoupling} con secuencia XY4 degradó severamente la señal experimental: de $PST = 82.8\%$ a $PST = 20.3\%$, una caída de 62.5~puntos porcentuales. La fidelidad se redujo de $\mathcal{F}_H = 0.547$ a $\mathcal{F}_H = 0.154$, y el número de \textit{outcomes} únicos aumentó de 30 a 163, indicando una dispersión sustancial de la probabilidad fuera de los picos teóricos. En esta configuración, el algoritmo fue incapaz de extraer los factores correctos.\\

Este comportamiento se explica por la interacción entre la secuencia DD y la estructura temporal del circuito en el procesador Heron~r1. El DD inserta pulsos compensatorios ($X$-$Y$-$X$-$Y$ en el caso de XY4) durante los periodos de inactividad (\textit{idle times}) de los qubits. Sin embargo, para un circuito de profundidad moderada (116~capas CZ) en un procesador con compuertas nativas de alta fidelidad ($\epsilon_{\text{CZ}} \approx 5 \times 10^{-3}$), los tiempos de inactividad son breves y la decoherencia acumulada durante estos intervalos no constituye la fuente dominante de error. En este régimen, los pulsos DD adicionales introducen errores de calibración propios --- cada pulso de compensación posee un error finito de implementación --- que exceden el beneficio de la refocalización del estado. Adicionalmente, la inserción de pulsos DD puede provocar \textit{crosstalk} con qubits adyacentes en la topología Heavy-Hex, amplificando el ruido correlacionado.

\subsubsection{Pauli Twirling}

El protocolo experimental del Capítulo~\ref{chap:metodologia} (Sección~\ref{sec:supresion_implementada}) prescribe un diseño factorial $2 \times 2$ que cruza la activación o desactivación de DD y PT, dando lugar a cuatro celdas: \emph{sin supresión}, \emph{solo DD}, \emph{solo PT} y \emph{DD+PT}. En el simulador con ruido (Fase~2) las cuatro celdas se ejecutaron formalmente bajo idénticas condiciones (Sección~\ref{subsec:efecto_mitigacion_ruido}). En hardware, en cambio, el cupo del plan abierto solo permitió ejecutar tres celdas formales dentro del estudio A6/A7 (\emph{sin supresión}, \emph{solo DD} y \emph{DD+PT}, todas reportadas en la Tabla~\ref{tab:mitigacion_hw}); la cuarta celda, \emph{solo PT}, se infiere indirectamente comparando la línea base sin supresión de A6/A7 ($PST = 82.8\%$) con la ejecución V1 inicial ($PT$ activado, sin DD, $PST = 75.5\%$) sobre el mismo circuito de 116~CZ:

\begin{table}[H]
    \centering
    \caption{Efecto del \textit{Pauli Twirling} para $a = 4$ en \texttt{ibm\_torino} (116~CZ, 4096 \textit{shots}).\\}
    \label{tab:pt_effect}
    \begin{tabular}{lcc}
        \hline
        Configuración          & $PST$ (\%) & Factores     \\
        \hline
        Sin mitigación (A6/A7) & 82.8       & $3 \times 5$ \\
        Solo PT (V1, inferida) & 75.5       & $3 \times 5$ \\
        \hline
    \end{tabular}
\end{table}

La celda \emph{Solo PT} inferida de V1 da $PST = 75.5\%$, ligeramente por debajo de la línea base sin supresión ($82.8\%$). La diferencia, de $\Delta PST \approx -7$~pp, debe leerse con cautela: V1 y A6/A7 fueron ejecutados en fechas distintas, con calibraciones del procesador potencialmente diferentes, y la deriva natural del \textit{backend} entre dos ventanas de calibración produce, en circuitos comparables, fluctuaciones que pueden acumular varios puntos porcentuales incluso bajo condiciones nominalmente idénticas. En consecuencia, parte de la brecha observada puede atribuirse a deriva de calibración y no exclusivamente al efecto del PT, y los 4096~\textit{shots} disponibles no permiten afirmar significancia estadística sobre el signo del efecto. La conclusión cualitativa, no obstante, es que \emph{el PT no produce una mejora robusta de la señal en este régimen}, en línea con la observación previa del simulador con ruido: para un circuito de profundidad $D_{2Q} = 116$ en el procesador Heron~r1, la contribución de los errores coherentes al error total es modesta, de modo que la conversión que el PT realiza de error coherente en ruido estocástico no se traduce en un beneficio observable sobre el $PST$.

\subsubsection{Combinación PT + DD}

La combinación de ambas técnicas produjo el peor resultado de todas las configuraciones: $PST = 0.76\%$ con $\mathcal{F}_H = 0.007$, una distribución efectivamente indistinguible del ruido uniforme. Se observaron 469~\textit{outcomes} únicos de los 512 posibles ($2^9$), lo cual confirma la pérdida casi total de la estructura periódica del QPE. El hecho de que los factores $\{3, 5\}$ fueran formalmente extraíbles en esta configuración (por la coincidencia estadística de que $r = 2$ fue identificado como candidato) no debe interpretarse como un éxito: la señal en los picos teóricos ($0.22\%$ en $|000000000\rangle$ y $0.54\%$ en $|100000000\rangle$) es despreciable frente al fondo de ruido.\\

Este resultado demuestra que la combinación PT + DD no es aditiva en sus beneficios: el DD amplifica el ruido por los mecanismos ya descritos, y el PT, al aleatorizar las compuertas de Pauli alrededor de cada CZ, modifica la estructura de cancelación de errores que el circuito original podría exhibir naturalmente, eliminando posibles interferencias destructivas de errores que serían favorables.

\subsubsection{Contraste con la simulación ruidosa}

La Tabla~\ref{tab:contraste_mitigacion} contrasta el efecto de las cuatro variantes de supresión entre la simulación con ruido FakeTorino (Sección~\ref{subsec:efecto_mitigacion_ruido}) y la ejecución en hardware real, sobre el mismo circuito $a = 4$, $\texttt{opt}=3$. La comparación es aproximadamente directa porque el \textit{snapshot} de FakeTorino se construyó a partir del propio procesador \texttt{ibm\_torino}, de modo que la diferencia entre las dos columnas mide específicamente lo que el \textit{fake backend} no modela.

\begin{table}[H]
    \centering
    \caption{Contraste del efecto de la mitigación entre el simulador ruidoso (FakeTorino, Fase~2) y el hardware real (\texttt{ibm\_torino}, Fase~3), para $a = 4$, $\texttt{opt}=3$. La celda \emph{Solo PT} en HW se obtiene por inferencia indirecta de V1.\\}
    \label{tab:contraste_mitigacion}
    \begin{tabular}{lcccc}
        \hline
        Configuración           & $PST_{\text{FT}}$ (\%) & $\mathcal{F}_H^{\text{FT}}$ & $PST_{\text{HW}}$ (\%) & $\mathcal{F}_H^{\text{HW}}$ \\
        \hline
        Sin mitigación          & 3.52                   & 0.0039                      & \textbf{82.8}          & \textbf{0.547}              \\
        Solo DD                 & 3.91                   & 0.0073                      & 20.3                   & 0.154                       \\
        Solo PT                 & 2.34                   & 0.0057                      & 75.5 (V1)              & ---                         \\
        DD + PT                 & 3.32                   & 0.0029                      & 0.8                    & 0.007                       \\
        \hline
    \end{tabular}
\end{table}

En el simulador FakeTorino, las cuatro variantes producen señales prácticamente indistinguibles entre sí ($PST \in [2.34, 3.91]$~\%, $\mathcal{F}_H \in [0.003, 0.007]$), confirmando la observación de la Sección~\ref{subsec:efecto_mitigacion_ruido}: el modelo de ruido del \textit{fake backend} es estocástico.\\

En hardware, en cambio, las diferencias entre celdas son de aproximadamente dos órdenes de magnitud: la línea base sin supresión alcanza el $82.8\%$, el PT solo se sitúa ligeramente por debajo ($75.5\%$, atribuible en parte a deriva de calibración), el DD la colapsa al $20.3\%$ y la combinación DD+PT la lleva por debajo del $1\%$. Esta sensibilidad es exactamente la firma del régimen NISQ que el \textit{fake backend} no captura: una vez que el ruido coherente y el \textit{crosstalk} entran en juego, la elección de la técnica de supresión deja de ser irrelevante y pasa a ser determinante. La conclusión metodológica es que, para esta clase de decisiones, la simulación con ruido funciona como un filtro de errores burdo, que no sustituye a la ejecución en hardware real.

%% ================================================================
%%  6.3.4 — ANÁLISIS Y DISCUSIÓN
%% ================================================================
\subsection{Análisis y discusión}
\label{subsec:discusion_hardware}

\subsubsection{Degradación de la señal: simulación ideal $\to$ FakeTorino $\to$ IBM Torino}

La Tabla~\ref{tab:degradacion_cascada} cuantifica la degradación de la señal cuántica a lo largo de los tres entornos de ejecución. Para las bases que cuentan con medición tanto en Fase~2 como en Fase~3 ($a \in \{4, 7\}$), la lectura de la cascada es directa; para las bases incluidas únicamente en el barrido de hardware ($a \in \{2, 8, 11, 13\}$) la columna FakeTorino se marca con ``---'', indicando ausencia de dato y no inferencia.

\begin{table}[H]
    \centering
    \caption{Degradación de la señal $PST$ (\%) desde la simulación ideal hasta el hardware real, $\texttt{opt} = 3$, $\texttt{approx} = 0.7$. La columna FakeTorino corresponde al \textit{baseline} sin mitigación de la Fase~2 y solo está disponible para $a \in \{4, 7\}$. La columna HW agrupa V1 ($a = 4, 7$, PT activado) y A3 ($a \in \{2, 8, 11, 13\}$, PT activado).\\}
    \label{tab:degradacion_cascada}
    \begin{tabular}{cccccccc}
        \hline
        $a$ & $r$ & $PST_{\text{ideal}}$ & $PST_{\text{FT}}$ & $PST_{\text{HW}}$ & $\Delta_{\text{ideal} \to \text{FT}}$ & $\Delta_{\text{ideal} \to \text{HW}}$ & $\mathcal{F}_H^{\text{HW}}$ \\
        \hline
        4   & 2 & 100.0 & 4.3   & 75.5 & 95.7 & 24.5 & ---   \\
        7   & 4 & 100.0 & 59.2  & 67.5 & 40.8 & 32.5 & ---   \\
        8   & 4 & 100.0 & ---   & 79.4 & ---  & 20.6 & 0.484 \\
        2   & 4 & 100.0 & ---   & 42.4 & ---  & 57.6 & 0.243 \\
        11  & 2 & 100.0 & ---   & 35.5 & ---  & 64.5 & 0.309 \\
        13  & 4 & 100.0 & ---   & 45.3 & ---  & 54.7 & 0.214 \\
        \hline
        \multicolumn{2}{c}{\textbf{Promedio}} & 100.0 & 31.8\textsuperscript{*} & 57.6 & 68.3\textsuperscript{*} & 42.4 & 0.310\textsuperscript{**} \\
        \hline
        \multicolumn{8}{l}{\footnotesize \textsuperscript{*} Promedio sobre las dos bases con dato FT. \textsuperscript{**} Promedio sobre las cuatro bases con $\mathcal{F}_H$ reportada.} \\
    \end{tabular}
\end{table}

\textit{El hallazgo central:} el \textit{fake backend} sobreestima la degradación en el régimen de circuitos cortos. En las dos bases con medición en ambos entornos, el hardware real \emph{supera} a la simulación con ruido. Para $a = 4$ (116~CZ en HW), la inversión es marcada: $PST_{\text{FT}} = 4.3\%$ frente a $PST_{\text{HW}} = 75.5\%$. Para $a = 7$ (244~CZ en HW), la inversión es leve pero del mismo signo: $PST_{\text{FT}} = 59.2\%$ frente a $PST_{\text{HW}} = 67.5\%$. La consistencia de la dirección sugiere que se trata de un patrón sistemático, no de una fluctuación, y admite una explicación física en tres componentes que no son mutuamente excluyentes:

\begin{enumerate}
    \item Como se declaró en la Sección~\ref{sec:protocolo_fases}, el \textit{snapshot} de FakeTorino reproduce despolarización, relajación térmica $T_1/T_2$ y matrices de confusión de lectura, pero no contiene información sobre errores coherentes ni \textit{crosstalk} dinámico. Este es, por construcción, un modelo de ruido \emph{aditivo y estocástico}: cada compuerta degrada la señal multiplicativamente y los errores no se cancelan entre sí, lo que tiende a producir cotas pesimistas en circuitos donde la compensación coherente es posible.
    \item El \textit{snapshot} reporta tasas de error de la ventana de calibración en que fue tomado; la ejecución en \texttt{ibm\_torino} se realizó con la calibración vigente del día, que para los qubits y enlaces seleccionados por SABRE pudo presentar fidelidades CZ superiores a las del \textit{snapshot}. La deriva natural de calibración entre dos ventanas, lejos de ser un artefacto, es la regla en NISQ.
    \item En un circuito como el de $a = 4$ con 116~CZ, las sobre-rotaciones sistemáticas y las fases espúreas asociadas a cada compuerta pueden interferir destructivamente entre sí en la cadena, recuperando parte de la señal que el modelo estocástico, por construcción, da por perdida. Este mecanismo se atenúa progresivamente al aumentar la profundidad, lo cual es consistente con la observación empírica: la inversión es muy fuerte en $a = 4$ ($\Delta_{\text{HW-FT}} \approx +71$~pp, 116~CZ) y se reduce a un valor pequeño en $a = 7$ ($\Delta_{\text{HW-FT}} \approx +8$~pp, 244~CZ).
\end{enumerate}

La lectura conjunta de estos tres factores conduce a un hallazgo metodológico de interés para futuros trabajos en NISQ: \emph{el \textit{fake backend} no es un predictor optimista de su propio hardware en el régimen de circuitos cortos; es un predictor pesimista}. Cualquier estimación de viabilidad basada en el simulador con ruido tenderá a subestimar la señal real cuando el circuito quepa cómodamente dentro del tiempo de coherencia, y las decisiones de ingeniería tomadas exclusivamente sobre el simulador --- como descartar una base $a$ por baja $PST$ simulada --- pueden ser injustificadamente conservadoras.\\

\textit{Comportamiento de las bases sin medición en Fase~2:} para $a \in \{2, 8, 11, 13\}$ no se dispone de medición en FakeTorino, lo que impide cuantificar la cascada completa. La columna $PST_{\text{HW}}$ permite, no obstante, ordenar las bases por dificultad efectiva en hardware: $a = 8$ y $a = 4$ encabezan el ranking con $PST > 75\%$, mientras que $a = 11$ ($D_{2Q} = 75$, la menor del estudio) aparece en el extremo opuesto con $PST = 35.5\%$. Esta no-monotonía respecto a la profundidad fue ya discutida y se atribuye a la heterogeneidad de la fidelidad CZ a lo largo de las cadenas Heavy-Hex que SABRE selecciona en cada ejecución.\\

La señal promedio en hardware real ($\overline{PST}_{\text{HW}} = 57.6\%$ sobre las seis bases con factores no triviales) es suficiente para la extracción exitosa de los factores de $N = 15$ en todos los casos, confirmando que el algoritmo de Shor opera en un régimen donde la señal cuántica es distinguible del ruido para esta instancia pedagógica.


\chapter{Conclusiones}
\label{cap:conclusiones}

El presente trabajo de grado se propuso evaluar la viabilidad técnica y el desempeño del algoritmo de Shor para la factorización del entero $N = 15$ en procesadores cuánticos NISQ de IBM, comparando sistemáticamente los resultados obtenidos en tres entornos de ejecución: simulación ideal, simuladores con modelos de ruido calibrados (\textit{Fake Backends}) y el procesador real \texttt{ibm\_torino} (familia Heron r1, 133~qubits superconductores). A lo largo de los capítulos precedentes se documentó la implementación completa del algoritmo ---desde los fundamentos teóricos hasta la ejecución experimental en hardware---, generando evidencia cuantitativa sobre la degradación de la señal cuántica y la efectividad de las estrategias de optimización y supresión de errores. En este capítulo final se sintetizan los hallazgos principales, se discuten las limitaciones del estudio y se proponen perspectivas futuras.

\section{Resultados principales}

La factorización $N = 15 = 3 \times 5$ se demostró experimentalmente en el procesador \texttt{ibm\_torino} con una señal máxima de $PST = 82.8\%$, correspondiente a la línea base sin supresión del estudio A6/A7 (\texttt{job\_id}~\texttt{d6agiah7}, $\mathcal{F}_H = 0.547$), y con una tasa de éxito cercana al $90\%$ en la extracción de los factores no triviales sobre las ejecuciones con bases coprimas válidas (Tabla~\ref{tab:jobs_hardware}). Hasta donde se tiene conocimiento, ésta constituye la primera demostración experimental del algoritmo de Shor en hardware cuántico real reportada desde la Universidad de Nariño.\\

La Transformada Cuántica de Fourier ($n = 9$ qubits) y los bloques controlados de exponenciación modular se implementaron en Qiskit y se verificaron numéricamente con errores de unitariedad $\lVert U^{\dagger} U - I \rVert_F < 10^{-13}$ en todas las variantes evaluadas (Tabla~\ref{tab:qft_metricas}). El mapeo a la topología Heavy-Hex de FakeTorino evidenció un \textit{overhead} sobre el conteo de compuertas de dos qubits que llega hasta un factor de $\sim 3.9\times$ con mapeo trivial respecto a la conectividad \textit{all-to-all} (Tabla~\ref{tab:transpilacion_faketorino}). Sobre esa base, el algoritmo SABRE redujo de forma sistemática el sobrecosto, con disminuciones de entre el $11\%$ y el $18\%$ en profundidad de dos qubits y entre el $19\%$ y el $27\%$ en el conteo de compuertas de dos qubits para las bases evaluadas (Tabla~\ref{tab:reduccion_sabre}).\\

Más allá de la transpilación, el factor con mayor impacto sobre la señal experimental resultó ser el nivel de optimización del transpilador de Qiskit. El estudio A2 (Tabla~\ref{tab:opt_level_hardware}, Sección~\ref{subsec:resultados_crudos_hw}) registró una transición de $PST = 1.1\%$ (\texttt{opt}~$=0$, 770~capas CZ y 1220~compuertas CZ con 113~SWAPs estimadas) a $PST = 70.8\%$ (\texttt{opt}~$=2$, 116~capas CZ sin SWAPs adicionales), una ganancia de aproximadamente $64\times$ obtenida al sólo modificar este parámetro. La transición adicional de \texttt{opt}~$=2$ a \texttt{opt}~$=3$ con \textit{Pauli Twirling} desactivado (línea base de A6/A7) elevó la señal hasta el máximo de $82.8\%$, atribuible al pase de re-síntesis con conciencia del ruido del \textit{backend}. Esta evidencia cuantifica por qué la compilación cuántica no constituye una optimización accesoria, sino una pieza estructural del algoritmo en la era NISQ.\\

La estrategia de pre-procesamiento clásico propuesta en la literatura reciente \cite{villarreal} se validó experimentalmente bajo condiciones equivalentes. El módulo desarrollado en Python caracterizó exhaustivamente el grupo $\mathbb{Z}_{15}^*$ e identificó las bases $a \in \{4, 11, 14\}$ con período $r = 2$ como las de menor complejidad estructural de circuito (Sección~\ref{sec:preprocesamiento_parametros}, Tabla~\ref{tab:caracterizacion_a}). En la serie V1 (\texttt{opt}~$=3$, PT activado), la base $a = 4$ ($r = 2$, 116~CZ) alcanzó $PST = 75.5\%$ frente a $a = 7$ ($r = 4$, 244~CZ) con $PST = 67.5\%$ ---una diferencia de 8.0~puntos porcentuales atribuible directamente a la profundidad del circuito (Tabla~\ref{tab:bases_hw})---, lo que convierte una predicción teórica sobre el escalamiento con $r$ en una ventaja medible y reproducible en hardware NISQ.\\

La cadena completa de degradación se cuantificó utilizando la fidelidad de Hellinger $\mathcal{F}_H$ y la probabilidad de éxito experimental $PST$. Para las dos bases con medición tanto en FakeTorino como en hardware real ($a \in \{4, 7\}$), la señal pasó de $PST = 100\%$ en simulación ideal a un promedio de $31.8\%$ en \textit{fake backend} y $57.6\%$ en hardware real, con una fidelidad promedio en hardware $\overline{\mathcal{F}}_H \approx 0.310$ sobre las cuatro bases con $\mathcal{F}_H$ reportada (Tabla~\ref{tab:degradacion_cascada}). El modelo de decaimiento exponencial $\mathcal{F} \sim (1 - \epsilon_{\text{CZ}})^{N_g}$ con $\epsilon_{\text{CZ}} \approx 5 \times 10^{-3}$ resultó un estimador razonable para circuitos de profundidad moderada en el procesador Heron r1: para $a = 4$ con 116~CZ predice $\mathcal{F}_{\text{est}} \approx 0.56$, frente a la fidelidad observada $\mathcal{F}_H = 0.529$ en el estudio A2 con \texttt{opt}~$=2$.\\

Quizá el hallazgo más inesperado, y a la vez más informativo desde un punto de vista físico, se obtuvo al evaluar las técnicas de supresión de errores en hardware real. El Desacoplamiento Dinámico con secuencia XY4, lejos de mejorar la señal, la degradó en $62.5$~puntos porcentuales: el $PST$ se redujo de $82.8\%$ (línea base) a $20.3\%$, la fidelidad cayó de $\mathcal{F}_H = 0.547$ a $0.154$, y la dispersión pasó de 30 a 163~\textit{outcomes} únicos (Tabla~\ref{tab:mitigacion_hw}). La combinación PT$+$DD acentuó este efecto hasta llevar la señal por debajo del $1\%$ ($PST = 0.8\%$, $\mathcal{F}_H = 0.007$, 469~\textit{outcomes} únicos sobre los 512 posibles), produciendo una distribución prácticamente indistinguible del ruido uniforme. Como se discutió en la Sección~\ref{subsec:mitigacion_hardware}, este comportamiento se atribuye a que, para circuitos de profundidad moderada ($D_{2Q} \sim 116$) en procesadores con compuertas nativas de alta fidelidad, los intervalos de inactividad son demasiado breves para que la refocalización compense los errores de calibración propios de los pulsos de DD; adicionalmente, el PT modifica patrones de cancelación coherente potencialmente favorables. La implicación práctica es directa: la aplicación indiscriminada de técnicas de supresión puede ser contraproducente, y su eficacia debe evaluarse empíricamente para cada régimen operativo.\\

A los resultados anteriores se añade una observación de carácter metodológico: en el régimen de circuitos cortos, los \textit{Fake Backends} subestiman la calidad del hardware real. El contraste de la Tabla~\ref{tab:contraste_mitigacion} muestra que, para el mismo circuito $a = 4$ con \texttt{opt}~$=3$, FakeTorino predice $PST_{\text{FT}} = 3.5\%$ sin supresión, mientras que \texttt{ibm\_torino} entrega $PST_{\text{HW}} = 82.8\%$. Además, las cuatro variantes de supresión producen señales prácticamente indistinguibles en el simulador ($PST \in [2.34, 3.91]\%$, $\mathcal{F}_H \in [0.003, 0.007]$), pero difieren en aproximadamente dos órdenes de magnitud en hardware. Esta firma del régimen NISQ ---ruido coherente, deriva de calibración entre ventanas y posibles cancelaciones destructivas en cadenas largas de CZ--- no la captura el modelo estocástico del \textit{snapshot}, lo que tiene implicaciones claras para el uso de simuladores ruidosos como filtro previo a la ejecución en hardware.

\section{Limitaciones del trabajo}

Es importante reconocer las limitaciones que acotan el alcance de los resultados obtenidos.\\

La limitación más fundamental es la restricción al caso $N = 15$, la instancia pedagógica más pequeña del algoritmo de Shor. Si bien este caso es canónico en la literatura y permite ilustrar de forma rigurosa la cadena completa de factorización, la extrapolación a enteros de relevancia criptográfica resulta inviable en el hardware actual.\\

En segundo lugar, el número de ejecuciones en hardware real estuvo condicionado por el cupo del IBM Quantum Open Plan (10~minutos mensuales de tiempo de procesador). Si bien los 13~\textit{jobs} ejecutados cubrieron de manera coherente los estudios planificados, una caracterización estadística más robusta (por ejemplo, repeticiones múltiples de cada configuración para cuantificar la variabilidad inter-ejecución, o réplicas de la línea base en distintos días) habría requerido un volumen de acceso fuera del alcance del plan abierto. La anomalía observada para $a = 11$ ($PST = 35.5\%$ a pesar de tener la menor profundidad 2Q del estudio, 75~capas, Tabla~\ref{tab:bases_hw}) ilustra esta limitación: sin réplicas, no es posible distinguir formalmente entre variabilidad en la calidad de los qubits asignados por SABRE y un efecto sistemático asociado a la base.\\

En tercer lugar, los simuladores con modelos de ruido (\textit{Fake Backends}) mostraron limitaciones importantes como predictores del comportamiento en hardware real. Como ya se documentó en las Secciones~\ref{subsec:efecto_mitigacion_ruido} y~\ref{subsec:mitigacion_hardware}, el \textit{snapshot} de FakeTorino reproduce despolarización, relajación térmica $T_1/T_2$ y matrices de confusión de lectura, pero no captura \textit{crosstalk} dinámico entre qubits vecinos, errores coherentes de compuerta ni la deriva temporal de las calibraciones. El resultado neto, sintetizado en la Tabla~\ref{tab:contraste_mitigacion}, es que el \textit{snapshot} resulta ser un predictor pesimista en circuitos cortos y produce un veredicto cualitativa y cuantitativamente incorrecto sobre la efectividad de DD y PT.\\

Finalmente, la implementación de la exponenciación modular se basó en matrices unitarias precomputadas (clase \texttt{UnitaryGate} de Qiskit). Esta elección fue deliberada y, dado el tamaño del registro objetivo ($t_t = 4$ qubits), permitió minimizar la profundidad del circuito y centrar la atención en las preguntas de transpilación y supresión que son el foco del trabajo. No obstante, no constituye una implementación escalable: una demostración general del algoritmo requeriría descomponer la exponenciación modular en compuertas elementales reversibles (Tabla~\ref{tab:exp_mod_bloques}), lo que incrementaría sustancialmente la profundidad y, por tanto, agudizaría los efectos de ruido aquí caracterizados.

\section{Perspectivas futuras}

Los resultados obtenidos abren varias líneas de investigación viables para futuros trabajos de grado o proyectos en la Universidad de Nariño.\\

Una extensión natural es la implementación de la estimación de fase iterativa (\textit{Iterative QPE}), que recicla un único qubit de control mediante mediciones intermedias y realimentación clásica, comprimiendo el registro de conteo de $t = 9$ qubits a un solo qubit \cite{Nielsen2024}. Esta técnica reduciría drásticamente la conectividad requerida y la profundidad efectiva del circuito, y abre la posibilidad concreta de ejecutar Shor para numeros mas grandes en el hardware actual.\\

Otra dirección prometedora, sugerida directamente por los resultados aquí obtenidos, es el estudio de Desacoplamiento Dinámico adaptativo. La caída de $PST = 82.8\%$ a $20.3\%$ con DD~XY4 (Tabla~\ref{tab:mitigacion_hw}) muestra que esta secuencia no es adecuada para circuitos de profundidad moderada con compuertas nativas de alta fidelidad. Sería valioso evaluar secuencias más cortas (como el eco de espín simple $X$-$X$) o protocolos de DD que se activen únicamente cuando los \textit{idle times} superen un umbral determinado por las tasas de error de compuerta. Un análisis análogo aplica al \textit{Pauli Twirling}: el régimen en el que su aleatorización resulta realmente beneficiosa puede caracterizarse experimentalmente como función de la profundidad del circuito y de la calidad de calibración.\\

También resultaría valioso extender el estudio a procesadores de generaciones más recientes de la hoja de ruta de IBM (Heron r2, Flamingo) o a otras arquitecturas con conectividad distinta, cuyas mejoras en fidelidad de compuerta y tiempos de coherencia podrían ampliar el umbral de profundidad viable y acercar la posibilidad de factorizar números mas grandes en hardware real.\\

Finalmente, la metodología reproducible documentada en este trabajo ---con el repositorio público de código, los \texttt{job\_id} de IBM Quantum y los protocolos de transpilación--- sienta las bases para que futuros estudiantes de la Universidad de Nariño repliquen, extiendan y mejoren estos experimentos, contribuyendo a la consolidación de competencias en computación cuántica experimental en la institución.

\section{Reflexión final}

Este trabajo de grado se concibió con el propósito de tender un puente entre la formulación teórica del algoritmo de Shor y su realidad práctica en el hardware cuántico contemporáneo. Los resultados aquí reportados confirman que, en la era NISQ, la ejecución exitosa de algoritmos cuánticos profundos no depende únicamente de la calidad del hardware ni de la corrección formal del algoritmo, sino de un entendimiento integral que abarca la ingeniería de compilación, la selección inteligente de parámetros clásicos, la comprensión del modelo de ruido del dispositivo y la evaluación empírica de las técnicas de supresión de errores.\\

La demostración de que el algoritmo de Shor puede factorizar $N = 15 = 3 \times 5$ con una señal del $82.8\%$ en un procesador real de IBM, accesible de forma remota desde Pasto, ilustra tanto el potencial como las limitaciones actuales de la computación cuántica. La brecha entre factorizar $N = 15$ (116~compuertas CZ, $T_{\text{circuito}}/T_2 \approx 0.31$) y un entero criptográficamente relevante (miles de millones de compuertas lógicas, qubits lógicos tolerantes a fallos) sigue siendo inmensa. Pero cada mejora en la fidelidad de las compuertas, cada reducción del \textit{overhead} de transpilación y cada avance en las técnicas de supresión y mitigación de errores acortan esa distancia.\\

Más allá del rigor técnico, este proyecto significó el paso del aula a la consola: del diagrama del circuito en el cuaderno al envío real del \textit{job} y a la espera, días después, del histograma medido a unos $50$~mK. Aprender a leer ese histograma ---a distinguir cuánto del pico es señal y cuánto es deriva de calibración, a aceptar que un experimento aparentemente contraintuitivo como la degradación con DD también es un resultado válido--- fue, en buena medida, parte central del aprendizaje que dejó el trabajo.\\

Se confía en que este aporte a la Universidad de Nariño y a su comunidad académica constituya un precedente concreto de investigación experimental en computación cuántica: una metodología documentada, un conjunto de datos reproducibles ---incluidos todos los \texttt{job\_id} y archivos de configuración--- y la evidencia de que es posible realizar investigación significativa en esta disciplina con los recursos de acceso abierto disponibles. Se espera que este estudio sirva como punto de partida para futuros trabajos que continúen explorando las fronteras de la computación cuántica desde la institución.

% =========================================================================
% BLOQUE B — Apéndice A "Preliminares matemáticos y formalismo cuántico"
% =========================================================================

\appendix
\chapter{Preliminares matemáticos y formalismo cuántico}
\label{app:ruido}
 
Este apéndice complementa la Sección~\ref{sec:nisq} con los elementos formales mínimos del modelo de ruido empleados para describir, simular y mitigar los errores observados en la ejecución del algoritmo de Shor sobre los procesadores de IBM Quantum.
 
\section{Canales cuánticos y representación de Kraus}
\label{app:kraus}
 
El formalismo general para describir procesos cuánticos ruidosos es el de los \textit{canales cuánticos}, mapas completamente positivos y que preservan la traza (CPTP) \cite{Nielsen2024, Watrous2018}. Todo canal cuántico $\mathcal{E}$ admite una \textit{representación de operadores de Kraus}
\begin{equation}
    \mathcal{E}(\rho) = \sum_{k} E_k\,\rho\,E_k^{\dagger}, \qquad \sum_k E_k^{\dagger}E_k = I,
    \label{eq:kraus}
\end{equation}
donde $\{E_k\}$ son los \textit{operadores de Kraus} \cite{Nielsen2024}. La evolución unitaria ideal corresponde al caso de un solo operador $E_0 = U$; los procesos no unitarios (relajación, desfase, despolarización) requieren múltiples operadores de Kraus.
 
\section{Canal despolarizante post-compuerta}
\label{app:despolarizante}
 
En un dispositivo NISQ, la aplicación de una compuerta cuántica $U$ no produce la transformación ideal $\rho \mapsto U\rho U^{\dagger}$, sino una versión ruidosa $\mathcal{E}_{\text{gate}}$ \cite{Preskill}. Componiendo el canal despolarizante de la Ec.~\ref{eq:canal_despolarizante} con la operación ideal se obtiene el modelo de error post-compuerta empleado en este trabajo:
\begin{equation}
    \mathcal{E}_{\text{gate}}(\rho) = \mathcal{E}_{\text{depol}}\big(U\rho U^{\dagger}\big) = (1-p)\,U\rho U^{\dagger} + \frac{p}{d}\,I.
    \label{eq:canal_gate}
\end{equation}
El parámetro $p$ se identifica con la tasa de error de compuerta reportada por IBM para cada qubit y par de qubits en sus calibraciones diarias \cite{AbuGhanem2025}. En las simulaciones clásicas con el \textit{noise model} de Qiskit Aer, este canal se aplica tras cada compuerta del circuito transpilado.



% =========================================================================
% BLOQUE B — Material migrado al Apéndice A (continuación)
% =========================================================================

% Las siguientes secciones se insertan dentro del mismo
% \chapter{Preliminares matemáticos y formalismo cuántico}
% creado para §4.2, después de las secciones ya existentes.

Las siguientes secciones recogen el desarrollo formal de los resultados algebraicos empleados en §\ref{sec:qft}, §\ref{sec:order_finding} y §\ref{sec:reduccion_clasica}, presentados allí en forma compacta.


\section{Exponenciación modular}
\label{app:exponenciacion-modular}

El efecto combinado de las operaciones controladas $c$-$U_{x,N}^{2^j}$ sobre los dos registros del algoritmo de búsqueda de orden es la transformación \cite{Nielsen2024}
\begin{align}
    |z\rangle|y\rangle & \to |z\rangle\, U_{x,N}^{z_{t}2^{t-1}} \cdots U_{x,N}^{z_1 2^0}\,|y\rangle \nonumber           \\
                       & = |z\rangle\, |x^{z_{t}2^{t-1}} \times \cdots \times x^{z_1 2^0}\, y \pmod{N}\rangle \nonumber \\
                       & = |z\rangle\, |x^z\, y \pmod{N}\rangle,
    \label{eq:exponenciacion_modular}
\end{align}
donde $z = z_{t}2^{t-1} + \cdots + z_0$ es la representación binaria del contenido del primer registro. La secuencia de operaciones controladas-$U^{2^j}$ equivale a multiplicar el contenido del segundo registro por la exponencial modular $x^z \pmod{N}$.\\

\textbf{Estrategia de cuadrado repetido}\\ 

El cálculo de $x^z \pmod{N}$ se descompone en dos etapas \cite{Nielsen2024}:
\begin{enumerate}
    \item \textit{Pre-computación de potencias.} se calculan sucesivamente $x^2 \pmod{N}$, $x^4 = (x^2)^2 \pmod{N}$, $x^8 = (x^4)^2 \pmod{N}$, y así hasta $x^{2^{t-1}} \pmod{N}$. Esto requiere $t - 1 = \mathcal{O}(L)$ operaciones de cuadrado modular, cada una con un costo de $\mathcal{O}(L^2)$ operaciones elementales, para un total de $\mathcal{O}(L^3)$.
    \item \textit{Multiplicaciones modulares.} una vez obtenidas las potencias, el exponencial se ensambla mediante la factorización
          \begin{equation}
              x^z \pmod{N} = \big(x^{z_{t}2^{t-1}} \pmod{N}\big) \big(x^{z_{t-1}2^{t-2}} \pmod{N}\big) \cdots \big(x^{z_1 2^{0}} \pmod{N}\big).
              \label{eq:repeated_squaring}
          \end{equation}
          Esto requiere $t - 1$ multiplicaciones modulares adicionales, cada una con costo $\mathcal{O}(L^2)$, para un total de $\mathcal{O}(L^3)$.
\end{enumerate}
El costo total de la exponenciación modular es, por tanto, $\mathcal{O}(L^3)$ compuertas, dominando el costo computacional del algoritmo completo \cite{Nielsen2024}.\\

\textbf{Requisito de reversibilidad}\\ 

Un computador cuántico exige que todas las operaciones sean unitarias. El circuito clásico de multiplicación modular $|z\rangle|y\rangle \to |z\rangle|x^z y \bmod N\rangle$ se implementa de forma reversible mediante la técnica de descomputación: se computa $x^z \pmod{N}$ en un registro auxiliar, se multiplica reversiblemente el segundo registro por este valor, y finalmente se borra el registro auxiliar invirtiendo la computación \cite{Nielsen2024}.\\

\textbf{Modos de fallo de la búsqueda de orden}\\

El éxito de la subrutina puede verse comprometido por dos modos de fallo. El primero ocurre cuando el procedimiento de estimación de fase produce una aproximación inexacta de $s/r$; teóricamente este error $\epsilon$ se mitiga aumentando el número de qubits auxiliares, pero en hardware NISQ esta solución es limitada por la acumulación de errores por ruido y decoherencia en circuitos profundos. El segundo fallo, de naturaleza algorítmica, ocurre cuando el entero aleatorio $s$ y el orden verdadero $r$ comparten un factor común, lo que provoca que el algoritmo de fracciones continuas retorne un divisor $r'$ de $r$ en lugar del orden exacto \cite{Nielsen2024}. Para solventar esta limitación se emplean tres estrategias: (1) repetir el procedimiento $2\log(N)$ veces para garantizar con alta probabilidad la observación de un par coprimo; (2) iterar el proceso clásico reemplazando la base $x$ por $x' \equiv x^{r'} \pmod N$ para deducir el orden de forma multiplicativa; o (3) ejecutar el circuito completo dos veces de forma independiente para obtener dos estimaciones $r_1'$ y $r_2'$. La tercera es la más eficiente, ya que la probabilidad de que los numeradores $s_1'$ y $s_2'$ no compartan factores comunes está acotada inferiormente por $1/4$, y el orden verdadero se extrae de manera determinista calculando $\mathrm{mcm}(r_1', r_2')$ \cite{Nielsen2024}.\\

\section{Extracción del orden mediante fracciones continuas}
\label{app:fracciones-continuas-teoria}

Tras la medición del registro de conteo se obtiene un entero $y \in \{0, 1, \ldots, 2^t - 1\}$. La fase estimada es $\tilde{\varphi} = y/2^t$, que aproxima al valor real $\varphi = s/r$ para algún $s \in \{0, \ldots, r-1\}$. La extracción del orden $r$ procede como sigue \cite{Nielsen2024}:
\begin{enumerate}
    \item Se calcula la expansión en fracción continua de $y/2^t$:
          \begin{equation}
              \frac{y}{2^t} = a_0 + \cfrac{1}{a_1 + \cfrac{1}{a_2 + \cfrac{1}{\ddots}}},
              \label{eq:fraccion_continua_expansion}
          \end{equation}
          donde $a_0 = 0$ y los coeficientes parciales $a_i$ se obtienen mediante el algoritmo de Euclides aplicado a $y$ y $2^t$.
    \item Se calculan los convergentes sucesivos $p_k/q_k$ de la expansión, definidos recursivamente por
          \begin{equation}
              p_k = a_k p_{k-1} + p_{k-2}, \qquad q_k = a_k q_{k-1} + q_{k-2},
              \label{eq:convergentes}
          \end{equation}
          con condiciones iniciales $p_{-1} = 1$, $p_{-2} = 0$, $q_{-1} = 0$, $q_{-2} = 1$. El Teorema~\ref{thm:fracciones_continuas} (Ec.~\ref{eq:fracciones_continuas}) garantiza que si $t \geq 2L + 1$, donde $L = \lceil \log_2 N \rceil$, entonces $s/r$ aparece como uno de los convergentes \cite{Nielsen2024}.
    \item Para cada convergente con denominador $q_k < N$, se verifica computacionalmente si $x^{q_k} \equiv 1 \pmod{N}$. Si la condición se cumple, $r = q_k$ (o $r$ divide a $q_k$); en caso contrario, se prueban los convergentes subsiguientes o se repite la ejecución cuántica con un nuevo resultado de medición.
\end{enumerate}

% =========================================================================
% BLOQUE B — Apéndice A: secciones nuevas sobre DD y PT formal
% =========================================================================

\section{Formalismo del Desacoplamiento Dinámico}
\label{app:dd-formal}

Esta sección complementa la presentación del Desacoplamiento Dinámico introducida en el cuerpo principal (Sección~\ref{sec:supresion_errores}). El propósito es presentar el formalismo subyacente: la descripción del Hamiltoniano de interacción sistema--entorno, la derivación de la cancelación a primer orden mediante pulsos $\pi$, el desarrollo explícito de las secuencias $X$-$X$ y XY4, y los detalles de planificación temporal en la implementación.\\

\textbf{Hamiltoniano sistema--baño}\\

Para formalizar la situación, se parte del Hamiltoniano total del sistema (qubit) acoplado a un baño (entorno) \cite{ViolaKnillLloyd1999, Ezzell2023}:
\begin{equation}
    H = H_S \otimes \mathbb{I}_B + \mathbb{I}_S \otimes H_B + H_{SB},
    \label{eq:hamiltoniano_total_dd}
\end{equation}
donde $H_S$ describe la dinámica propia del qubit, $H_B$ la del entorno y $H_{SB}$ el acoplamiento indeseado entre ambos. Es precisamente el término $H_{SB}$ el responsable de la decoherencia: sin intervención, este término entrelaza irreversiblemente el qubit con los grados de libertad del baño, provocando una pérdida neta de información cuántica. En qubits superconductores tipo \textit{transmon}, $H_{SB}$ engloba en particular la interacción $ZZ$ residual siempre activa con qubits vecinos, el \textit{crosstalk} estático y el acoplamiento con modos del resonador \cite{Ezzell2023, AbuGhanem2025}. El Desacoplamiento Dinámico no intenta eliminar físicamente este acoplamiento, sino \textit{promediarlo a cero} mediante la aplicación rápida de pulsos de control sobre el qubit \cite{ViolaKnillLloyd1999}.\\

\textbf{Eco de espín y secuencia $X$-$X$}\\

La idea más sencilla de DD se remonta al célebre \textit{eco de espín} (\textit{spin echo}), descubierto por Hahn en 1950 en el contexto de resonancia magnética nuclear (NMR) \cite{Hahn1950}. El principio es directo: si durante un intervalo de tiempo $\tau$ el qubit acumula una fase errónea $+\phi$ debido al ruido, la aplicación de un pulso $\pi$ (una rotación de $180°$ alrededor del eje $x$, equivalente a una compuerta $X$) invierte el signo de la evolución del ruido. Si a continuación se deja pasar otro intervalo $\tau$, la fase errónea acumulada es $-\phi$, y ambas contribuciones se cancelan. El resultado neto es que el efecto del ruido se ``deshace'' a primer orden en $\tau$.\\

La extensión directa del eco de espín es la secuencia $X$-$X$ (también conocida como CPMG, una implementación concreta de la idea $X$-$X$ según \cite{Rahman2024}), que aplica dos pulsos $\pi$ alrededor del eje $x$ de forma simétrica durante el intervalo de inactividad \cite{Ezzell2023, Rahman2024}:
\begin{equation}
    X\text{-}X \equiv \frac{\tau}{2} - X - \tau - X - \frac{\tau}{2}.
    \label{eq:xx_sequence}
\end{equation}
Esta secuencia cancela a primer orden el ruido de defasamiento ($\sigma_z$), ya que $X \sigma_z X = -\sigma_z$: la primera mitad del intervalo acumula una fase errónea que la segunda mitad invierte \cite{Ezzell2023}. Sin embargo, al rotar exclusivamente alrededor del eje $x$, la secuencia $X$-$X$ no cancela errores que conmutan con el pulso, como ciertos acoplamientos $\sigma_x$ estáticos, y ofrece solo protección parcial \cite{Ezzell2023}. Pese a esta limitación, su sencillez la convierte en una opción eficaz y de bajo costo para suprimir el defasamiento dominante en qubits superconductores, y fue la secuencia empleada en las simulaciones con modelo de ruido de este trabajo.\\

\begin{figure}[htbp]
    \centering
    \includegraphics[width=0.9\textwidth]{Imagenes/ejemploDD.png}
    \caption{Ejemplo de DD de una secuencia simple X-X en un circuito sencillo \cite{ibm_error_mitigation}.}
    \label{fig:dd_ejemplo}
\end{figure}

\textbf{Cancelación universal a primer orden con XY4}\\

En la práctica, el ruido en un procesador real no se limita a un solo eje. Los qubits superconductores están sujetos a errores tanto de fase ($\sigma_z$) como de \textit{bit-flip} ($\sigma_x$) e interacciones cruzadas ($\sigma_y$). La secuencia XY4 (Ec.~\ref{eq:xy4}) es \textit{universal} de primer orden: en la idealización de pulsos instantáneos y perfectos, cancela interacciones arbitrarias del qubit con el baño hasta $\mathcal{O}(\tau^2)$ \cite{Ezzell2023, ViolaKnillLloyd1999}. Formalmente, el Hamiltoniano efectivo de error tras un ciclo completo de XY4 se reduce a \cite{Ezzell2023}:
\begin{equation}
    H_{\text{XY4}}^{\text{eff}} = \mathbb{I}_S \otimes \tilde{B} + \mathcal{O}(\tau^2),
    \label{eq:heff_xy4}
\end{equation}
donde $\tilde{B}$ es un operador que actúa exclusivamente sobre el baño. En otras palabras, a primer orden el qubit queda completamente desacoplado del entorno: la evolución residual solo afecta al baño, sin contaminar la información cuántica del sistema. Este resultado es una consecuencia directa del formalismo de simetrización grupal desarrollado en \cite{ViolaKnillLloyd1999}: en esencia, los pulsos DD aplicados periódicamente generan un promedio del Hamiltoniano de interacción sobre un grupo finito de operadores, de modo que los términos indeseados se anulan en el límite de pulsos rápidos.\\

\section{Formalismo del Pauli Twirling}
\label{app:pt-formal}

Esta sección complementa la presentación del Pauli Twirling introducida en el cuerpo principal (Sección~\ref{sec:supresion_errores}). El propósito es presentar el formalismo del \textit{twirling} sobre el grupo de Pauli, demostrar la conversión de un canal de error general en un canal estocástico de Pauli, describir el conjunto finito de pares válidos para las compuertas CNOT y ECR, y detallar la implementación \textit{circuit-by-circuit} en Qiskit Runtime.\\

\textbf{Demostración de la conversión a canal de Pauli}\\

Para formalizar la idea, considérese un canal de ruido general $\mathcal{E}$ asociado a una compuerta, cuya acción sobre el estado $\rho$ se escribe como $\mathcal{E}(\rho) = \sum_M M\rho M^{\dagger}$, donde los operadores $M$ representan los distintos mecanismos de error. El \textit{twirling exacto} sobre un conjunto de operadores de Pauli $W \subseteq \{I, X, Y, Z\}^{\otimes n}$ se define como el promedio \cite{CaiBenjamin2018}:
\begin{equation}
    \mathcal{T}_W(\mathcal{E})(\rho) = \frac{1}{|W|} \sum_{w \in W} w\, \mathcal{E}\!\left(w \rho\, w^{\dagger}\right) w^{\dagger}.
    \label{eq:twirling_exacto}
\end{equation}
En otras palabras, cada vez que se ejecuta el circuito, se conjugan los operadores de error con un par de compuertas de Pauli $w$ distinto, elegido del conjunto $W$. Al promediar sobre todos los elementos del conjunto (o sobre suficientes muestras aleatorias), las componentes fuera de la diagonal del canal de error en la base de Pauli se cancelan. El resultado es que el canal de ruido original se convierte en un \textit{canal estocástico de Pauli} \cite{CaiBenjamin2018, Wallman2016}:
\begin{equation}
    \mathcal{T}_W(\mathcal{E})(\rho) = \sum_{g \in G} p_g \, g\, \rho\, g^{\dagger},
    \label{eq:canal_pauli}
\end{equation}
donde $G = \{I, X, Y, Z\}^{\otimes n}$ es el grupo de Pauli sobre $n$ qubits y $p_g$ es la probabilidad de que ocurra el error de Pauli $g$, con $\sum_g p_g = 1$. Este tipo de canal es diagonal en la base de Pauli y se comporta de forma estadísticamente simple: los errores se acumulan linealmente con el número de compuertas, en lugar de cuadráticamente como ocurre con los errores coherentes \cite{Wallman2016}.\\

\textbf{Pares de Pauli válidos para CNOT y ECR}\\

Como se indicó en el cuerpo principal (Ec.~\ref{eq:pt_preservation}), los pares $(P_{\text{antes}}, P_{\text{después}})$ utilizados en el twirling no son arbitrarios: deben preservar la operación lógica de la compuerta original $\mathcal{U}$ salvo por una fase global irrelevante. Para cada compuerta nativa de dos qubits existe, en consecuencia, un conjunto finito de combinaciones válidas, que se enumeran en la Figura~\ref{fig:pt_cnot_ecr} para los casos CNOT y ECR. El motor de ejecución de Qiskit selecciona aleatoriamente entre estas combinaciones cada vez que genera una instancia del circuito.\\

\begin{figure}[htbp]
    \centering
    \includegraphics[width=0.85\textwidth]{Imagenes/pauli_twirling_principio.png}
    \caption{Principio algebraico del Pauli Twirling. Primera fila: una compuerta de
    Pauli $\Lambda_P$ puede absorberse en los operadores envolventes, simplificando
    el circuito. Segunda fila: una compuerta ideal $U$ envuelta entre pares de Pauli
    $P$--$P'$ equivale a $U$ con pares actualizados, preservando la operación lógica.
    Tercera fila: cuando la compuerta presenta un error coherente $\tilde{U} = U \cdot E$,
    el canal resultante tras el twirling es equivalente a $\Lambda_P \cdot U$, convirtiendo
    el error coherente en un factor de Pauli aleatorizado \cite{ibm_error_mitigation}.}
    \label{fig:pt_principio}
\end{figure}

\begin{figure}[htbp]
    \centering
    \includegraphics[width=0.9\textwidth]{Imagenes/pauli_twirling_cnot_ecr.png}
    \caption{Pauli Twirling para las compuertas nativas CNOT (filas superiores)
    y ECR (filas inferiores) en procesadores IBM Quantum. Cada conjunto de circuitos
    dentro de una fila es lógicamente equivalente, es decir, produce el mismo efecto
    ideal sobre los qubits, al satisfacer la condición
    $P_{\text{después}} \cdot U \cdot P_{\text{antes}} = e^{i\varphi} U$.
    Los pares de Pauli mostrados ($X$, $Y$, $Z$ en los qubits de control y objetivo)
    constituyen el conjunto finito de combinaciones válidas del que Qiskit Runtime
    selecciona aleatoriamente en cada instancia del circuito \cite{ibm_error_mitigation}.}
    \label{fig:pt_cnot_ecr}
\end{figure}

% =========================================================================
% BLOQUE B — Apéndice B "Arquitectura del software e implementación"
% =========================================================================

\chapter{Arquitectura del software e implementación}
\label{app:arquitectura-software-cap}

Este apéndice documenta los aspectos de implementación del software empleado en el trabajo. Su contenido complementa el Capítulo~\ref{chap:metodologia} y permite la reproducibilidad técnica del estudio sin sobrecargar el discurso metodológico del cuerpo principal. Las secciones siguen el orden lógico del flujo de desarrollo: migración técnica respecto al repositorio base, estructura modular del proyecto, configuración detallada del transpilador y del algoritmo SABRE, configuración de las técnicas de supresión, \textit{scripts} de ejecución por fase, módulos de post-procesamiento y, finalmente, versiones del entorno de software.

% =====================================================================
% B.1 MIGRACIÓN TÉCNICA
% =====================================================================
\section{Migración técnica del repositorio base a Qiskit~2.x}
\label{app:migracion-qiskit}

El repositorio proporcionado por \cite{villarreal} aportó la arquitectura general de clases (\texttt{RegisterQC}, \texttt{Shor}, \texttt{Factorizer}, \texttt{Transpiler}) y el flujo de ejecución basado en archivos de configuración. La presente sección detalla las modificaciones técnicas concretas que fueron necesarias para adaptar dicha base al ecosistema actual de IBM Quantum (Qiskit SDK $\geq$ 2.x, primitivas V2 y \textit{Fake Backends} migrados al paquete \path{qiskit_ibm_runtime.fake_provider}). La Tabla~\ref{tab:delimitacion-aportes} sintetiza el conjunto de sustituciones aplicadas.

\newpage

\begin{table}[H]
    \centering
    \caption{Delimitación de aportes y resumen de la migración a Qiskit~2.x: repositorio base frente a contribución propia.\\}
    \label{tab:delimitacion-aportes}
    \footnotesize
    \begin{tabular}{p{0.22\textwidth} p{0.34\textwidth} p{0.34\textwidth}}
        \toprule
        \textbf{Elemento} & \textbf{Aporte del repositorio base \cite{villarreal}} & \textbf{Contribución de este trabajo} \\
        \midrule
        Arquitectura y SDK
            & Arquitectura general de clases y flujo de ejecución gobernado por archivos de configuración, basado en versiones previas del SDK de Qiskit.
            & Migración integral al ecosistema Qiskit~2.x, conservando la arquitectura general pero reescribiendo los módulos incompatibles con las primitivas V2 y los nuevos proveedores. \\
        \midrule
        Acceso a \textit{Fake Backends}
            & \texttt{qiskit.providers.fake\_pro\-vider} (deprecado).
            & \texttt{qiskit\_ibm\_runtime.fake\_pro\-vider}, con instanciación directa \texttt{FakeTorino()} correspondiente al procesador Heron~r1. \\
        \midrule
        Simulador con ruido
            & Composición manual del modelo de ruido a partir del \textit{provider} antiguo.
            & \texttt{AerSimulator.from\_backend\-(FakeTorino())}, que carga automáticamente el \textit{snapshot} de calibración. \\
        \midrule
        Parámetro \texttt{approximation\_de\-gree}
            & Uso ambiguo: aplicado sin distinguir entre su acepción de truncamiento entero (AQFT lógica) y su acepción de umbral flotante en la síntesis del transpilador.
            & Identificación y corrección de la ambigüedad: \texttt{approximation\_degree}$=0{,}7$ se fija exclusivamente en la etapa de transpilación (flotante en $[0,1]$); la \texttt{QFT} lógica se instancia sin truncamiento adicional. \\
        \midrule
        Técnicas de supresión (DD y \textit{Pauli Twirling})
            & Configuración por defecto: secuencia DD heredada sin selección explícita y sin activación de \textit{Pauli Twirling}.
            & Configuración manual mediante \texttt{SamplerOptions}, conforme a las exigencias propias de \texttt{SamplerV2}. \\
        \midrule
        Conversión \textit{bitstring}-a-entero
            & Conversión implícita y dependencia del \textit{Endianness} sin documentación ni verificación.
            & Conversión explícita con la convención \textit{little-endian} de Qiskit, verificada empíricamente en la Fase~1 contra los picos teóricos $y_s = \lfloor s\,2^t/r \rceil$ obtenidos por simulación ideal. \\
        \midrule
        Protocolo experimental
            & Ausencia de un protocolo estructurado por fases.
            & Diseño e implementación de un protocolo en tres fases (simulación ideal, simulación con ruido y hardware real), con \textit{scripts} autocontenidos y configuración independiente para cada una. \\
        \midrule
        Post-procesamiento por fracciones continuas
            & Utilidades dispersas en el repositorio y post-procesamiento elemental.
            & Módulo propio diferenciado por fase (\path{fracciones_continuas_*.py}) con interfaz común, verificación del orden $r$, extracción de factores no triviales e integración con la verificación del \textit{Endianness}. \\
        \bottomrule
    \end{tabular}
\end{table}

\newpage

% =====================================================================
% B.2 ESTRUCTURA MODULAR
% =====================================================================
\section{Estructura modular del proyecto}
\label{app:arquitectura-software}

La arquitectura del proyecto se estructura en cinco capas funcionales que reflejan la separación lógica entre las responsabilidades del algoritmo, los servicios de ejecución y las utilidades de configuración. Todo el código se encuentra en el repositorio público del trabajo \cite{compuertas_CNOT}; la organización por carpetas se ilustra en el árbol siguiente.

\begin{verbatim}
quantum-shors-algorithm/
+-- main.py                # Punto de entrada
+-- common/
|   +-- settings.py/settings.yaml    # Configuracion global
|   \-- backend.py                   # Factory de backends
+-- algorithm/
|   +-- circuit.py            # Clase RegisterQC
|   \-- phase_estimation.py   # Clase Shor (orquestador)
+-- services/
|   +-- arithmetics.py       # Factorizer(envoltura clasica)
|   +-- sampling.py          # AerSampler / IBMRuntimeSampler
|   \-- transpilation.py     # Transpiler(PassManager+SABRE)
+-- utilities/
|   +-- optimal_a.py           # Seleccion de bases validas
|   +-- job_results.py       # Recuperacion asincrona de jobs
|   \-- arguments_parser.py    # Validacion de argumentos
+-- simulacion_ideal/          # Fase 1: scripts y resultados
+-- simulacion_con_ruido/      # Fase 2: scripts y resultados
\-- ejecucion_hardware/        # Fase 3: scripts y resultados
\end{verbatim}

\vspace{0.5cm}

\textit{Capa de entrada y configuración} (\texttt{main.py}, \texttt{common/settings.py}, \texttt{utiliti\-es/arguments\_parser.py}). El punto de entrada del programa carga la configuración centralizada desde un archivo YAML y procesa los argumentos de la línea de comandos o, alternativamente, los parámetros de un archivo \path{.ini} de configuración masiva. Este esquema sigue la filosofía de los \textit{Qiskit Patterns}, que promueven un flujo cíclico: mapear el problema, optimizar el circuito, ejecutar y post-procesar.\\

\textit{Capa algorítmica} (\texttt{algorithm/circuit.py}, \texttt{algorithm/phase\_estimati\-on.py}). Contiene la lógica de construcción del circuito cuántico (clase \texttt{RegisterQC}) y la orquestación de la estimación de fase cuántica (clase \texttt{Shor}). La clase \texttt{Shor} coordina todo el flujo experimental: configura el \textit{backend}, invoca al transpilador, activa las opciones de supresión de errores en la primitiva \texttt{SamplerV2}, envía el trabajo (\textit{job}) y, en el caso de simulaciones locales, procesa los resultados sincrónicamente.\\

\textit{Capa de servicios} (\texttt{services/arithmetics.py}, \texttt{services/sampling.py}, \texttt{services/transpilation.py}). Implementa las operaciones intermedias. La clase \texttt{Factorizer} encapsula la envoltura clásica del algoritmo de Shor (verificaciones de paridad, potencias primas y cálculo del $\gcd$), delegando la subrutina cuántica a la clase \texttt{Shor}. La clase \texttt{Transpiler} envuelve el \texttt{PassManager} de Qiskit y gestiona la inyección de parámetros SABRE personalizados. Las clases \texttt{AerSampler} e \texttt{IBMRuntimeSampler} abstraen la ejecución en simulador local y en hardware real, respectivamente, unificando la interfaz de muestreo.\\

\textit{Capa de \textit{backends}} (\texttt{common/backend.py}). Implementa el patrón de diseño \textit{Factory} para instanciar tres tipos de \textit{backend} de manera intercambiable, alineados con las recomendaciones de Qiskit~2.x: \texttt{AerSimulator} sin modelo de ruido (simulación ideal); \texttt{AerSimulator.from\_backend(FakeTorino())} para simulación con el modelo de ruido calibrado del procesador Heron~r1 (que reemplaza al antiguo \texttt{FakeProvider}); y \texttt{QiskitRuntimeService}~$+$~\texttt{SamplerV2} en modo \texttt{ibm\_quantum} para ejecución en hardware real.\\

\textit{Capa de utilidades.} Proporciona herramientas de pre-procesamiento clásico (selección de bases óptimas $a$ a través de \path{get_optimal_a_for_valid_periods}), post-procesamiento (fracciones continuas, extracción de factores, graficación de resultados) y validación de argumentos.\\

\textit{Conexión con IBM Quantum.} Se establece mediante el servicio \texttt{QiskitRuntimeSer\-vice}, que utiliza las credenciales almacenadas localmente bajo un nombre de cuenta configurable (por ejemplo, \path{ibm_quantum}). La ejecución se realiza a través de la primitiva \texttt{SamplerV2}. A diferencia de \texttt{EstimatorV2}, \texttt{SamplerV2} no expone un parámetro \texttt{resilience\_level}: todas las técnicas de supresión deben configurarse explícitamente a través de \texttt{SamplerOptions} (cf. Sección~\ref{app:supresion-config}).

% =====================================================================
% B.3 CONFIGURACIÓN DETALLADA DEL TRANSPILADOR Y SABRE AVANZADO
% =====================================================================
\section{Configuración detallada del transpilador y SABRE avanzado}
\label{app:sabre-avanzado}

La Tabla~\ref{tab:transpilador_config} consolida los parámetros del transpilador empleados en este trabajo, con los valores nominales que se aplican en todas las fases salvo indicación explícita. Estos parámetros se inyectan directamente al método \texttt{transpile} de Qiskit a través de la clase envolvente \texttt{Transpiler}.

\newpage

\begin{table}[H]
    \centering
    \caption{Parámetros del transpilador empleados en el trabajo.}
    \label{tab:transpilador_config}
    \footnotesize
    \begin{tabular}{p{0.32\textwidth} p{0.30\textwidth} p{0.30\textwidth}}
        \toprule
        \textbf{Parámetro} & \textbf{Valor} & \textbf{Función} \\
        \midrule
        \texttt{optimization\_level}
            & $\{0, 1, 2, 3\}$ (barrido)
            & Agresividad de las transformaciones aplicadas al circuito. \\
        \midrule
        \texttt{basis\_gates}
            & \texttt{cz}, \texttt{rzz}, \texttt{id}, \texttt{rx}, \texttt{rz}, \texttt{sx}, \texttt{x}
            & Conjunto nativo de Heron~r1; \texttt{cz} es la única compuerta de dos qubits. \\
        \midrule
        \texttt{layout\_method}
            & \texttt{sabre}
            & Asignación inicial de qubits lógicos a qubits físicos. \\
        \midrule
        \texttt{routing\_method}
            & \texttt{sabre}
            & Inserción de SWAPs durante el enrutamiento. \\
        \midrule
        \texttt{approximation\_degree}
            & $0{,}7$
            & Umbral de fidelidad aceptado durante la síntesis de bloques de dos qubits. \\
        \midrule
        \texttt{seed\_transpiler}
            & $457$
            & Semilla fija para la reproducibilidad de las heurísticas de SABRE. \\
        \midrule
        \texttt{scheduling\_method}
            & \texttt{alap} (solo Fase~3)
            & Planifica las instrucciones \textit{as late as possible} para concentrar los \texttt{Delay} al inicio del circuito. \\
        \bottomrule
    \end{tabular}
\end{table}

Adicionalmente, la clase \texttt{Transpiler} permite inyectar parámetros avanzados del algoritmo SABRE directamente en el \textit{pipeline}. El método \path{_set_transpile_optimization_params} define tres hiperparámetros que controlan la exploración heurística del espacio de soluciones, recogidos en la Tabla~\ref{tab:sabre_avanzado}.

\begin{table}[H]
    \centering
    \caption{Hiperparámetros avanzados de SABRE.}
    \label{tab:sabre_avanzado}
    \footnotesize
    \begin{tabular}{p{0.25\textwidth} p{0.15\textwidth} p{0.50\textwidth}}
        \toprule
        \textbf{Hiperparámetro} & \textbf{Valor} & \textbf{Función} \\
        \midrule
        \texttt{max\_iterations}
            & $10$
            & Número de iteraciones bidireccionales del algoritmo SABRE. Cada iteración recorre el circuito en ambas direcciones, refinando el mapeo lógico--físico. \\
        \midrule
        \texttt{layout\_trials}
            & $200$
            & Número de semillas aleatorias para la búsqueda del \textit{layout} inicial. A mayor número de ensayos, mayor probabilidad de encontrar un mapeo con menor distancia topológica promedio. \\
        \midrule
        \texttt{swap\_trials}
            & $200$
            & Número de candidatos SWAP evaluados en cada paso de enrutamiento. \\
        \bottomrule
    \end{tabular}
\end{table}

Cuando esta optimización está habilitada, \texttt{Transpiler} construye un objeto \texttt{SabreLayout} con el mapa de acoplamiento (\texttt{CouplingMap}) extraído del \textit{backend} y lo inyecta directamente en la etapa de \textit{layout} del \texttt{PassManager}, reemplazando la implementación por defecto. Este procedimiento se implementa en el método \path{get_sabre_layout}, que maneja explícitamente los casos en que el \textit{backend} no expone un mapa de acoplamiento (como en el simulador ideal).\\

Tras la transpilación, el método \path{get_isa_statistics} extrae las métricas cuantitativas descritas en el cuerpo principal (conteo por tipo de compuerta, número total y profundidad de dos qubits $D_{2Q}$). Como salvaguarda de seguridad en el envío a hardware, si el circuito transpilado excede las $10^5$ compuertas de dos qubits se excluye automáticamente del envío para no desperdiciar recursos en una ejecución cuya señal sería inevitablemente colapsada por el ruido acumulado.

% =====================================================================
% B.4 CONFIGURACIÓN DETALLADA DE LAS TÉCNICAS DE SUPRESIÓN
% =====================================================================
\section{Configuración detallada de las técnicas de supresión}
\label{app:supresion-config}

Las técnicas de supresión de errores se configuran a través de los objetos \texttt{DynamicalDecou\-plingOptions} y \texttt{TwirlingOptions}, que se ensamblan en un único \texttt{SamplerOptions} y se pasan a la primitiva \texttt{SamplerV2} en el momento de la ejecución. Las Tablas~\ref{tab:dd_config} y~\ref{tab:pt_config} consolidan los parámetros adoptados.

\begin{table}[H]
    \centering
    \caption{Parámetros de \texttt{DynamicalDecouplingOptions} empleados.}
    \label{tab:dd_config}
    \footnotesize
    \begin{tabular}{p{0.30\textwidth} p{0.20\textwidth} p{0.45\textwidth}}
        \toprule
        \textbf{Parámetro} & \textbf{Valor} & \textbf{Función} \\
        \midrule
        \texttt{enable}
            & \texttt{True} / \texttt{False}
            & Activa o desactiva la inserción automática de secuencias DD. \\
        \midrule
        \texttt{sequence\_type}
            & \texttt{XY4}
            & Secuencia de pulsos aplicada en los intervalos ociosos. \texttt{XY4} cancela ruido de fase y de bit a primer orden, frente a \texttt{XpXm} heredado del repositorio base. \\
        \midrule
        \texttt{scheduling\_method}
            & \texttt{alap}
            & Planifica las instrucciones lo más tarde posible, concentrando los \texttt{Delay} al inicio del circuito y maximizando el espacio donde insertar pulsos DD. \\
        \bottomrule
    \end{tabular}
\end{table}

\begin{table}[H]
    \centering
    \caption{Parámetros de \texttt{TwirlingOptions} empleados.}
    \label{tab:pt_config}
    \footnotesize
    \begin{tabular}{p{0.30\textwidth} p{0.20\textwidth} p{0.45\textwidth}}
        \toprule
        \textbf{Parámetro} & \textbf{Valor} & \textbf{Función} \\
        \midrule
        \texttt{enable\_gates}
            & \texttt{True} / \texttt{False}
            & Inserta pares de operadores de Pauli aleatorios alrededor de cada compuerta de dos qubits. \\
        \midrule
        \texttt{enable\_measure}
            & \texttt{False} (por defecto)
            & Aplica \textit{twirling} también a las operaciones de medición; opcional. \\
        \midrule
        \texttt{strategy}
            & \texttt{active}
            & El \textit{twirling} se aplica únicamente a los qubits involucrados en cada compuerta. \\
        \midrule
        \texttt{num\_randomizations}
            & \texttt{auto}
            & El número de instancias aleatorias lo determina automáticamente el \textit{runtime} en función del número total de \textit{shots}. \\
        \bottomrule
    \end{tabular}
\end{table}

Cuando ambas opciones están habilitadas simultáneamente (configuración DD$+$PT), el \textit{runtime} de IBM Quantum aplica primero la transpilación con la planificación \texttt{alap} y la inserción de secuencias \texttt{XY4} en los intervalos ociosos; sobre el circuito así programado, genera múltiples copias con distintas selecciones aleatorias de pares de Pauli para las compuertas de dos qubits, y agrega los conteos de todas las copias en la distribución final reportada al usuario.

% =====================================================================
% B.5 SCRIPTS DE EJECUCIÓN POR FASE
% =====================================================================
\section{Scripts de ejecución por fase}
\label{app:scripts-fases}

Los \textit{scripts} que ejecutan de manera autocontenida las tres fases del estudio constituyen el núcleo experimental de este trabajo y no pertenecen al repositorio original. Cada \textit{script} carga su propio archivo \texttt{.ini} de configuración y produce reportes automáticos en formato JSON y Markdown junto con las figuras correspondientes. La Tabla~\ref{tab:scripts_fases} resume los \textit{scripts} principales por fase.

\begin{table}[H]
    \centering
    \caption{\textit{Scripts} de ejecución por fase del protocolo experimental.}
    \label{tab:scripts_fases}
    \footnotesize
    \begin{tabular}{p{0.13\textwidth} p{0.40\textwidth} p{0.42\textwidth}}
        \toprule
        \textbf{Fase} & \textbf{Script} & \textbf{Función} \\
        \midrule
        Fase 1
            & \path{ejecutar_shor_ideal_puro.py}
            & Simulación ideal sobre \texttt{AerSimulator} sin modelo de ruido. \\
        \midrule
        Fase 1
            & \path{ejecutar_shor_faketorino_ideal.py}
            & Variante complementaria sobre \texttt{AerSimulator.from\_backend(Fake\-Torino())} con el modelo de ruido desactivado (\texttt{set\_options(noise\_model=None)}); aísla el \textit{overhead} topológico. \\
        \midrule
        Fase 2
            & \path{ejecutar_shor_faketorino_ruido.py}
            & Línea base ruidosa sin mitigación. Barre \texttt{opt}~$\in \{0, 3\}$ para $a \in \{4, 7\}$. \\
        \midrule
        Fase 2
            & \path{calcular_fidelidad_ruido.py}
            & Cálculo detallado de $\mathcal{F}_H$, $D_H$ y $D_{KL}$ frente a la distribución teórica. \\
        \midrule
        Fase 2
            & \path{comparar_mitigacion.py}
            & Diseño factorial $2 \times 2$ de las cuatro variantes (sin supresión, DD, PT, DD$+$PT). \\
        \midrule
        Fase 2
            & \path{fracciones_continuas_ruido.py}
            & Post-procesamiento detallado bajo DD$+$PT con clasificación señal/ruido por medición. \\
        \midrule
        Fase 3
            & \path{ejecutar_estudio_por_lotes_hardware.py}
            & Envío y registro masivo de \textit{jobs} sobre \texttt{ibm\_torino}. Agrupa los estudios A2, A3 y A6/A7. \\
        \midrule
        Fase 3
            & \path{ejecutar_analisis_transpilacion_hardware.py}
            & Análisis comparativo de la transpilación al \textit{backend} real previo al envío. \\
        \midrule
        Fase 3
            & \path{analizar_caso_a4_hardware.py}
            & Estudio detallado del caso $a = 4$ (mayor número de configuraciones por estudio A2 y A6/A7). \\
        \bottomrule
    \end{tabular}
\end{table}

La Tabla~\ref{tab:jobs_hardware_apendice} desglosa los 13~\textit{jobs} ejecutados en \texttt{ibm\_torino}, organizados en los tres estudios complementarios mencionados en el cuerpo principal. Cada \textit{job} se identificó con su correspondiente \texttt{job\_id} para permitir la recuperación asíncrona de los resultados días después del envío.

\begin{table}[H]
    \centering
    \caption{Organización de los 13 \textit{jobs} ejecutados en \texttt{ibm\_torino}.}
    \label{tab:jobs_hardware_apendice}
    \footnotesize
    \begin{tabular}{p{0.18\textwidth} p{0.30\textwidth} p{0.45\textwidth}}
        \toprule
        \textbf{Estudio} & \textbf{Configuración} & \textbf{Propósito} \\
        \midrule
        V1 (inicial)
            & PT activado, \texttt{opt}$=3$, $a \in \{4, 7, 14\}$
            & Calibración inicial del flujo de envío y verificación cualitativa de la señal. \\
        \midrule
        A2 (\texttt{opt} barrido)
            & $a = 4$ fijo, \texttt{opt}~$\in \{0, 1, 2\}$ con PT activado
            & Cuantificar el impacto del nivel de optimización del transpilador en la señal en hardware. \\
        \midrule
        A3 (base $a$)
            & \texttt{opt}$=3$, PT activado, $a \in \{2, 4, 7, 8, 11, 13, 14\}$
            & Caracterizar la respuesta del procesador en función de la base. Las bases $a = 11$, $a = 14$ se incluyen con propósito diagnóstico (cf. Sección~\ref{sec:protocolo_fases}). \\
        \midrule
        A6/A7 (supresión)
            & $a = 4$, \texttt{opt}$=3$; sin supresión, solo DD (XY4), DD$+$PT
            & Evaluar el impacto individual y combinado de las técnicas de supresión sobre $\mathcal{F}_H$ y $PST$. \\
        \bottomrule
    \end{tabular}
\end{table}

Los identificadores \texttt{job\_id} concretos asignados por el servicio de IBM Quantum y las métricas obtenidas para cada \textit{job} se reportan en la Tabla~\ref{tab:jobs_hardware} del Capítulo~\ref{chap:resultados}, donde se discute el análisis de los resultados experimentales.

% =====================================================================
% B.6 MÓDULOS DE POST-PROCESAMIENTO
% =====================================================================
\section{Módulos de post-procesamiento clásico}
\label{app:postproc-modulos}

El post-procesamiento se implementa en tres módulos paralelos, uno por cada fase experimental: \path{fracciones_continuas_simulacion_ideal.py}, \path{fracciones_continuas_simulacion_ruido.py} y \path{fracciones_continuas_hardware.py}. Los tres comparten la misma estructura algorítmica documentada en la Sección~\ref{sec:postprocesamiento} del cuerpo principal y se diferencian únicamente en los detalles de carga de datos: el módulo de simulación lee los conteos del objeto devuelto por \texttt{AerSampler}, mientras que el módulo de hardware recupera los \textit{jobs} a partir del \path{job_id} registrado y accede a la distribución de conteos mediante \texttt{job.result()[0].data.output.get\_counts()}.\\

La interfaz común de los tres módulos consiste en una función \texttt{procesar\_distribucion\-(counts, a, N, t)} que devuelve un diccionario estructurado con los siguientes campos: la lista ordenada de los $k$ \textit{bitstrings} más frecuentes, la fase $\tilde{\varphi}$ asociada a cada uno, los convergentes $p_k/q_k$ devueltos por la rutina de fracciones continuas, los candidatos a $r$ que satisfacen la verificación clásica, los factores no triviales obtenidos mediante $\gcd(a^{q_k/2} \pm 1, N)$ y, finalmente, las métricas $\mathcal{F}_H$ y $PST$ calculadas conforme a la Sección~\ref{sec:postprocesamiento}.\\

La aplicación del algoritmo de fracciones continuas se realiza con \texttt{Fraction(phase).li\-mit\_denominator(N)} de la biblioteca estándar de Python, que devuelve los convergentes con denominador menor o igual a $N$. La conversión \textit{bitstring}-a-entero se aplica antes del paso de fracciones continuas y se efectúa de acuerdo con la convención \textit{little-endian} de Qiskit, verificada empíricamente en la Fase~1 contra los picos teóricos $y_s = \lfloor s \cdot 2^t / r \rceil$.

% =====================================================================
% B.7 ENTORNO Y VERSIONES
% =====================================================================
\section{Configuración del entorno y versiones de software}

La Tabla~\ref{tab:entorno_software} resume las versiones específicas de cada paquete utilizado en los experimentos. La fijación explícita de estas versiones es necesaria para la reproducibilidad: el ecosistema de Qiskit ha experimentado, en los últimos dos años, sustituciones de API significativas entre versiones consecutivas, especialmente en la transición a las primitivas V2 y a la nueva ubicación de los \textit{Fake Backends}.

\begin{table}[H]
    \centering
    \caption{Entorno de software utilizado en los experimentos.\\}
    \label{tab:entorno_software}
    \begin{tabular}{ll}
        \toprule
        \textbf{Paquete / Herramienta}           & \textbf{Versión} \\
        \midrule
        Python                                   & 3.11   \\
        \texttt{qiskit}                          & 2.3.0  \\
        \texttt{qiskit-aer}                      & 0.16.0 \\
        \texttt{qiskit-ibm-runtime}              & 0.34.0 \\
        \texttt{qiskit-ibm-transpiler}           & 0.11.1 \\
        \texttt{numpy}                           & 2.2.2  \\
        \texttt{scipy}                           & 1.15.1 \\
        \texttt{matplotlib}                      & 3.10.0 \\
        \texttt{pandas}                          & 2.2.3  \\
        \texttt{seaborn}                         & 0.13.2 \\
        \bottomrule
    \end{tabular}
\end{table}

% =========================================================================
% BLOQUE B — Apéndice C "Tablas de validación complementaria"
% (incluir Tabla 6.1, Tabla 6.2 y Figura 6.1)
% =========================================================================

\chapter{Tablas complementarias}
\label{app:caracterizacion-bloques}

Este apéndice consolida las tablas y figuras de validación numérica del capítulo 6. Su contenido respalda las afirmaciones de los capítulos anteriores.

\section{Verificación de la Transformada Cuántica de Fourier}

La QFT sobre $n = 9$ qubits de control se implementó conforme a la descomposición estándar en compuertas Hadamard y rotaciones de fase controladas $CR_k$ (Ec.~\ref{eq:rotacion_fase}), requiriendo teóricamente $n = 9$ compuertas Hadamard, $\binom{9}{2} = 36$ rotaciones $CR_k$, y $\lfloor 9/2 \rfloor = 4$ compuertas SWAP para la inversión del orden de los qubits de salida. La corrección del operador se verificó numéricamente calculando la norma de Frobenius del error residual $\|U^{\dagger}U - I\|_F$, tanto para la QFT directa como para su inversa. Adicionalmente, se verificó la relación de inversión $\mathrm{QFT} \cdot \mathrm{QFT}^{-1} = I$ obteniendo $\|\mathrm{QFT} \cdot \mathrm{QFT}^{-1} - I\|_F = 3.82 \times 10^{-14}$, lo cual confirma que ambos operadores son mutuamente inversos con precisión numérica de máquina. La discrepancia entre la profundidad descompuesta (Depth Total $= 18$) y la transpilada (Depth Total $= 61$) refleja el costo de traducir las rotaciones $CR_k$ abstractas a la base nativa $\{\mathrm{RZ}, \mathrm{SX}, \mathrm{X}, \mathrm{ECR}\}$, un factor que se amplifica significativamente al considerar la topología restringida del hardware.

\begin{table}[H]
    \centering
    \caption{Métricas de profundidad y verificación unitaria de la QFT para $n = 9$ qubits. La columna $\|U^{\dagger}U - I\|_F$ cuantifica la desviación de la unitariedad; valores del orden de $10^{-14}$ confirman precisión numérica de punto flotante.}
    \label{tab:qft_metricas}
    \begin{tabular}{lcccr}
        \hline
        \textbf{Variante}         & \textbf{Depth Total} & \textbf{Depth 2Q} & \textbf{Gates Totales} & $\boldsymbol{\|U^{\dagger}U - I\|_F}$ \\
        \hline
        QFT directa (decompose)   & 18                   & 16                & 49                     & $3.84 \times 10^{-14}$                \\
        QFT$^{-1}$ (decompose)    & 18                   & 16                & 49                     & $3.79 \times 10^{-14}$                \\
        QFT directa (transpilada) & 61                   & 30                & 173                    & $3.84 \times 10^{-14}$                \\
        QFT$^{-1}$ (transpilada)  & 61                   & 30                & 189                    & $3.79 \times 10^{-14}$                \\
        \hline
    \end{tabular}
\end{table}

\section{Bloques de exponenciación modular controlada}

El operador unitario de multiplicación modular $U_a |y\rangle = |ay \bmod N\rangle$ se construyó como una matriz de permutación de dimensión $2^L \times 2^L = 16 \times 16$ para cada base $a$. Como se demostró en la Sección~\ref{sec:construccion_circuitos}, este operador satisface la propiedad de periodicidad $(U_a)^r = I$, donde $r = \mathrm{ord}_N(a)$ es el orden multiplicativo. Esta propiedad se verificó numéricamente para todas las bases, obteniendo $\|(U_a)^r - I\|_F = 0.00$ en todos los casos. La Tabla~\ref{tab:exp_mod_bloques} resume la complejidad de los bloques individuales $C$-$U_{a^{2^k} \bmod N}$ para cada base.

\begin{table}[H]
    \centering
    \caption{Profundidad de dos qubits (Depth 2Q) y conteo de compuertas de dos qubits (Gates 2Q) para el circuito completo de exponenciación modular controlada, por base. La profundidad total es la suma de los bloques no triviales $C$-$U_{a^{2^k}}$ con $a^{2^k} \bmod N \neq 1$.\\}
    \label{tab:exp_mod_bloques}
    \begin{tabular}{ccccc}
        \hline
        \textbf{Base $a$} & $\boldsymbol{r = \mathrm{ord}_{15}(a)}$ & \textbf{Bloques no triviales}        & \textbf{Depth 2Q total} & \textbf{Gates 2Q total} \\
        \hline
        4                 & 2                                       & $k = 0$ ($b = 4$)                    & 233                     & 236                     \\
        7                 & 4                                       & $k = 0$ ($b = 7$), $k = 1$ ($b = 4$) & 466                     & 472                     \\
        11                & 2                                       & $k = 0$ ($b = 11$)                   & 230                     & 233                     \\
        14                & 2                                       & $k = 0$ ($b = 14$)                   & 231                     & 234                     \\
        \hline
    \end{tabular}
\end{table}

\newpage

\begin{figure}[H]
    \centering
    \includegraphics[width=\textwidth]{simulacion_ideal/caracterizacion_bloques_N15.png}
    \caption{Caracterización de los bloques fundamentales del algoritmo de Shor para $N = 15$. Se muestra la profundidad y el conteo de compuertas para la QFT ($n = 9$ qubits) y la exponenciación modular controlada $C$-$U_{a^{2^k}}$ para cada base $a \in \{4, 7, 11, 14\}$. La verificación numérica confirma la unitariedad de todos los operadores ($\|U^{\dagger}U - I\|_F < 10^{-10}$) y la periodicidad $(U_a)^r = I$, validando la corrección de la implementación antes de la ejecución del circuito completo.}
    \label{fig:caracterizacion_bloques}
\end{figure}

\newpage

\section{Tablas de resultados en simulación ideal}

\begin{table}[htpb]
    \centering
    \caption{Métricas del circuito completo \texttt{RegisterQC} (13 Qubits) para $N = 15$ antes (pre-transpilación) y después (post-transpilación a la base $CX + U3$, sin restricción topológica). La \textit{Depth Total (pre)} corresponde a las 12 capas secuenciales del circuito de alto nivel: una compuerta $X$ de inicialización, nueve pares $H-C-U_{a^{2k}}$ aplicados en secuencia sobre cada qubit de control, una $\text{QFT}^{-1}$ y una medición final. La post-transpilación revela el costo real en compuertas de dos qubits. Entiéndase (D.) como Depth (profundidad).}
    \label{tab:transpilacion_pre}
    \begin{tabular}{cccccc}
        \toprule
        \textbf{Base} & \textbf{D.Total (pre)} & \textbf{D.2Q (pre)} & \textbf{D.Total (post)} & \textbf{D.2Q (post)} & \textbf{CX (post)} \\
        \midrule
        4  & 12 & 0 & 462 & 233 & 308 \\
        7  & 12 & 0 & 919 & 463 & 541 \\
        11 & 12 & 0 & 461 & 230 & 305 \\
        14 & 12 & 0 & 463 & 231 & 306 \\
        \bottomrule
    \end{tabular}
\end{table}

\newpage

\begin{table}[htpb]
    \centering
    \caption{Métricas del circuito \texttt{RegisterQC} ejecutado en simulación ideal pura (\textit{all-to-all}) para $N = 15$, $t = 9$ y $4096 \text{ shots}$ (transpilado con \texttt{AerSimulator}), evaluadas bajo los cuatro niveles de optimización del transpilador de Qiskit. Se reportan la profundidad total ($D$), la profundidad restringida a compuertas de dos qubits ($D_{2Q}$), el conteo de compuertas de dos qubits ($G_{2Q}$), el conteo total de compuertas ($G_{\text{tot}}$) y la probabilidad de señal ($PST$). El estado $\Delta$ indica que, pese a la correcta identificación del orden $r$, los factores clásicos son triviales.}
    \label{tab:transpilacion_ideal}
    \begin{tabular}{cccccccc}
        \toprule
        $a$ & \textbf{Opt} & \textbf{Depth} & \textbf{Depth$_{2Q}$} & \textbf{Gates$_{2Q}$} & \textbf{Gates$_{\text{tot}}$} & \textbf{PST (\%)} & \textbf{Estado} \\
        \midrule
        4  & 0 & 463 & 233 & 236 & 734  & 100.0 & \checkmark \\
        4  & 1 & 462 & 233 & 236 & 702  & 100.0 & \checkmark \\
        4  & 2 & 462 & 233 & 236 & 687  & 100.0 & \checkmark \\
        4  & 3 & 462 & 233 & 236 & 687  & 100.0 & \checkmark \\
        7  & 0 & 927 & 466 & 472 & 1403 & 100.0 & \checkmark \\
        7  & 1 & 925 & 466 & 472 & 1336 & 100.0 & \checkmark \\
        7  & 2 & 914 & 461 & 466 & 1296 & 100.0 & \checkmark \\
        7  & 3 & 914 & 461 & 466 & 1296 & 100.0 & \checkmark \\
        11 & 0 & 464 & 232 & 235 & 732  & 100.0 & \checkmark \\
        11 & 1 & 463 & 232 & 235 & 700  & 100.0 & \checkmark \\
        11 & 2 & 461 & 230 & 233 & 687  & 100.0 & \checkmark \\
        11 & 3 & 461 & 230 & 233 & 687  & 100.0 & \checkmark \\
        14 & 0 & 466 & 233 & 236 & 734  & 100.0 & $\Delta$ \\
        14 & 1 & 465 & 233 & 236 & 705  & 100.0 & $\Delta$ \\
        14 & 2 & 463 & 231 & 234 & 687  & 100.0 & $\Delta$ \\
        14 & 3 & 463 & 231 & 234 & 687  & 100.0 & $\Delta$ \\
        \bottomrule
    \end{tabular}
\end{table}

\newpage

\begin{table}[htpb]
    \centering
    \caption{Métricas del circuito \texttt{RegisterQC} transpilado a la topología Heavy-Hex de \texttt{FakeTorino} (Heron r1) y ejecutado en simulación ideal (sin modelo de decoherencia ni errores de compuerta) para $N = 15$, $t = 9$ y $4096 \text{ shots}$. Se reportan la profundidad y el conteo de compuertas de dos qubits ($D_{2Q}$, $G_{2Q}$), la probabilidad de señal ($PST$), la probabilidad de ruido (masa fuera de los picos teóricos del QPE) y los factores extraídos. El estado $\Delta$ marca los casos con factores triviales por $a \equiv -1 \pmod N$.}
    \label{tab:transpilacion_faketorino_trivial}
    \begin{tabular}{cccccccc}
        \toprule
        $a$ & \textbf{Opt} & \textbf{Depth$_{2Q}$} & \textbf{Gates$_{2Q}$} & \textbf{PST (\%)} & \textbf{Ruido (\%)} & \textbf{Factores} & \textbf{Estado} \\
        \midrule
        4  & 0 & 570  & 923  & 100.0 & 0.0 & 3, 5 & \checkmark \\
        4  & 1 & 526  & 779  & 100.0 & 0.0 & 3, 5 & \checkmark \\
        4  & 2 & 433  & 584  & 100.0 & 0.0 & 3, 5 & \checkmark \\
        4  & 3 & 439  & 600  & 100.0 & 0.0 & 3, 5 & \checkmark \\
        7  & 0 & 1180 & 1648 & 100.0 & 0.0 & 3, 5 & \checkmark \\
        7  & 1 & 1042 & 1261 & 100.0 & 0.0 & 3, 5 & \checkmark \\
        7  & 2 & 874  & 1075 & 100.0 & 0.0 & 3, 5 & \checkmark \\
        7  & 3 & 869  & 1071 & 100.0 & 0.0 & 3, 5 & \checkmark \\
        11 & 0 & 569  & 922  & 100.0 & 0.0 & 3, 5 & \checkmark \\
        11 & 1 & 525  & 778  & 100.0 & 0.0 & 3, 5 & \checkmark \\
        11 & 2 & 418  & 569  & 100.0 & 0.0 & 3, 5 & \checkmark \\
        11 & 3 & 461  & 607  & 100.0 & 0.0 & 3, 5 & \checkmark \\
        14 & 0 & 570  & 923  & 100.0 & 0.0 & Triviales & $\Delta$ \\
        14 & 1 & 526  & 779  & 100.0 & 0.0 & Triviales & $\Delta$ \\
        14 & 2 & 427  & 588  & 100.0 & 0.0 & Triviales & $\Delta$ \\
        14 & 3 & 427  & 588  & 100.0 & 0.0 & Triviales & $\Delta$ \\
        \bottomrule
    \end{tabular}
\end{table}

\newpage

\begin{table}[htpb]
    \centering
    \caption{Comparativa del \textit{overhead} de ruteo Heavy-Hex: profundidad y conteo de compuertas de dos qubits entre la simulación pura (\textit{all-to-all}) y la transpilación a \texttt{FakeTorino} bajo los cuatro niveles de optimización del transpilador de Qiskit. Los ratios $D_{2Q}^{\text{Torino}}/D_{2Q}^{\text{Puro}}$ y $G_{2Q}^{\text{Torino}}/G_{2Q}^{\text{Puro}}$ cuantifican el factor de amplificación introducido por la topología.}
    \label{tab:transpilacion_faketorino}
    \begin{tabular}{cccccccc}
        \toprule
        $a$ & \textbf{Opt} & \textbf{D.Puro$_{2Q}$} & \textbf{D.Torino$_{2Q}$} & \textbf{Ratio $D$} & \textbf{G.Puro$_{2Q}$} & \textbf{G.Torino$_{2Q}$} & \textbf{Ratio $G$} \\
        \midrule
        4  & 0 & 233 & 570  & 2.45$\times$ & 236 & 923  & 3.91$\times$ \\
        4  & 1 & 233 & 526  & 2.26$\times$ & 236 & 779  & 3.30$\times$ \\
        4  & 2 & 233 & 433  & 1.86$\times$ & 236 & 584  & 2.47$\times$ \\
        4  & 3 & 233 & 439  & 1.88$\times$ & 236 & 600  & 2.54$\times$ \\
        7  & 0 & 466 & 1180 & 2.53$\times$ & 472 & 1648 & 3.49$\times$ \\
        7  & 1 & 466 & 1042 & 2.24$\times$ & 472 & 1261 & 2.67$\times$ \\
        7  & 2 & 461 & 874  & 1.90$\times$ & 466 & 1075 & 2.31$\times$ \\
        7  & 3 & 461 & 869  & 1.89$\times$ & 466 & 1071 & 2.30$\times$ \\
        11 & 0 & 232 & 569  & 2.45$\times$ & 235 & 922  & 3.92$\times$ \\
        11 & 1 & 232 & 525  & 2.26$\times$ & 235 & 778  & 3.31$\times$ \\
        11 & 2 & 230 & 418  & 1.82$\times$ & 233 & 569  & 2.44$\times$ \\
        11 & 3 & 230 & 461  & 2.00$\times$ & 233 & 607  & 2.61$\times$ \\
        14 & 0 & 233 & 570  & 2.45$\times$ & 236 & 923  & 3.91$\times$ \\
        14 & 1 & 233 & 526  & 2.26$\times$ & 236 & 779  & 3.30$\times$ \\
        14 & 2 & 231 & 427  & 1.85$\times$ & 234 & 588  & 2.51$\times$ \\
        14 & 3 & 231 & 427  & 1.85$\times$ & 234 & 588  & 2.51$\times$ \\
        \bottomrule
    \end{tabular}
\end{table}

\begin{table}[H]
    \centering
    \caption{Métricas de transpilación para el circuito \texttt{RegisterQC} ($N = 15$) bajo los
    mapeos trivial y SABRE, transpilado sobre la topología Heavy-Hex de \texttt{FakeTorino} con
    4096 \textit{shots} y semilla 457. El conteo de SWAPs explícitas reportado por Qiskit es nulo en
    todas las configuraciones debido a que el pase de traducción al conjunto de compuertas
    nativas (\texttt{BasisTranslator}) descompone cada SWAP en tres compuertas de dos qubits,
    que quedan contabilizadas en Gates$_{2Q}$.}
    \label{tab:transpilacion_faketorino_sabre}
    \small
    \begin{tabular}{cclcccccc}
        \hline
        $a$ & Layout  & Opt & D.\ total & Depth$_{2Q}$ & Gates$_{2Q}$ & $PST$ (\%) & $\mathcal{F}_H$ \\
        \hline
        4  & trivial & 0 & 2615 & 570  & 923  & 100.0 & 1.0000 \\
        4  & trivial & 1 & 1841 & 533  & 833  & 100.0 & 0.9998 \\
        4  & trivial & 2 & 1839 & 521  & 793  & 100.0 & 1.0000 \\
        4  & trivial & 3 & 1843 & 521  & 793  & 100.0 & 0.9999 \\
        4  & SABRE   & 0 & 2509 & 504  & 749  & 100.0 & 0.9999 \\
        4  & SABRE   & 1 & 1869 & 526  & 779  & 100.0 & 1.0000 \\
        4  & SABRE   & 2 & 1584 & 433  & 584  & 100.0 & 1.0000 \\
        4  & SABRE   & 3 & 1598 & 439  & 600  & 100.0 & 1.0000 \\
        \hline
        7  & trivial & 0 & 5521 & 1180 & 1648 & 100.0 & 0.9999 \\
        7  & trivial & 1 & 3758 & 1094 & 1432 & 100.0 & 0.9999 \\
        7  & trivial & 2 & 3599 & 1019 & 1349 & 100.0 & 0.9999 \\
        7  & trivial & 3 & 3608 & 1019 & 1349 & 100.0 & 0.9999 \\
        7  & SABRE   & 0 & 4952 & 1015 & 1255 & 100.0 & 0.9999 \\
        7  & SABRE   & 1 & 3576 & 1042 & 1261 & 100.0 & 1.0000 \\
        7  & SABRE   & 2 & 3186 & 874  & 1075 & 100.0 & 1.0000 \\
        7  & SABRE   & 3 & 3172 & 869  & 1071 & 100.0 & 0.9996 \\
        \hline
        11 & trivial & 0 & 2607 & 569  & 922  & 100.0 & 1.0000 \\
        11 & trivial & 1 & 1845 & 532  & 832  & 100.0 & 0.9999 \\
        11 & trivial & 2 & 1785 & 508  & 778  & 100.0 & 0.9999 \\
        11 & trivial & 3 & 1787 & 508  & 778  & 100.0 & 1.0000 \\
        11 & SABRE   & 0 & 2499 & 503  & 748  & 100.0 & 1.0000 \\
        11 & SABRE   & 1 & 1864 & 525  & 778  & 100.0 & 1.0000 \\
        11 & SABRE   & 2 & 1543 & 418  & 569  & 100.0 & 0.9999 \\
        11 & SABRE   & 3 & 1636 & 461  & 607  & 100.0 & 1.0000 \\
        \hline
        14 & trivial & 0 & 2618 & 570  & 923  & 100.0 & 0.9999 \\
        14 & trivial & 1 & 1861 & 533  & 833  & 100.0 & 0.9999 \\
        14 & trivial & 2 & 1802 & 512  & 785  & 100.0 & 1.0000 \\
        14 & trivial & 3 & 1803 & 512  & 785  & 100.0 & 0.9998 \\
        14 & SABRE   & 0 & 2509 & 504  & 749  & 100.0 & 1.0000 \\
        14 & SABRE   & 1 & 1869 & 526  & 779  & 100.0 & 1.0000 \\
        14 & SABRE   & 2 & 1567 & 427  & 588  & 100.0 & 1.0000 \\
        14 & SABRE   & 3 & 1567 & 427  & 588  & 100.0 & 1.0000 \\
        \hline
    \end{tabular}
\end{table}

\begin{table}[H]
    \centering
    \caption{Reducción porcentual lograda por SABRE respecto al mapeo trivial en
    profundidad de dos qubits ($D_{2Q}$) y conteo de compuertas de dos qubits ($G_{2Q}$),
    para cada base $a$ y nivel de optimización. Se observa una dependencia no monótona
    con el nivel de optimización: un mínimo local en Opt=1 y máximos en Opt=2.}
    \label{tab:reduccion_sabre}
    \begin{tabular}{cccccccc}
        \hline
        $a$ & Opt & $D_{2Q}$ Triv. & $D_{2Q}$ SABRE & Red.\ $D$ (\%) & $G_{2Q}$ Triv. & $G_{2Q}$ SABRE & Red.\ $G$ (\%) \\
        \hline
        4  & 0 & 570  & 504  & 11.6 & 923  & 749  & 18.9 \\
        4  & 1 & 533  & 526  &  1.3 & 833  & 779  &  6.5 \\
        4  & 2 & 521  & 433  & 16.9 & 793  & 584  & 26.4 \\
        4  & 3 & 521  & 439  & 15.7 & 793  & 600  & 24.3 \\
        \hline
        7  & 0 & 1180 & 1015 & 14.0 & 1648 & 1255 & 23.8 \\
        7  & 1 & 1094 & 1042 &  4.8 & 1432 & 1261 & 11.9 \\
        7  & 2 & 1019 &  874 & 14.2 & 1349 & 1075 & 20.3 \\
        7  & 3 & 1019 &  869 & 14.7 & 1349 & 1071 & 20.6 \\
        \hline
        11 & 0 &  569 &  503 & 11.6 &  922 &  748 & 18.9 \\
        11 & 1 &  532 &  525 &  1.3 &  832 &  778 &  6.5 \\
        11 & 2 &  508 &  418 & 17.7 &  778 &  569 & 26.9 \\
        11 & 3 &  508 &  461 &  9.3 &  778 &  607 & 22.0 \\
        \hline
        14 & 0 &  570 &  504 & 11.6 &  923 &  749 & 18.9 \\
        14 & 1 &  533 &  526 &  1.3 &  833 &  779 &  6.5 \\
        14 & 2 &  512 &  427 & 16.6 &  785 &  588 & 25.1 \\
        14 & 3 &  512 &  427 & 16.6 &  785 &  588 & 25.1 \\
        \hline
    \end{tabular}
\end{table}

\section{Tablas de resultados en simulación con Ruido}

\begin{table}[H]
    \centering
    \caption{Métricas de la simulación con ruido (FakeTorino, sin mitigación) para $N = 15$.
    Se reportan la profundidad de compuertas de dos qubits (D2Q), el conteo de compuertas de
    dos qubits (G2Q), la probabilidad de señal teórica ($PST$), la fidelidad de Hellinger
    ($\mathcal{F}_H$) y los factores extraídos.}
    \label{tab:faketorino_sin_mitigacion}
    \begin{tabular}{cccccccc}
        \hline
        Base $a$ & opt & D2Q & G2Q & $PST$ (\%) & $\mathcal{F}_H$ & Factores & Estado \\
        \hline
        4 & 0 & 504  & 749  & 25.00 & 0.1482 & $\{3,5\}$ & \checkmark \\
        4 & 3 & 439  & 600  &  4.30 & 0.0010 & $\{3,5\}$ & \checkmark \\
        7 & 0 & 1015 & 1255 & 43.16 & 0.2285 & $\{3,5\}$ & \checkmark \\
        7 & 3 & 869  & 1071 & 59.18 & 0.4133 & $\{3,5\}$ & \checkmark \\
        \hline
    \end{tabular}
\end{table}

\begin{table}[H]
    \centering
    \caption{Análisis detallado de fidelidad con ruido ($\texttt{opt} = 3$, FakeTorino).
    Se reportan la fidelidad de Hellinger ruidosa e ideal, la degradación
    $\Delta\mathcal{F}_H = \mathcal{F}_H^{\text{ideal}} - \mathcal{F}_H^{\text{ruido}}$.}
    \label{tab:faketorino_fidelidad_detalle}
    \begin{tabular}{ccccccc}
        \hline
        $a$ & $r$ & $\mathcal{F}_H^{\text{ruido}}$ & $\mathcal{F}_H^{\text{ideal}}$
            & $\Delta\mathcal{F}_H$ & Prob.\ picos \\
        \hline
        4 & 2 & 0.0010 & 1.0000 & 0.9990 & $0.20\%$ \\
        7 & 4 & 0.4559 & 0.9997 & 0.5438 & $46.29\%$ \\
        \hline
    \end{tabular}
\end{table}

\begin{table}[H]
    \centering
    \caption{Comparación de técnicas de supresión ($\texttt{opt} = 3$, FakeTorino,
    512~\textit{shots}). Se reportan $PST$, $\mathcal{F}_H$ y los factores extraídos para cada variante.}
    \label{tab:faketorino_mitigacion}
    \begin{tabular}{clcccc}
        \hline
        $a$ & Variante & $PST$ (\%) & $\mathcal{F}_H$ & Factores & Estado \\
        \hline
        4 & Sin supresión      &  3.52 & 0.0039 & $\{3,5\}$ & \checkmark \\
        4 & DD ($XY4$)      &  3.91 & 0.0073 & $\{3,5\}$ & \checkmark \\
        4 & PT (activo)         &  2.34 & 0.0057 & $\{3,5\}$ & \checkmark \\
        4 & DD + PT             &  3.32 & 0.0029 & $\{3,5\}$ & \checkmark \\
        \hline
        7 & Sin supresión      & 56.64 & 0.4029 & $\{3,5\}$ & \checkmark \\
        7 & DD ($XY4$)      & 49.61 & 0.3300 & $\{3,5\}$ & \checkmark \\
        7 & PT (activo)         & 54.69 & 0.3998 & $\{3,5\}$ & \checkmark \\
        7 & DD + PT             & 46.48 & 0.2906 & $\{3,5\}$ & \checkmark \\
        \hline
    \end{tabular}
\end{table}

\begin{table}[H]
    \centering
    \caption{Mediciones más frecuentes del QPE con ruido para $a = 7$, $r = 4$
    (DD + PT, $\texttt{opt} = 3$). Las fases esperadas son
    $\varphi \in \{0, 0.25, 0.50, 0.75\}$. Se reporta la clasificación de cada
    medición como señal (fase cercana a un valor teórico) o ruido (fase espúrea).}
    \label{tab:qpe_ruido_a7_ddpt}
    \begin{tabular}{cccccccc}
        \hline
        \# & $|y\rangle$ & $\varphi = y/2^{9}$ & Aprox.\ $s/r$ & $r_{\text{cand}}$ & Tipo & Prob. \\
        \hline
        1 & $|000000000\rangle$ & 0.0000 & $0/1$ & 1 & señal & $8.98\%$ \\
        2 & $|100000000\rangle$ & 0.5000 & $1/2$ & 2 & señal & $7.23\%$ \\
        3 & $|010000000\rangle$ & 0.2500 & $1/4$ & 4 & señal & $7.03\%$ \\
        4 & $|110000000\rangle$ & 0.7500 & $3/4$ & 4 & señal & $6.84\%$ \\
        5 & $|100000001\rangle$ & 0.5020 & $1/2$ & 2 & señal & $2.15\%$ \\
        6 & $|110000001\rangle$ & 0.7520 & $3/4$ & 4 & señal & $1.76\%$ \\
        7 & $|100100000\rangle$ & 0.5625 & $5/9$ & 9 & ruido & $1.76\%$ \\
        8 & $|011000000\rangle$ & 0.3750 & $3/8$ & 8 & ruido & $1.76\%$ \\
        \hline
    \end{tabular}
\end{table}

\section{Tablas de resultados de ejecución en hardware IBM}

\begin{table}[H]
    \centering
    \caption{Impacto del nivel de optimización del transpilador en hardware real
    (\texttt{ibm\_torino}, $a = 4$, $\texttt{approx} = 0.7$, 4096~\textit{shots},
    Pauli Twirling activado). Se reportan la profundidad 2Q, el conteo de compuertas
    CZ, el número de SWAPs estimados, la señal $PST$ y la fidelidad $\mathcal{F}_H$.}
    \label{tab:opt_level_hardware}
    \begin{tabular}{ccccccc}
        \hline
        Nivel opt. & Depth 2Q & CZ   & SWAPs est. & $PST_{\text{HW}}$ (\%) & $\mathcal{F}_H^{\text{HW}}$ & Outcomes únicos \\
        \hline
        0 & 770 & 1220 & 113 &  1.1 & 0.011 & 512 \\
        1 & 526 &  779 &  65 &  3.1 & 0.031 & 505 \\
        2 & 116 &  116 &   0 & 70.8 & 0.529 &  57 \\
        \hline
    \end{tabular}
\end{table}

\bibliography{bibliografia}
\end{document}



















